% ATTENZIONE!!! 
% Per far funzionare i collegamenti ipertestuali si raccomanda di usare
%	\classname{nomedellaclasse}
% per le classi dello stesso package mentre invece 
%	\hyperref[nomedellaclasse]{\ttfamily{}nomequalificatodellaclasse}
% per le classi che non sono dello stesso package e che hanno il nome completo

% **************************************************
% Macro specifiche per il documento corrente
% **************************************************
% Nome
\newcommand{\docName}{Definizione di prodotto}
% Nome file
\newcommand{\docFileName}{definizione\_di\_prodotto.2.0.pdf}
% Versione
\newcommand{\docVers}{2.0}
% Data creazione
\newcommand{\creationDate}{2013-03-04}
% Data ultima modifica
\newcommand{\modificationDate}{2013-03-19}
% Stato in {Approvato, Non approvato}
\newcommand{\docState}{Approvato}
% Uso in {Interno, Esterno}
\newcommand{\docUsage}{Esterno}
% Destinatari da specificare come nome1\\ &nome2\\ ecc.
\newcommand{\docDistributionList}{Prof. Tullio Vardanega\\&Prof. Riccardo Cardin\\&Dott. Gregorio Piccoli\\&Team SoftwareSynthesis}
% Redattori da specificare come nome1\\ &nome2\\ ecc.
\newcommand{\docAuthors}{Andrea Rizzi\\&Elena Zecchinato\\&Marco Schivo \\&Stefano Farronato}
% Approvato da
\newcommand{\approvedBy}{Marco Schivo}
% Verificatori
\newcommand{\verifiedBy}{Andrea Meneghinello}
% Perscorso (relativo o assoluto) che punta alla directory contenente shared/
% come sua sottodirectory (per comodità chiamiamola 'doc root').
\newcommand{\docRoot}{..}
% definire se si vuole l'indice delle tabelle
\def\INDICETABELLE{false}
% definire se si vuole l'indice delle figure
\def\INDICEFIGURE{false}

% importa il preambolo condiviso da tutti i documenti
% shared/preamble.tex
%
% Questo documento contiene la parte del preambolo condivisa e viene pertanto
% richiamato nel 'master' di tutti i documenti di progetto.  Al suo interno
% contiene le inclusioni (e le configurazioni) di tutti i package richiesti per
% la compilazione dei documenti, le macro di carattere generale e la definizione
% degli stili di pagina.

\documentclass[a4paper,10pt,openright]{article}

% **************************************************
% Macro generiche
% **************************************************
\newcommand{\team}{Software Synthesis}                    % chi siamo
\newcommand{\email}{software.synthesis@gmail.com}         % e-mail
\newcommand{\caName}{}                                    % titolo capitolato
\newcommand{\caDescr}{}                                   % descrizione
\newcommand{\inglese}[1]{\textit{#1}}

% **************************************************
% Codifica e lingua dei documenti
% **************************************************
\usepackage[utf8x]{inputenc}                              % codifica caratteri dei documenti sorgenti
\usepackage[english,italian]{babel}                       % localizzazione ai fini di sillabazione e cross-references
\usepackage[T1]{fontenc}                                  % codifica font di output

% **************************************************
% Definizione geometria della pagina
% **************************************************
\usepackage[a4paper,head=4cm,top=4.5cm,bottom=3cm,left=3cm,right=3cm,bindingoffset=5mm]{geometry}

% *************************************************
% Intestazioni e piè di pagina personalizzati
% *************************************************
\usepackage{fancyhdr}
% stile normale
\fancypagestyle{normal}{
\fancyhead{}                                              % intestazione
\fancyhead[RE,RO]{
\begin{picture}(0,0)
  \put(-410,0){\includegraphics[width=1.02\textwidth]{header_logo}}
  \put(-410,10){\sffamily\large\leftmark}
\end{picture}
\vspace{-4pt}
}
\renewcommand{\headrulewidth}{.4pt}                       % riga sotto l'intestazione
\cfoot{}                                                  % piè di pagina
\fancyfoot[RO,LE]{\sffamily
  pag.~\thepage{} di \pageref{LastPage}}                  % a dx nelle pag. dispari e a sx in quelle pari
\fancyfoot[RE,LO]{\sffamily\docFileName{} -- v.\docVers}
\renewcommand{\footrulewidth}{.4pt}                       % riga sopra il piè di pagina
}
% stile per gli indici
\fancypagestyle{toc}{
\fancyhead{}                                              % intestazione
\fancyhead[RE,RO]{
\begin{picture}(0,0)
  \put(-410,0){\includegraphics[width=1.02\textwidth]{header_logo}}
\end{picture}
}
\renewcommand{\headrule}{}                                % nessuna riga sotto l'intestazione
\cfoot{}                                                  % piè di pagina
\fancyfoot[RO,LE]{\sffamily\thepage{}}                    % a dx nelle pag. dispari e a sx in quelle pari
\fancyfoot[RE,LO]{\sffamily\docFileName{} -- v.\docVers}
\renewcommand{\footrulewidth}{.4pt}                       % riga sopra il piè di pagina
}

\pagestyle{fancy}                                         % premetto: non so usare bene le marche:
\renewcommand{\sectionmark}[1]{\markboth{#1}{#1}}         % se qualcuno ha idee migliori si faccia avanti!

% **************************************************
% Tabelle
% **************************************************
\usepackage{tabularx}                                     % tabelle di larghezza fissa con una o più colonne variabili
\usepackage{multirow}                                     % colonne con colonne che si estendono per più righe
\usepackage{booktabs}                                     % per inserire l'ambiente table e le righe orizz. nelle tabelle
\usepackage{longtable}			                          % tabelle oltre i limiti di pagina

% **************************************************
% Cross-references e collegamenti ipertestuali
% **************************************************
\usepackage[hidelinks]{hyperref}
\hypersetup{%
  colorlinks=false, linktocpage=false, pdfborder={0,0,0}, pdfstartpage=3, pdfstartview=FitV,%
  urlcolor=Cyan, linkcolor=Cyan, citecolor=Black, %pagecolor=Black,%
  pdftitle={\docName}, pdfauthor={\team}, pdfsubject={}, pdfkeywords={},%
  pdfcreator={pdflatex}, pdfproducer={pdflatex with hyperref package}%
}

% **************************************************
% Immagini e grafica
% **************************************************
\usepackage{graphicx}                                     % supporto ad aspetti avanzati delle immagini
\graphicspath{{\docRoot/pics/}}                           % percorso contenente tutti i file immagini
\usepackage{color}                                        % permette di colorare facilmente il testo

% **************************************************
% Altri pacchetti opzionali
% **************************************************     
\usepackage{lastpage}                                     % per sapere il numero totale di pagine
\usepackage{lipsum}                                       % genera "dummy text" per prove di impaginazione
\usepackage{eurosym}                                      % per il simbolo dell'euro usare \EUR{x} dove x è l'importo


% macro specifiche per il documento corrente
\newcommand{\classsection}[1]{\subsubsection{#1}\label{#1}}
\newcommand{\classname}[1]{\hyperref[#1]{\ttfamily#1}}

% Fine del preambolo e inizio del documento
\begin{document}

% Inclusione della prima pagina
% shared/firstpage.tex
%
% Questo documento definisce il contenuto della prima pagina, che si suppone
% essere uguale in tutti i documenti.  Oltre al logo e al titolo, la prima
% pagina contiene i metadati relativi al documento in cui viene inclusa.


% rimuove intestazioni e piè di pagina
\pagestyle{empty}

\begin{center}

% logo del gruppo
\includegraphics[width=1.5\textwidth]{logo}

\vspace{1in}

% titolo del documento
{\Huge\bfseries \docName}

\vspace{1in}

% tabella riepilogativa
\begin{tabularx}{.7\textwidth}{>{\bfseries\sffamily}l>{\sffamily}l}
\toprule
\multicolumn{2}{>{\sffamily}c}{Informazioni sul documento}\\
\midrule
Nome file:            & \docFileName\\
Versione:             & \docVers\\
Data creazione:       & \creationDate\\
Data ultima modifica: & \modificationDate\\
Stato:                & \docState\\
Uso:                  & \docUsage\\
Redattori:            & \docAuthors\\
Approvato da:         & \approvedBy\\
Verificatori:         & \verifiedBy\\
\bottomrule
\end{tabularx}

\end{center}

\newpage


%---------------------------RUOLI----------------------------
%FASE 1:
%Progettisti: STEFANO,RIZZI,SCHIVO,ELENA

%VerificatorI: MENE,TRES
%Responsabile finale supremo: DIEGO
%------------------------------------------------------------

% Storico delle modifiche
\section*{Storia delle modifiche}
\begin{center}
\begin{longtable}{lp{.32\textwidth}lll}
\toprule
Versione & Descrizione intervento & Membro & Ruolo & Data\\
\midrule % inserire qui il contenuto della tabella
1.0 & Approvazione documento &Diego Beraldin &Responsabile  & 2013-03-09\\
0.10 & Correzioni errori rilevati da verificatori &Schivo Marco  &  Progettista& 2013-03-08\\
0.9 & Verifica totale documento &Andrea Meneghinello & Verificatore & 2013-03-08\\
0.8 & Descritto classi di abook.authentication, server.authentication.servlet, server.connection & Elena Zecchinato &  Progettista& 2013-03-07\\
0.7 & Descritto PresenterMediator, CallHistoryPanelPresenter, SearchResultPanelPresenter, ToolsPanelPresenter, AccountSettingsPanelPresenter &Schivo Marco &  Progettista& 2013-03-07\\
0.7 & Descritto AddressBookPanelPresenter, GroupPanelPresenter, LoginPanelPresenter &Schivo Marco &  Progettista& 2013-03-06\\
0.6 & Descritto classi di server.abook.servlet, server.call, server.message &Elena Zecchinato &  Progettista& 2013-03-06\\
0.5 & Descritto RegisterPanelPresenter, CommunicationPanelPresenter, ContactPanelPresenter, MainPanelPresenter &Stefano Farronato &  Progettista& 2013-03-05\\
0.4 & Descritte classi package server.dao e server.abook &Andrea Rizzi &  Progettista& 2013-03-05\\
0.3 & Inizio descrizione parte clientpresenter &Stefano Farronato &  Progettista& 2013-03-04\\
0.2 & Stesura sezione ``Standard di progetto'' ed inizio descrizione server &Andrea Rizzi &  Progettista& 2013-03-04\\
0.1 & Creazione del documento e stesura delle sezioni ``Introduzione'' e ``Riferimenti'' &Elena Zecchinato &Progettista  & 2013-03-04\\
\bottomrule
\end{longtable}
\end{center}
\newpage

% inclusione dell'indice
% shared/toc.tex
%
% Questo file contiene le istruzioni che generano l'indice o gli indici del
% documento (utile nel caso in cui decidessimo di avere anche un indice delle
% tabelle e/o un indice delle figure).

\pagestyle{toc}
\pagenumbering{roman}

\tableofcontents

\newpage


% Alcuni aggiustamenti per le pagine
\pagenumbering{arabic}
\setcounter{page}{1}
\pagestyle{normal}

% Qui ha inizio il documento vero e proprio...
\newpage

\section{Introduzione}
\subsection{Scopo del prodotto}
\purpose

\subsection{Scopo del documento}
Il presente documento presenta una descrizione dettagliata dell'architettura del sistema software destinata alla realizzazione del prodotto \caName{} coerentemente con la progettazione ad alto livello descritta nell'allegato \textit{specifica\_tecnica.2.0.pdf}.

A tal fine si riporta per ognuno dei componenti definiti nel documento di specifica tecnica una descrizione delle classi in termini di operazioni disponibili, proprietà, responsabilità e collaborazioni. Il contenuto del presente documento ha inoltre valore vincolante per i programmatori, pertanto avranno l'obbligo di attenersi alle disposizioni in esso contenute senza alcuna possibilità di deroga.

\subsection{Glossario}
\glossaryIntro

\subsection{Convenzioni di scrittura}
% vedere issue #56 al riguardo
Al fine di rendere quanto più agevole possibile la consultazione del documento da parte dei programmatori e del committente, è stata adottata una serie di accorgimenti sia a livello di riferimenti sulla nomenclatura delle classi sia a livello cromatico per campi dati e metodi. 
Tali norme possono essere consultate in dettaglio nel documento \textit{norme\_di\_progetto.3.0.pdf} allegato.
\clearpage

\section{Riferimenti}
\subsection{Normativi}
\begin{itemize}
\item[] \textit{piano\_di\_qualifica.3.0.pdf} allegato.
\item[] \textit{norme\_di\_progetto.3.0.pdf} allegato.
\item[] \textit{specifica\_tecnica.2.0.pdf} allegato
\end{itemize}

\subsection{Informativi}
\begin{itemize}
\item[] Capitolato d'appalto: \caName{}, v1.0, redatto e rilasciato dal proponente Zucchetti s.r.l. reperibile all'indirizzo \url{http://www.math.unipd.it/~tullio/IS-1/2012/Progetto/C1.pdf};
\item[] testo di consultazione: \textit{Software Engineering (8th edition) Ian Sommerville, Pearson Education | Addison Wesley};
\item[] manuale all'utilizzo dei design pattens: \textit{Design Patterns, Elementi per il riuso di software a oggetti -- (1/Ed. italiana) Eric Gamma, Richard Helm, Ralph Johnson, John Vlissides, Pearson Education};
\item[] \textit{glossario.3.0.pdf} allegato.
\end{itemize}
\clearpage

\section{Standard di progetto}

\subsection{Standard di progettazione architetturale}
Lo sviluppo del progetto ha seguito le regole architetturali specificate nel documento  \textit{norme\_di\_progetto.3.0.pdf} allegato.

\subsection{Standard di documentazione del codice}
Le regole che definiscono la documentazione del codice relativo al funzionamento del prodotto sono specificate nel documento \textit{norme\_di\_progetto.3.0.pdf} allegato.

\subsection{Standard di denominazione di entità e relazioni}
Le convenzioni relative alla denominazione delle entità e le relative relazioni sono specificate nel documento \textit{norme\_di\_progetto.3.0.pdf} allegato.

\subsection{Standard di programmazione}
Le regole relative agli standard di programmazione sono enunciate nel documento \textit{norme\_di\_progetto.3.0.pdf} allegato.

\subsection{Strumenti di lavoro}
Gli strumenti utilizzati per la stesura e lo sviluppo sono specificati nei documenti \textit{norme\_di\_progetto.3.0.pdf} e \textit{piano\_di\_qualifica.3.0.pdf} allegati.

\clearpage

\section{Specifica sotto-architettura sever}\label{sec:serverarchitecture}

\subsection{Package org.softwaresynthesis.mytalk.server}\label{sec:server}

\classsection{IMyTalkObject}

\subsubsection*{Funzione}
Interfaccia usata per rappresentare un \texttt{transfer object} mappato mediante \texttt{Hibernate}.

\subsubsection*{Relazioni d'uso}

Nessuna relazione evidenziata

\subsubsection*{Classi estese ed interfacce implementate}

Nessuna relazione evidenziata.

\subsubsection*{Metodi}
\begin{description}
	\item{\method{\underline{+ getId(): Long}}}\\
	Restituisce l'identificativo dell'oggetto.
	\item{\method{\underline{+ setId(id: Long): void}}}\\
	Imposta l'identificativo dell'oggetto.
\end{description}

\classsection{IController}

\subsubsection*{Funzione}
Interfaccia rappresentante il comportamento di un generico \inglese{controller}.

\subsubsection*{Relazioni d'uso}

Nessuna relazione evidenziata

\subsubsection*{Classi estese ed interfacce implementate}

Nessuna relazione evidenziata.

\subsubsection*{Metodi}
\begin{description}
	\item{\method{-- execute(request: HttpServletRequest, response: HttpServletResponse): void}}\\
	Metodo usato per l'esecuzione del comando associato al \inglese{controller}.
\end{description}

\classsection{AbstractController}

\subsubsection*{Funzione}
Implementazione di \classname{IController}. Classe astratta che rappresenta un generico controller. Richiede la ridefinizione del metodo \method{doAction()} da parte delle classi figlie.

\subsubsection*{Relazioni d'uso}
\begin{itemize}
	\item \texttt{java.io.IOException}: eccezione richiamabile dal metodo\method{doAction()}.
	\item \texttt{java.io.PrintWriter}: classe istanziata all'interno del metodo \method{doAction()}. Usata per scrivere l'output della \inglese{servlet}.
	\item \texttt{javax.servlet.ServletException}: eccezione sollevabile dai metodi \method{doAction()}.
	\item \texttt{javax.servlet.http.HttpServletRequest}: classe usata per interagire con le richieste \texttt{ajax} inoltrate dal client.
	\item \texttt{javax.servlet.http.HttpServletResponse}: classe usata per interagire con le richieste \texttt{ajax} inoltrate dal client.
	\item \classname{authentication.security.ISecurityStrategy}: usata per la criptazione dei dati.
\end{itemize}

\subsubsection*{Classi estese ed interfacce implementate}

\begin{itemize}
	\item \classname{IController}: interfaccia d'implementazione.
\end{itemize}

\subsubsection*{Attributi}

Nessun attributo evidenziato

\subsubsection*{Metodi}

\begin{description}
	\item{\method{-- check(request: HttpServletRequest): boolean}}\\
	Metodo usato per verificare se l'utente che richiede l'esecuzione di un operazione, è autenticato nel sistema. Il metodo definisce inizialmente una variabile booleana result, usata per ritornare una segnalazione della buona riuscita del metodo. result è impostata di \inglese{default} a \inglese{true}. Quindi si procede controllando se la mail dell'utente (\memberdata{userMail}) corrisponde a quella presente nel campo username della sessione corrente (session.getAttribute(''username'')). se il \inglese{match} è negativo, result viene impostata a \inglese{false}. Il metodo termina ritornando il contenuto di result.
	\item{\method{-- execute(request: HttpServletRequest, response: HttpServletResponse): void}}\\
	Il metodo inizia definendo un oggetto \texttt{PrintWriter} writer. Quindi viene effettuato un controllo sul richiedente mediante il richiamo del metodo \method{check()}. Se tale metodo ritorna \inglese{true} allora il flusso principale prosegue richiamando il metodo \method{doAction}. In caso contrario il metodo termina salvando nel \texttt{writer} il valore ''null''.
\end{description}

\classsection{ControllerManager}

\subsubsection*{Funzione}
Estensione di \texttt{org.apache.catalina.websocket.WebSocketServlet} ed implementazione di \inglese{servlet}. ControllerManager rappresenta l'implementazione del pattern \texttt{front controller}. Esso fornisce un unico punto d'accesso al sistema. Tale implementazione avviene definendo la classe come una \inglese{servlet} il cui compito sarà quello di delegare al controller corretto la richiesta inoltrata dal client.

\subsubsection*{Relazioni d'uso}
\begin{itemize}
	\item \texttt{java.io.IOException}: eccezione sollevabile in seguito ad un errato input utente;
	\item \texttt{java.util.Enumeration}: usata per definire gli attributi usati dalla classe;
	\item \texttt{java.util.HashMap}: usata per definire gli attributi usati dalla classe;
	\item \texttt{java.util.Hashtable}: usata per definire gli attributi usati dalla classe;
	\item \texttt{java.util.Map}: usata per definire gli attributi usati dalla classe;
	\item \texttt{java.util.ResourceBundle}: usata per definire gli attributi usati dalla classe;
	\item \texttt{javax.servlet.ServletConfig}: usata per eseguire delle operazioni di configurazione della \inglese{servlet};
	\item \texttt{javax.servlet.ServletException}: eccezione che può essere sollevata in seguito all'esecuzione del metodo \method{doPost()} della classe;
	\item \texttt{javax.servlet.annotation.WebServlet}: classe necessaria per il corretto funzionamento della \inglese{servlet};
	\item \texttt{javax.servlet.http.HttpServletRequest}: usata per le richieste ajax inoltrate dal client alla \inglese{servlet};
	\item \texttt{javax.servlet.http.HttpServletResponse}: usata per le richieste ajax inoltrate dal client alla \inglese{servlet};
	\item \texttt{org.apache.catalina.websocket.StreamInbound}: classe generica che definisce un canale di comunicazione client-client;
	\item \texttt{org.softwaresynthesis.mytalk.server.connection.PushInbound}: specializzazione della classe \texttt{org.apache.catalina.websocket.StreamInbound};
	\item \texttt{org.softwaresynthesis.mytalk.server.connection.PushInbound.State}: classe atta a definire lo stato di un canale di comunicazione, usata per gestire lo stato dell'utente connesso.
\end{itemize}

\subsubsection*{Classi estese ed interfacce implementate}

\begin{itemize}
	\item \classname{javax.servlet.Servlet}: interfaccia d'implementazione.
	\item \classname{org.apache.catalina.websocket.WebSocketServlet}: classe estesa.
\end{itemize}

\subsubsection*{Attributi}

\begin{itemize}
	\item{\memberdata{\underline{-- serialVersionUID: long}}}\\
	Attributo usato per definire un id univoco alla \inglese{servlet}.
	\item{\memberdata{-- clients: Map<Long, PushInbound>}}\\
	Map usata per definire un associazione tra gli utenti connessi (riconosciuti mediante ID) e i \classname{PushInbound} a loro associati.
	\item{\memberdata{-- controllers: Hashtable<String, String>}}\\
	Hashtablet usata per mappare le associazioni funzionalità-controller associato.
\end{itemize}

\subsubsection*{Metodi}

\begin{description}
	\item{\method{+ ControllerManager(): void}}\\
	Costruttore pubblico della \inglese{servlet}. Inizializza \memberdata{controllers}.
	
	\item{\method{+ getControllers(): Hashtable<String, String>}}\\
	Metodo usato per ottenere il contenuto del campo \memberdata{controllers}. Il metodo restituisce \memberdata{controllers}.
	
	\item{\method{+ createController(classPath: string): IController}}\\
	Metodo usato per creare una nuova istanza di uno specifico controller a partire dal nome \texttt{classPath}. A tale scopo lo sviluppatore dovrà usare il metodo forName di \texttt{Class}, passando per l'appunto il nome della classe. Successivamente a partire da tale classe, sarà creata un istanza di \classname{IController}. Il metodo termina restituendo l'istanza di \classname{IController}. Si osservi che la firma del metodo dovrà specificare la possibilità che vengano invocate le seguenti eccezioni:
	\begin{itemize}
		\item \texttt{ClassNotFoundException}: se la classe non viene trovata dal class loader;
		\item \texttt{IllegalAccessException}: se non si hanno i privilegi di accesso a tale classe;
		\item \texttt{InstantiationException}: se non si riesce ad instanziare un oggetto di tale classe;
	\end{itemize}
	
	Tali eccezioni possono essere invocate a causa di un errato utilizzo di \texttt{Class}.
	
	\item{\method{+ init(configuration: ServletConfiguration): void}}\\
	Metodo usato per inizializzare la \inglese{servlet} caricando in memoria la lista dei controller. Il metodo inizia richiamando l'init di super. Quindi crea un istanza di \texttt{ResourceBundler} a partire da una chiamata a metodo \texttt{ResorceBundler.getBundler()} (passando come parametro la stringa: ''org.softwaresynthesis.mytalk.server.controllerList''). Il passo successivo consiste nel creare un oggetto di tipo \texttt{Enumeration<String>} (nominato keys), in cui salvare ciò che viene restituito da una chiamata a getKeys() (chiamato a partire dall'istanza di \texttt{ResorceBundler}). A questo punto si attuerà, all'interno di un ciclo \verb|while(keys.hasMoreElements())|, una procedura atta a salvare in controllers le coppie chiave-valore contenute in keys.
	
	\item{\method{+ destroy(): void}}\\
	Metodo usato per eliminare il contenuto di \memberdata{controllers}. Si dovrà usare il metodo \texttt{clear()} di \memberdata{controllers}.
	
	\item{\method{\# createWebSocketInbound(): StreamInbound}}\\
	Metodo usato per creare un canale di comunicazione di tipo \classname{PushInbound}. Il metodo dovrà impostare mediante setState, lo stato del canale al valore \verb|State.AVAILABLE|. Il metodo ritorna il canale di comunicazione cosi creato.
	
	\item{\method{+ putClient(n: Long, c: PushInBound): void}}\\
	Metodo usato per aggiungere in \memberdata{clients} la coppia idUtente (n) - canale PushInBound di comunicazione (c). Il metodo è composto dall'unica istruzione \verb|clients.put(n,c)|.
	
	\item{\method{+ removeClient(c: PushInbound): void}}\\
	Metodo usato per rimuovere un client dall'attributo \memberdata{clients}. Il metodo scorre il contenuto dell'attributo verificando con un costrutto if, la presenza del canale \texttt{PushInbound} c. Se lo trova, rimuove il client associato a c da \memberdata{clients}.
	
	\item{\method{\underline{+ getState(identifier: Long): String}}}\\
	Il metodo definisce una \texttt{String} result da restituire alla fine del metodo stesso. Dopo la creazione, il metodo cerca di ottenere da \memberdata{clients}, il cliente avente id uguale ad identifier. Quindi se tale client esiste realmente, il metodo procede ottenendo lo stato in cui si trova il client (client.getState()) e richiamando il suo metodo \texttt{toString()}, salva il suo valore in result. Se l'utente con id identifier non esiste in \memberdata{clients}, result = ''offline''. Il metodo termina restituendo result;
	
	\item{\method{\# doPost(request: HttpServletRequest, response: HttpServletResponse): void}}\\
	Metodo doPost() della servlet. Tale metodo è l'unico invocato direttamente dall'esterno da parte di un client, mediante un richiasta \texttt{HttpServletRequest}. Il metodo definisce localmente un \classname{IController} in cui memorizzare il controller a cui delegare l'operazione richiesta dal client (che per l'appunto ha invocato il ControllerManager). All'interno di un costrutto try-catch (nel blocco try), si memorizza l'operazione contenuta in request mediante un getParameter(''operation''). Tale sintassi è obbligatoria. Quindi se l'operazione impostata è diversa da \inglese{null}, si procede ricercando in \memberdata{controllers} il controller avente come chiave il valore di tipo String specificato come operazione richiesta dall'utente. A questo punto si crea il controller giusto (usando il metodo \method{createController(controllerName)}) e si richiama il suo metodo \texttt{execute()} passando come parametri request e response. I catch del blocco try-catch, atti a definire del blocchi per catturare le eccezioni: 
	\begin{itemize}
		\item \texttt{ClassNotFoundException};
		\item \texttt{IllegalAccessException};
		\item \texttt{InstantiationException};
	\end{itemize}
	
	hanno al loro interno delle procedure per scrivere l'ora e il nome dell'errore (inteso come nome dell'eccezione sollevata) in un file di log. Non sono fatte ulteriori segnalazioni di errore al client, poiché al fine sarà la chiamata controller.execute() (fatta in precedenza e eventuale causa del sollevamento delle eccezioni), a restituire al client una segnalazione dell'errore riscontrato. Il metodo termina distruggendo l'oggetto \classname{IController} creato per la gestione della richiesta.
	
\end{description}

\subsection{Package org.softwaresynthesis.mytalk.server.dao}\label{sec:dao}

\classsection{ISessionManager}

\subsubsection*{Funzione}
Interfaccia rappresentante il comportamento di un generico classe adibita alla creazione e alla gestione delle sessioni \texttt{Hibernate}.

\subsubsection*{Relazioni d'uso}

Nessuna relazione evidenziata

\subsubsection*{Classi estese ed interfacce implementate}

Nessuna relazione evidenziata.


\subsubsection*{Metodi}
\begin{description}
	\item{\method{\underline{+ getInstance(): ISessionManager}}}\\
	Metodo d'implementazione del pattern \texttt{Singleton}. Usato per restituire l'unica istanza creabile.
	\item{\method{\underline{+ getSessionFactory(): SessionFactory}}}\\
	Metodo usato per ottenere una sessione verso il database gestito mediante \texttt{Hibernate}.
\end{description}

\classsection{SessionManager}

\subsubsection*{Funzione}
Implementazione di \classname{ISessionManager}. Inizializza un'unica \inglese{factory} per le sessioni, utilizzate da \textit{Hibernate}, per comunicare con il database.

\subsubsection*{Relazioni d'uso}

\begin{itemize}
	\item \classname{org.hibernate.SessionFactory}: necessaria per interrogare il database.
	\item \classname{org.hibernate.cfg.Configuration}: definisce i parametri necessari per la connessione con il database. Inoltre definisce i \textit{mapping} necessari tra le classi \inglese{transfer object} e le relative tabelle nel database.
\end{itemize}

\subsubsection*{Classi estese ed interfacce implementate}

\begin{itemize}
	\item \classname{dao.ISessionManager}: interfaccia d'implementazione del package \textit{dao} usata per definire il comportamento di un generico gestore Hibernate. Con gestore Hibernate intendiamo una classe in grado di creare delle sessioni con il database.
\end{itemize}

\subsubsection*{Attributi}

\begin{itemize}
	\item{\memberdata{\underline{-- instance: ISessionManager}}}\\
	Attributo usato per implementare il \inglese{pattern} singleton. Tale istanza verrà inizializzata tramite il metodo statico \method{getIstance()}, assicurando che l'attributo non sia già stato inizializzato in precedenza.
	\item{\memberdata{-- sessionFactory: SessionFactory}}\\
	Attributo contenente la \inglese{factory} delle sessioni verso il database.
\end{itemize}

\subsubsection*{Metodi}

\begin{description}
	\item{\method{-- SessionManager()}}\\
	Costruttore privato della classe, definito \inglese{private} in correlazione all'applicazione del \inglese{pattern} singleton.
	\item{\method{\underline{+ getInstance(): SessionManager}}}\\
	Metodo pubblico che ritorna l'istanza \memberdata{istance}. Il metodo controlla se \memberdata{istance} è già stata inizializzata, nel caso in cui non lo sia il metodo dovrà generare un istanza di SessionManager richiamando il costruttore privato \method{SessionManager()} e assegnare il valore ritornato all'attributo \memberdata{istance}. Il programma termina restituendo \memberdata{istance}.
	\item{\method{\underline{+ getSessionFactory(): SessionFactory}}}\\
	Metodo che ritorna l'attributo \memberdata{sessionFactory}.
\end{description}

\classsection{DataPersistanceManager}

\subsubsection*{Funzione}
Classe che rappresenta un punto d'accesso al database del sistema. Implementata mediante le funzionalità di \texttt{Hibernate}, la classe permette di attuare operazioni \texttt{CRUD} sulle tabelle della base di dati.

\subsubsection*{Relazioni d'uso}

\begin{itemize}
	\item \texttt{org.softwaresynthesis.mytalk.server.IMyTalkObject};
	\item \texttt{org.softwaresynthesis.mytalk.server.abook.IGroup};
	\item \texttt{org.softwaresynthesis.mytalk.server.abook.IUserData};
	\item \texttt{org.softwaresynthesis.mytalk.server.dao.util.GetUtil};
	\item \texttt{org.softwaresynthesis.mytalk.server.dao.util.ModifyUtil};
	\item \texttt{org.softwaresynthesis.mytalk.server.dao.util.UtilFactory};
	\item \texttt{org.softwaresynthesis.mytalk.server.message.IMessage};
\end{itemize}

\subsubsection*{Classi estese ed interfacce implementate}

Nessuna relazione evidenziata.

\subsubsection*{Attributi}

\begin{itemize}
	\item{\memberdata{-- manager: ISessionManager}}\\
	Attributo usato per mantenere un riferimento alla singola istanza di SessionManagaer (estensione di \classname{ISessionManagaer}.
	\item{\memberdata{-- factory: UtilFactory}}\\
	Attributo usato per mantenere un riferimento ad un'istanza di UtilFactory, usata per l'esecuzione delle operazioni relative alla manipolazione dei dati.
\end{itemize}

\subsubsection*{Metodi}

\begin{description}
	\item{\method{+ DataPersistanceManager()}}\\
	Costruttore pubblico della classe. Richiama il costruttore \method{DataPersistanceManager(ISessionManager manager, UtilFactory factory} passando come parametri \texttt{SessionMAnager.getIstance()} e una nuova istanza di \classname{UtilFactory}.
	
	\item{\method{\# DataPersistanceManager(manager: ISessionManager, factory: UtilFactory)}}\\
	Costruttore protetto della classe. Inizializza gli attributi della classe con i parametri d'ingresso.
	
	\item{\method{+ delete(object: IMyTalkObject): boolean}}\\
	Metodo per la manipolazione dei dati. Nello specifico permette di eseguire un operazione di eliminazione in una delle tabelle del database (tabella corrispondente all'oggetto object: IMyTalkObject). Il metodo definisce una variabile result di tipo \texttt{boolean} usata per evidenziare la buona riuscita o meno dell'operazione. Di default result è impostata a \inglese{false}. Il metodo crea un istanza di \classname{ModifyUtil} a partire da una chiamata \verb|this.factory.getDeleteUtil(this.manager)|. Quindi si salva in result il valore ritornato dall'operazione \verb|execute(object)| (dove object, parametro d'ingresso del metodo, rappresenta la tabella su cui eseguire l'operazione). \verb|execute(object)| è richiamata a partire dall'istanza \classname{ModifyUtil} creata in precedenza). Il metodo deve terminare ritornando il valore di result.
	
	
	\item{\method{+ insert(object: IMyTalkObject): boolean}}\\
	Metodo per la manipolazione dei dati. Nello specifico permette di eseguire un operazione di inserimento in una delle tabelle del database (tabella corrispondente all'oggetto object: IMyTalkObject). Il metodo definisce una variabile result di tipo \texttt{boolean} usata per evidenziare la buona riuscita o meno dell'operazione. Di default result è impostata a \inglese{false}. Il metodo crea un istanza di \classname{ModifyUtil} a partire da una chiamata \verb|this.factory.getInsertUtil(this.manager)|. Quindi si salva in result il valore ritornato dall'operazione \verb|execute(object)| (dove object, parametro d'ingresso del metodo, rappresenta la tabella su cui eseguire l'operazione). \verb|execute(object)| è richiamata a partire dall'istanza \classname{ModifyUtil} creata in precedenza). Il metodo deve terminare ritornando il valore di result.
	
	\item{\method{+ update(object: IMyTalkObject): boolean}}\\
	Metodo per la manipolazione dei dati. Nello specifico permette di eseguire un operazione di modifica dei campi dati di una delle tabelle del database (tabella corrispondente all'oggetto object: IMyTalkObject). Il metodo definisce una variabile result di tipo \texttt{boolean} usata per evidenziare la buona riuscita o meno dell'operazione. Di default result è impostata a \inglese{false}. Il metodo crea un istanza di \classname{ModifyUtil} a partire da una chiamata \verb|this.factory.getUpdateUtil(this.manager)|. Quindi si salva in result il valore ritornato dall'operazione \verb|execute(object)| (dove object, parametro d'ingresso del metodo, rappresenta la tabella su cui eseguire l'operazione). \verb|execute(object)| è richiamata a partire dall'istanza \classname{ModifyUtil} creata in precedenza). Il metodo deve terminare ritornando il valore di result.
	
	\item{\method{+ getGroup(id: Long): IGroup}}\\
	Metodo usato per l'esecuzione di una query SQL atta ad garantire la restituzione dei dati inerenti il gruppo avente come primary key il valore contenuto in id. Il metodo definisce un istanza di \classname{GetUtil} mediante una chiamata \verb|this.factory.getGroupUtil(this.manager)|. Quindi definisce un oggetto \classname{IGroup} denominato result. Successivamente si definisce una lista \texttt{List<IMyTalkObject>}, da usare per memorizzare temporaneamente i dati ottenuti dall'esecuzione della query:\\
	
	\verb|from Groups as g where g.id = VALORE_id|\\
	
	Tale operazione si ottiene con l'istruzione \verb|select.execute(query)|. Quindi si controlla se la lista cosi popolata contiene almeno un elemento (che per come è stata progettata la query, ossia un interrogazione su chiave primaria, deve essere al più uno). Nel caso si imposta result con una chiamata a metodo get() richiamato dalla lista usata. Il metodo termina restituendo result.
	
		\item{\method{+ getGroup(owner: IUserData): List<IGroup>}}\\
	Metodo usato per l'esecuzione di una query SQL atta ad garantire la restituzione dei gruppi aventi come possessore l'utente owner. Il metodo definisce un istanza di \classname{GetUtil} mediante una chiamata \verb|this.factory.getGroupUtil(this.manager)|. Quindi definisce un oggetto \classname{IGroup} denominato result. Successivamente si definisce una lista \texttt{List<IMyTalkObject>}, da usare per memorizzare temporaneamente i dati ottenuti dall'esecuzione della query:\\
	
	\verb|from Groups as g where g.owner = VALORE_owner|\\
	
	Tale operazione si ottiene con l'istruzione \verb|select.execute(query)|. Quindi si controlla se la lista cosi popolata contiene almeno un elemento. Nel caso si imposta result uguale alla lista usata. Il metodo termina restituendo result.
	
	\item{\method{+ getMessageNewKey(): Long}}\\
	Metodo usato per l'esecuzione di una query SQL atta ad garantire la restituzione della più alta chiave primaria usata nella tabella \texttt{message}. Il metodo restituirà tale valore maggiorato di 1. Il metodo definisce un istanza di \classname{GetUtil} mediante una chiamata \verb|this.factory.getGroupUtil(this.manager)|. Quindi definisce un oggetto \classname{Long} denominato result. Successivamente si definisce un long nominato id, da usare per memorizzare temporaneamente il dato ottenuto dall'esecuzione della query:\\
	
	\verb|max(id) from Messages|\\
	
	Tale operazione si ottiene con l'istruzione \verb|select.execute(query)|. Quindi si controlla se la variabile id contiene alemno un elemento. Nel caso si imposta result uguale alla uguale al contenuto di id + 1. Il metodo termina restituendo result.
	
	\item{\method{+ getMessage(id: Long): IMessage}}\\
	Metodo usato per l'esecuzione di una query SQL atta a garantire la restituzione dei dati inerenti il messaggio avente come primary key il valore contenuto in id. Il metodo definisce un istanza di \classname{GetUtil} mediante una chiamata \verb|this.factory.getGenericUtil(this.manager)|. Quindi definisce un oggetto \classname{IGroup} denominato result. Successivamente si definisce una lista \texttt{List<IMyTalkObject>}, da usare per memorizzare temporaneamente i dati ottenuti dall'esecuzione della query:\\
	
	\verb|from Messages as m where m.id = VALORE_id|\\
	
	Tale operazione si ottiene con l'istruzione \verb|select.execute(query)|. Quindi si controlla se la lista cosi popolata contiene almeno un elemento (che per come è stata progettata la query, ossia un interrogazione su chiave primaria, deve essere al più uno). Nel caso si imposta result con una chiamata a metodo get() richiamato dalla lista usata. Il metodo termina restituendo result.
	
		\item{\method{+ getMessages(receiver: IUserData): List<IMessages>}}\\
	Metodo usato per l'esecuzione di una query SQL atta ad garantire la restituzione dei messaggi aventi come ricevente l'utente receiver. Il metodo definisce un istanza di \classname{GetUtil} mediante una chiamata \verb|this.factory.getGenericUtil(this.manager)|. Quindi definisce un oggetto \classname{IMessages} denominato result. Successivamente si definisce una lista \texttt{List<IMyTalkObject>}, da usare per memorizzare temporaneamente i dati ottenuti dall'esecuzione della query:\\
	
	\verb|from Messages as m where m.receiver = VALORE_receiver|\\
	
	Tale operazione si ottiene con l'istruzione \verb|select.execute(query)|. Quindi si controlla se la lista cosi popolata contiene almeno un elemento. Nel caso si imposta result uguale alla lista usata. Il metodo termina restituendo result.	
	
	\item{\method{+ getUserData(mail: String): IUserData}}\\
	Metodo usato per l'esecuzione di una query SQL atta a garantire la restituzione dei dati inerenti l'utente avente come indirizzo mail il valore contenuto in mail. Il metodo definisce un istanza di \classname{GetUtil} mediante una chiamata \verb|this.factory.getUserDataUtil(this.manager)|. Quindi definisce un oggetto \classname{IUserData} denominato result. Successivamente si definisce una lista \texttt{List<IMyTalkObject>}, da usare per memorizzare temporaneamente i dati ottenuti dall'esecuzione della query:\\
	
	\verb|from UserData as u where u.mail = VALORE_mail|\\
	
	Tale operazione si ottiene con l'istruzione \verb|select.execute(query)|. Quindi si controlla se la lista cosi popolata contiene almeno un elemento (che per come è stata progettata la query, ossia un interrogazione su chiave primaria, deve essere al più uno). Nel caso si imposta result con una chiamata a metodo get() richiamato dalla lista usata. Il metodo termina restituendo result.	
	
	\item{\method{+ getUserData(id: Long): IUserData}}\\
	Metodo usato per l'esecuzione di una query SQL atta a garantire la restituzione dei dati inerenti l'utente avente come chiave primaria il valore contenuto in id. Il metodo definisce un istanza di \classname{GetUtil} mediante una chiamata \verb|this.factory.getUserDataUtil(this.manager)|. Quindi definisce un oggetto \classname{IUserData} denominato result. Successivamente si definisce una lista \texttt{List<IMyTalkObject>}, da usare per memorizzare temporaneamente i dati ottenuti dall'esecuzione della query:\\
	
	\verb|from UserData as u where u.id = VALORE_id|\\
	
	Tale operazione si ottiene con l'istruzione \verb|select.execute(query)|. Quindi si controlla se la lista cosi popolata contiene almeno un elemento (che per come è stata progettata la query, ossia un interrogazione su chiave primaria, deve essere al più uno). Nel caso si imposta result con una chiamata a metodo get() richiamato dalla lista usata. Il metodo termina restituendo result.
	
	\item{\method{+ getUserDatas(mail: String, name: String, surname: String): List<IUserData>}}\\
	Metodo usato per l'esecuzione di una query SQL atta ad garantire la restituzione degli utenti aventi come valore dei campi omonimi ai parametri del metodo, il valore di quest'ultimi. Il metodo definisce un istanza di \classname{GetUtil} mediante una chiamata \\ \verb|this.factory.getUserDataUtil(this.manager)|. Quindi definisce un oggetto \classname{IUserData} denominato result. Successivamente si definisce una lista \texttt{List<IMyTalkObject>}, da usare per memorizzare temporaneamente i dati ottenuti dall'esecuzione della query:\\
	
	\verb|from UserData as u where u.mail like VALORE_receiver|\\
	\verb|or u.name like VALORE_name or u.surname like VALORE_surname|\\
	
	Tale operazione si ottiene con l'istruzione \verb|select.execute(query)|. Quindi si controlla se la lista cosi popolata contiene almeno un elemento. Nel caso si imposta result uguale alla lista usata. Il metodo termina restituendo result.
	
\end{description}

\subsection{Package org.softwaresynthesis.mytalk.server.dao.util}\label{sec:daoUtil}

\classsection{UtilFactory}

\subsubsection*{Funzione}
Classe usata per la creazione delle classi di manipolazione del database.

\subsubsection*{Relazioni d'uso}

\begin{itemize}
		\item{org.softwaresynthesis.mytalk.server.dao.ISessionManager}: usata per definire un riferimento con il sistema centrale per la gestione del database.
\end{itemize}

\subsubsection*{Classi estese ed interfacce implementate}

Nessuna relazione evidenziata.

\subsubsection*{Attributi}

Nessun attributo evidenziare.

\subsubsection*{Metodi}

\begin{description}
	\item{\method{+ getDeleteUtil(manager: ISessionManager): ModifyUtil}}\\
	Metodo usato per ritornare un \inglese{utility} per la cancellazione di un record. Il metodo ritorna \verb|new DeleteUtil(manager)|.
	
	\item{\method{+ getInsertUtil(manager: ISessionManager): ModifyUtil}}\\
	Metodo usato per ritornare una \inglese{utility} per l'inserimento di un record. Il metodo ritorna \verb|new InsertUtil(manager)|.
	
	\item{\method{+ getUpdateUtil(manager: ISessionManager): ModifyUtil}}\\
	Metodo usato per ritornare una \inglese{utility} per l'aggiornamento di un record. Il metodo ritorna \verb|new UpdateUtil(manager)|.
	
	\item{\method{+ getCallUtil(manager: ISessionManager): GetUtil}}\\
	Metodo usato per ritornare una \inglese{utility} per inizializzare correttamente oggetti rappresentanti una chiamata. Il metodo ritorna \verb|new GetCallUtil(manager)|.	
	
	\item{\method{+ getGroupUtil(manager: ISessionManager): GetUtil}}
		Metodo usato per ritornare una \inglese{utility} per inizializzare correttamente oggetti rappresentanti gruppi di utenti. Il metodo ritorna \verb|new GetGroupUtil(manager)|.	
	
	\item{\method{+ getUserDataUtil(manager: ISessionManager): GetUtil}}\\
		Metodo usato per ritornare una \inglese{utility} per inizializzare correttamente oggetti rappresentanti utenti del sistema. Il metodo ritorna \verb|new GetUserDataUtil(manager)|.	

	\item{\method{+ getGenericUtil(manager: ISessionManager): GetUtil}}\\
		Metodo usato per ritornare una \inglese{utility} per estrarre dal database oggetti che non contengono campi dati di tipo ''collezione'' che devono essere inizializzate. Il metodo ritorna \verb|new NotInitialize(manager)|.	
		
\end{description}		

\classsection{ModifyUtil}

\subsubsection*{Funzione}
Classe astratta per definire il comportamento generico di una classe per la manipolazione del database. La classe costituisce un implementazione del \textit{template method}, poiché definisce un iter per l'esecuzione di un algoritmo di modifica, lasciando alle classi figlie il compito di ridefinire il metodo \method{doAction()}.

\subsubsection*{Relazioni d'uso}

\begin{itemize}
		\item \texttt{org.hibernate.Session};
		\item \texttt{org.hibernate.SessionFactory};
		\item \texttt{org.hibernate.Transaction};
		\item \texttt{org.softwaresynthesis.mytalk.server.IMyTalkObject};
		\item \texttt{org.softwaresynthesis.mytalk.server.dao.ISessionManager};
\end{itemize}

\subsubsection*{Classi estese ed interfacce implementate}

Nessuna relazione evidenziata.

\subsubsection*{Attributi}

\begin{itemize}
	\item{\memberdata{\underline{-- manager: ISessionManager}}}\\
	Attributo usato per comunicare con la classe rappresentate il ''cuore'' del sistema di gestione del database.
\end{itemize}

\subsubsection*{Metodi}

\begin{description}
	\item{\method{+ ModifyUtil(manager: ISessionManager)}}\\
	Costruttore pubblico della classe, usato per inizializzare \memberdata{manager} con il parametro d'ingresso del metodo.
	
	\item{\method{+ execute(object: IMyTalkObject): boolean}}\\
	Metodo che definisce l'iter per l'esecuzione dell'operazione di manipolazione del database. Il metodo definisce quelle procedure comuni a tutte le operazioni di modifica. Inizialmente il metodo definisce i seguenti oggetti:
	\begin{itemize}
		\item boolean result = false;
		\item Session session = null;
		\item SessionFactory factory = null;
		\item Transaction transaction = null;
	\end{itemize}
	
Quindi viene definito un blocco try-catch. Nel blocco try, si salva in factory un riferimento a \classname{SessionManager}. Quindi in session si riporta la sessione ottenuta da una chiamata \verb|factory.openSession()|. Si definisce l'apertura di una transizione transaction mediante una chiamata \verb|session.beginTransaction()|. Il metodo procede eseguendo il metodo \method{doAction()} passando come parametro d'ingresso session e object. Si esegue il commit sulla transizione e si imposta result a \inglese{true}. Nel blocco catch si esegue un'operazione di rollback e si imposta result a \inglese{false}.

Il metodo termina chiudendo la sessione e ritornando il contenuto di result.

	\item{\method{+ \{abstract\} doAction(session: ISessionManager, object: IMyTalkObject): void}}\\
	Metodo astratto non definito in questa classe. Le classi figlie di \classname{ModifyUtil} ridefiniranno questo metodo.

\end{description}

\classsection{InsertUtil}

\subsubsection*{Funzione}
Classe figlia di \classname{ModifyUtil}. Ridefinisce il metodo \method{doAction} per eseguire un operazione di insert in una tabella del database.

\subsubsection*{Relazioni d'uso}

\begin{itemize}
		\item \texttt{org.hibernate.Session};
		\item \texttt{org.softwaresynthesis.mytalk.server.IMyTalkObject};
		\item \texttt{org.softwaresynthesis.mytalk.server.dao.ISessionManager};
\end{itemize}

\subsubsection*{Classi estese ed interfacce implementate}

\begin{itemize}
	\item \classname{dao.util.ModifyUtil}: classe astratta estesa.
\end{itemize}

\subsubsection*{Attributi}

Nessun attributo evidenziato.

\subsubsection*{Metodi}

\begin{description}
	\item{\method{+ InsertUtil(manager: ISessionManager)}}\\
	Costruttore pubblico della classe, usato per inizializzare \memberdata{manager} con il parametro d'ingresso del metodo. Per farlo richiama il costruttore della classe padre.

	\item{\method{+ \{final\} doAction(session: ISessionManager, object: IMyTalkObject): void}}\\
	Metodo usato per eseguire un operazione di insert nel database. Il metodo esegue l'unica operazione \verb|session.save(object)|.

\end{description}

\classsection{DeleteUtil}

\subsubsection*{Funzione}
Classe figlia di \classname{ModifyUtil}. Ridefinisce il metodo \method{doAction} per eseguire un operazione di delete in una tabella del database.

\subsubsection*{Relazioni d'uso}

\begin{itemize}
		\item \texttt{org.hibernate.Session};
		\item \texttt{org.softwaresynthesis.mytalk.server.IMyTalkObject};
		\item \texttt{org.softwaresynthesis.mytalk.server.dao.ISessionManager};
\end{itemize}

\subsubsection*{Classi estese ed interfacce implementate}

\begin{itemize}
	\item \classname{dao.util.ModifyUtil}: classe astratta estesa.
\end{itemize}

\subsubsection*{Attributi}

Nessun attributo evidenziato.

\subsubsection*{Metodi}

\begin{description}
	\item{\method{+ DeleteUtil(manager: ISessionManager)}}\\
	Costruttore pubblico della classe, usato per inizializzare \memberdata{manager} con il parametro d'ingresso del metodo. Per farlo richiama il costruttore della classe padre.

	\item{\method{+ \{final\} doAction(session: ISessionManager, object: IMyTalkObject): void}}\\
	Metodo usato per eseguire un operazione di delete nel database. Il metodo esegue l'unica operazione \verb|session.delete(object)|.

\end{description}

\classsection{UpdateUtil}

\subsubsection*{Funzione}
Classe figlia di \classname{ModifyUtil}. Ridefinisce il metodo \method{doAction} per eseguire un operazione di aggiornamento (intesa come modifica del valore di alcuni campi dati) in una tabella del database.

\subsubsection*{Relazioni d'uso}

\begin{itemize}
		\item \texttt{org.hibernate.Session};
		\item \texttt{org.softwaresynthesis.mytalk.server.IMyTalkObject};
		\item \texttt{org.softwaresynthesis.mytalk.server.dao.ISessionManager};
\end{itemize}

\subsubsection*{Classi estese ed interfacce implementate}

\begin{itemize}
	\item \classname{dao.util.ModifyUtil}: classe astratta estesa.
\end{itemize}

\subsubsection*{Attributi}

Nessun attributo evidenziato.

\subsubsection*{Metodi}

\begin{description}
	\item{\method{+ UpdateUtil(manager: ISessionManager)}}\\
	Costruttore pubblico della classe, usato per inizializzare \memberdata{manager} con il parametro d'ingresso del metodo. Per farlo richiama il costruttore della classe padre.

	\item{\method{+ \{final\} doAction(session: ISessionManager, object: IMyTalkObject): void}}\\
	Metodo usato per eseguire un operazione di update nel database. Il metodo esegue l'unica operazione \verb|session.update(object)|.

\end{description}

\classsection{GetUtil}

\subsubsection*{Funzione}
Classe astratta per definire il comportamento generico di una classe per l'esecuzione di ''interrogazioni'' (query di tipo \texttt{select}) al database. La classe costituisce un implementazione del \textit{template method}, poiché definisce un iter per l'esecuzione di un algoritmo, lasciando alle classi figlie il compito di ridefinire il metodo \method{doInitialize()}.

\subsubsection*{Relazioni d'uso}

\begin{itemize}
		\item \texttt{java.util.List};
		\item \texttt{org.hibernate.Query};
		\item \texttt{org.hibernate.Session};
		\item \texttt{org.hibernate.SessionFactory};
		\item \texttt{org.hibernate.Transaction};
		\item \texttt{org.softwaresynthesis.mytalk.server.IMyTalkObject};
		\item \texttt{org.softwaresynthesis.mytalk.server.dao.ISessionManager};
\end{itemize}

\subsubsection*{Classi estese ed interfacce implementate}

Nessuna relazione evidenziata.

\subsubsection*{Attributi}

\begin{itemize}
	\item{\memberdata{\underline{-- manager: ISessionManager}}}\\
	Attributo usato per comunicare con la classe rappresentate il ''cuore'' del sistema di gestione del database.
\end{itemize}

\subsubsection*{Metodi}

\begin{description}
	\item{\method{+ GetUtil(manager: ISessionManager)}}\\
	Costruttore pubblico della classe, usato per inizializzare \memberdata{manager} con il parametro d'ingresso del metodo.
	
	\item{\method{+ execute(query: String): List<IMyTalkObject>}}\\
	Metodo che definisce l'iter per l'esecuzione dell'operazione d'interrogazione del database. Il metodo definisce quelle procedure comuni a tutte le operazioni di lettura dei dati. Inizialmente il metodo definisce i seguenti oggetti:
	\begin{itemize}
		\item Query hqlQuery = null;
		\item List<IMyTalkObject> collection = null;
		\item Session session = null;
		\item SessionFactory factory = null;
		\item Transaction transaction = null;
	\end{itemize}
	
Quindi viene definito un blocco try-catch. Nel blocco try, si salva in factory un riferimento a \classname{SessionManager}. Quindi in session si riporta la sessione ottenuta da una chiamata \verb|factory.openSession()|. Si definisce la query d'interrogazione mediante una chiamata \verb|session.createQuery(query)|. Si definisce l'apertura di una transizione transaction mediante una chiamata \verb|session.beginTransaction()|. Il metodo procede prima impostando collection a \verb|(List<IMyTalkObject>)hqlQuery.list();|, e poi eseguendo il metodo \method{doInitialize()} passando come parametro d'ingresso collection. Si esegue il commit sulla transizione e si imposta result a \inglese{true}. Nel blocco catch si esegue un'operazione di rollback e si imposta collection a \inglese{null}.

Il metodo termina chiudendo la sessione e ritornando il contenuto di collection.

	\item{\method{+ uniqueResult(query: String): Long}}\\
	Metodo che definisce l'iter per l'esecuzione dell'operazione d'interrogazione del database (quelle che prevedono un unico risultato). Il metodo definisce quelle procedure comuni a tutte le operazioni di lettura dei dati. Inizialmente il metodo definisce i seguenti oggetti:
	\begin{itemize}
		\item Query hqlQuery = null;
		\item Long result = null;
		\item Session session = null;
		\item SessionFactory factory = null;
		\item Transaction transaction = null;
	\end{itemize}
	
Quindi viene definito un blocco try-catch. Nel blocco try, si salva in factory un riferimento a \classname{SessionManager}. Quindi in session si riporta la sessione ottenuta da una chiamata \verb|factory.openSession()|. Si definisce la query d'interrogazione mediante una chiamata \verb|session.createQuery(query)|. Si definisce l'apertura di una transizione transaction mediante una chiamata \verb|session.beginTransaction()|. Il metodo procede prima impostando result a \verb|(Long)hqlQuery.uniqueResult()|, e poi eseguendo l'operazione di commit a partire dall'oggetto transaction. Nel blocco catch si esegue un'operazione di rollback e si imposta result a \inglese{null}.

Il metodo termina chiudendo la sessione e ritornando il contenuto di result.

	\item{\method{+ \{abstract\} doInitialize(collection: List<IMyTalkObject>): void}}\\
	Metodo astratto non definito in questa classe. Le classi figlie di \classname{GetUtil} ridefiniranno questo metodo.

\end{description}

\classsection{GetCallUtil}

\subsubsection*{Funzione}
Classe figlia di \classname{GetUtil}. Rappresenta l'operazione necessaria per ottenere le chiamate registrate nel database.

\subsubsection*{Relazioni d'uso}

\begin{itemize}
		\item \texttt{java.util.Iterator};
		\item \texttt{import java.util.List};
		\item \texttt{import org.hibernate.Hibernate};
		\item \texttt{org.hibernate.Session};
		\item \texttt{org.softwaresynthesis.mytalk.server.IMyTalkObject};
		\item \texttt{org.softwaresynthesis.mytalk.server.dao.ISessionManager};
\end{itemize}

\subsubsection*{Classi estese ed interfacce implementate}

\begin{itemize}
	\item \classname{dao.util.GetUtil}: classe astratta estesa.
\end{itemize}

\subsubsection*{Attributi}

Nessun attributo evidenziato.

\subsubsection*{Metodi}

\begin{description}
	\item{\method{+ GetCallUtil(manager: ISessionManager)}}\\
	Costruttore pubblico della classe, usato per inizializzare \memberdata{manager} con il parametro d'ingresso del metodo. Per farlo richiama il costruttore della classe padre.

	\item{\method{+ \{final\} doInitialize(collection: List<IMyTalkObject>): void}}\\
	Il metodo definisce un istanza di \classname{ICall} call inizializzandola a null. Quindi procede definendo un iteratore Iterator<IMyTalkObject> a partire da collection. All'interno di un ciclo while che itera per ogni elemento dell'iteratore, si sovrascrive il contenuto di call con quanto estratto dall'iteratore (uso del metodo \texttt{next()}), e poi deve essere eseguita l'istruzione:\\
	
	\verb|Hibernate.initialize(call.getCalls());|

\end{description}

\classsection{GetGroupUtil}

\subsubsection*{Funzione}
Classe figlia di \classname{GetUtil}. Rappresenta l'operazione necessaria per ottenere i gruppi presenti nel database.

\subsubsection*{Relazioni d'uso}

\begin{itemize}
		\item \texttt{java.util.Iterator};
		\item \texttt{import java.util.List};
		\item \texttt{import org.hibernate.Hibernate};
		\item \texttt{org.hibernate.Session};
		\item \texttt{org.softwaresynthesis.mytalk.server.IMyTalkObject};
		\item \texttt{org.softwaresynthesis.mytalk.server.dao.ISessionManager};
\end{itemize}

\subsubsection*{Classi estese ed interfacce implementate}

\begin{itemize}
	\item \classname{dao.util.GetUtil}: classe astratta estesa.
\end{itemize}

\subsubsection*{Attributi}

Nessun attributo evidenziato.

\subsubsection*{Metodi}

\begin{description}
	\item{\method{+ GetGroupUtil(manager: ISessionManager)}}\\
	Costruttore pubblico della classe, usato per inizializzare \memberdata{manager} con il parametro d'ingresso del metodo. Per farlo richiama il costruttore della classe padre.

	\item{\method{+ \{final\} doInitialize(collection: List<IMyTalkObject>): void}}\\
	Il metodo definisce un istanza di \classname{IGroup} group inizializzandola a null. Quindi procede definendo un iteratore Iterator<IMyTalkObject> a partire da collection. All'interno di un ciclo while che itera per ogni elemento dell'iteratore, si sovrascrive il contenuto di group con quanto estratto dall'iteratore (uso del metodo \texttt{next()}), e poi deve essere eseguita l'istruzione:\\
	
	\verb|Hibernate.initialize(group.getAddressBook());|

\end{description}

\classsection{GetUserDataUtil}

\subsubsection*{Funzione}
Classe figlia di \classname{GetUtil}. Rappresenta l'operazione necessaria per ottenere gli utenti presenti nel database.

\subsubsection*{Relazioni d'uso}

\begin{itemize}
		\item \texttt{java.util.Iterator};
		\item \texttt{import java.util.List};
		\item \texttt{import org.hibernate.Hibernate};
		\item \texttt{org.hibernate.Session};
		\item \texttt{org.softwaresynthesis.mytalk.server.IMyTalkObject};
		\item \texttt{org.softwaresynthesis.mytalk.server.dao.ISessionManager};
\end{itemize}

\subsubsection*{Classi estese ed interfacce implementate}

\begin{itemize}
	\item \classname{dao.util.GetUtil}: classe astratta estesa.
\end{itemize}

\subsubsection*{Attributi}

Nessun attributo evidenziato.

\subsubsection*{Metodi}

\begin{description}
	\item{\method{+ GetUserDataUtil(manager: ISessionManager)}}\\
	Costruttore pubblico della classe, usato per inizializzare \memberdata{manager} con il parametro d'ingresso del metodo. Per farlo richiama il costruttore della classe padre.

	\item{\method{+ \{final\} doInitialize(collection: List<IMyTalkObject>): void}}\\
	Il metodo definisce un istanza di \classname{IUserData} user inizializzandola a null. Quindi procede definendo un iteratore Iterator<IMyTalkObject> a partire da collection. All'interno di un ciclo while che itera per ogni elemento dell'iteratore, si sovrascrive il contenuto di user con quanto estratto dall'iteratore (uso del metodo \texttt{next()}), e poi devono essere eseguite le istruzioni:\\
	
	\verb|Hibernate.initialize(user.getAddressBook());|\\
	\verb|Hibernate.initialize(user.getMessages());|\\
	\verb|Hibernate.initialize(user.getCalls());|

\end{description}

\classsection{NotInitialize}

\subsubsection*{Funzione}
Classe figlia di \classname{GetUtil}.

\subsubsection*{Relazioni d'uso}

\begin{itemize}
		\item \texttt{org.hibernate.Session};
		\item \texttt{org.softwaresynthesis.mytalk.server.IMyTalkObject};
		\item \texttt{org.softwaresynthesis.mytalk.server.dao.ISessionManager};
\end{itemize}

\subsubsection*{Classi estese ed interfacce implementate}

\begin{itemize}
	\item \classname{dao.util.GetUtil}: classe astratta estesa.
\end{itemize}

\subsubsection*{Attributi}

Nessun attributo evidenziato.

\subsubsection*{Metodi}

\begin{description}
	\item{\method{+ DeleteUtil(manager: ISessionManager)}}\\
	Costruttore pubblico della classe, usato per inizializzare \memberdata{manager} con il parametro d'ingresso del metodo. Per farlo richiama il costruttore della classe padre.

	\item{\method{+ \{final\} doInitialize(collection: List<IMyTalkObject>): void}}\\
	Metodo dal corpo vuoto.

\end{description}

\subsection{Package org.softwaresynthesis.mytalk.server.abook}\label{sec:abook}

\classsection{IUserData}

\subsubsection*{Funzione}
Interfaccia rappresentante il comportamento di un generico utente del sistema. L'interfaccia dovrà definire dei metodi di tipo \textit{get} e \textit{set} per i dati d'interesse.

\subsubsection*{Relazioni d'uso}

\begin{itemize}
	\item \classname{AddressBookEntry}: l'interfaccia definisce dei metodi per la manipolazione di dati \classname{AddressBookEntry}.
\end{itemize}

\subsubsection*{Classi estese ed interfacce implementate}
\begin{itemize}
	\item \texttt{org.softwaresynthesis.mytalk.server.IMyTalkObject}: interfaccia da estendere. Ogni oggetto che implementerà l'interfaccia \classname{IUserData} dovrà essere in grado di convertire il proprio contenuto informativo in formato \textit{Json}.
\end{itemize}


\subsubsection*{Metodi}
\begin{description}
	\item{\method{+ getId(): Long}}\\
	Restituisce l'identificatore univoco di uno \classname{IUserData}.
	\item{\method{+ getEmail(): String}}\\
	Restituisce l'indirizzo e-mail con cui uno \classname{IUserData} si è registrato nel sistema \caName.
	\item{\method{+ setEmail(mail: String): void}}\\ 
	Imposta l'indirizzo e-mail con cui si registra nel sistema \caName uno \classname{IUserData}.
	\item{\method{+ getPassword(): String}}\\
	Restituisce la password di accesso al sistema \caName di uno \classname{IUserData}.
	\item{\method{+ setPassword(password: String): void}}\\
	Imposta la password di accesso al sistema di uno \classname{IUserData}.
	\item{\method{+ getQuestion(): String}}\\
	Restituisce la domanda segreta, scelta da uno \classname{IUserData}, per il recupero della \textit{password} smarrita di accesso al sistema \caName.
	\item{\method{+ setQuestion(question: String): void}}\\
	Imposta la domanda segreta, scelta da uno \classname{IUserData}, per il recupero della \textit{password} smarrita di accesso al sistema \caName.
	\item{\method{+ getAnswer(): String}}\\
	Restituisce la risposta alla domanda per il recupero della \textit{password} smarrita di accesso al sistema \caName.
	\item{\method{+ setAnswer(answer: String): void}}\\
	Imposta la risposta alla domanda segreta per il recupero della \textit{password} di accesso al sistema \caName.
	\item{\method{+ getName(): String}}\\
	Restituisce il nome di uno \classname{IUserData}.
	\item{\method{+ setName(name: String): void}}\\
	Imposta il nome di uno \classname{IUserData}.
	\item{\method{+ getSurname(): String}}\\
	Restituisce il cognome di uno \classname{IUserData}.
	\item{\method{+ setSurname(surname: String): void}}\\
	Imposta il cognome di uno \classname{IUserData}.
	\item{\method{+ getPicturePath(): String}}\\
	Restituisce una stringa con il percorso dell'immagine del profilo di uno \classname{IUserData}.
	\item{\method{+ setPicturePath(path: String): void}}\\
	Imposta il percorso dell'immagine profilo di uno \classname{IUserData}.
	\item{\method{+ getAddressBook(): Set<AddressBookEntry>}}\\
	Metodo che ritorna il la rubrica dell'utente sotto forma d'insieme di \classname{AddressBookEntry}.
	\item{\method{+ addAddressBookEntry(entry: AddressBookEntry): void}}\\
	Metodo usato per aggiungere una nuova \classname{AddressBookEntry} all'insieme di \inglese{entry} che costituisce la rubrica utente.
\end{description}

\classsection{IGroup}

\subsubsection*{Funzione}
Interfaccia rappresentante un gruppo di una rubrica utente del sistema \caName.

\subsubsection*{Relazioni d'uso}

Nessuna relazione evidenziata.

\subsubsection*{Classi estese ed interfacce implementate}
\begin{itemize}
		\item \texttt{org.softwaresynthesis.mytalk.server.IMyTalkObject}: interfaccia da estendere. Ogni oggetto che implementerà l'interfaccia \classname{IGroup} dovrà essere in grado di convertire il proprio contenuto informativo in formato \textit{Json}.
\end{itemize}

\subsubsection*{Metodi}
\begin{description}
	\item{\method{+ getId(): Long}}\\
	Restituisce l'identificativo univoco di uno gruppo di una rubrica utente.
	\item{\method{+ getName(): String}}\\
	Restituisce il nome di un gruppo di una rubrica utente.
	\item{\method{+ setName(name: String): void}}\\
	Imposta il nome di un gruppo di una rubrica utente.
\end{description}

\classsection{IAddressBookEntry}

\subsubsection*{Funzione}
Interfaccia rappresentante una \textit{entry} di una rubrica utente del sistema mytalk.

\subsubsection*{Relazioni d'uso}
\begin{itemize}
	\item \classname{IUserData}: l'interfaccia \classname{IAddressBookEntry} definisce più metodi che restituiscono oggetti aventi tipo di ritorno \classname{IUserData}, essi sono i metodi \inglese{get} per ottenere il ``proprietario'' della rubrica e per ottenere l'utente registrato nella rubrica. Analogamente \classname{IUserData} viene usato come parametro d'ingresso per i metodi \inglese{set} collegati ai metodi già citati.
\end{itemize}

\subsubsection*{Classi estese ed interfacce implementate}
\begin{itemize}
		\item \texttt{org.softwaresynthesis.mytalk.server.IMyTalkObject}: interfaccia da estendere. Ogni oggetto che implementerà l'interfaccia \classname{IAddressBookEntry} dovrà essere in grado di convertire il proprio contenuto informativo in formato \textit{Json}.
\end{itemize}

\subsubsection*{Metodi}
\begin{description}
	\item{\method{+ getId(): Long}}\\
	Restituisce l'identificativo univoco di una \inglese{entry} di una rubrica utente del sistema \caName.
	\item{\method{+ getEntry(): IUserData}}\\
	Restituisce un istanza di un oggetto avente tipo \classname{IUserData} rappresentante un contatto della rubrica.
	\item{\method{+ setEntry(contact: IUserData): void}}\\
	Imposta l'utente \classname{IUserData} (passato come parametro d'ingresso) come contatto della rubrica.
	\item{\method{+ getGroup(): IGroup}}\\
	Restituisce il gruppo a cui appartiene lo \classname{IUserData} registrato nella rubrica.
	\item{\method{+ setGroup(group: IGroup): void}}\\
	Imposta il gruppo di appartenenza dello \classname{IUserData} registrato nella rubrica.
	\item{\method{+ getOwner(): IUserData}}\\
	Restituisce lo \classname{IUserData} possessore dell'\inglese{entry} corrente della rubrica
	\item{\method{+ setOwner(owner: IUserData ): void}}\\
	Imposta l'utente \classname{IUserData} possessore della \inglese{entry} della rubrica.
\end{description}

\classsection{UserData}

\subsubsection*{Funzione}
Implementazione dell'interfaccia \classname{IUserData}. Un istanza della classe dovrà rappresentare un generico utente del sistema definendone gli attributi e i metodi per impostare ed ottenere il contenuto dei medesimi.

\subsubsection*{Relazioni d'uso}

Nessuna relazione evidenziata

\subsubsection*{Classi estese ed interfacce implementate}
\begin{itemize}
		\item{IUserData}: interfaccia da implementare.
\end{itemize}

\subsubsection*{Attributi}

\begin{itemize}
	\item{\memberdata{-- id: long}}
	Attributo che definisce il codice identificativo con il quale l'utente è registrato nel database del sistema.
	\item{\memberdata{-- mail: String}}
	Attributo che definisce l'indirizzo e-mail con il quale l'utente si è registrato nel sistema.
	\item{\memberdata{-- password: String}}
	Attributo che definisce la password per il \inglese{login} dell'utente nel sistema.
	\item{\memberdata{-- question: String}}
	Attributo che definisce la domanda segreta usata dall'utente in caso di smarrimento della password.
	\item{\memberdata{-- answer: String}}
	Attributo che definisce la risposta alla domanda segreta definita nell'attributo \memberdata{question}.
	\item{\memberdata{-- name: String}}
	Attributo che definisce il nome dell'utente.
	\item{\memberdata{-- surname: String}}
	Attributo che definisce il cognome dell'utente.
	\item{\memberdata{-- path: String}}
	Attributo che definisce il percorso (sul server) in cui è memorizzata l'immagine del profilo dell'utente.
	\item{\memberdata{-- addressBook: Set<AddressBookEntry>}}
	Attributo che definisce l'insieme di \classname{AddressBookEntry} che costituiscono la rubrica dell'utente.
	
\end{itemize}

\subsubsection*{Metodi}

\begin{description}
	\item{\method{+ getId(): Long}}\\
	Restituisce l'identificatore univoco di un utente, ritornando l'attributo \memberdata{id}.
	\item{\method{\# setId(id: long): void}}\\
	Imposta l'indirizzo \textit{id} con cui l'utente si registra nel sistema \caName. Il metodo sovrascrive il contenuto dell'attributo \memberdata{id} con il valore tipo \texttt{long} ricevuto come parametro d'ingresso. ritornando l'attributo \memberdata{id}.
	\item{\method{+ getEmail(): String}}\\
	Restituisce l'indirizzo e-mail con cui uno l'utente si è registrato nel sistema \caName, ritornando il contenuto dell'attributo \memberdata{mail}.
	\item{\method{+ setEmail(mail: String): void}}\\ 
	Imposta l'indirizzo e-mail con cui l'utente si registra nel sistema \caName. Il metodo non fa altro che sovrascrivere il contenuto dell'attributo \memberdata{mail} con il valore tipo \texttt{String} ricevuto come parametro d'ingresso.
	\item{\method{+ getPassword(): String}}\\
	Restituisce la password dell'utente, ritornando il valore contenuto nell'attributo \memberdata{password}.
	\item{\method{+ setPassword(password: String): void}}\\
	Imposta la password di accesso al sistema, sovrascrivendo il contenuto dell'attributo \memberdata{password} con il valore di tipo \texttt{String} ricevuto come parametro d'ingresso.
	\item{\method{+ getQuestion(): String}}\\
	Restituisce la domanda segreta, scelta dall'utente, per il recupero della password di accesso al sistema \caName. Nello specifico il metodo restituisce il contenuto dell'attributo \memberdata{question}.
	\item{\method{+ setQuestion(question: String): void}}\\
	Imposta la domanda segreta da inserire in caso di smarrimento della password. Il metodo sovrascrive il contenuto dell'attributo \memberdata{question} con il valore tipo \texttt{String} ricevuto come parametro d'ingresso
	\item{\method{+ getAnswer(): String}}\\
	Restituisce la risposta alla domanda per il recupero della password (smarrita) di accesso al sistema \caName. Il metodo ritorna il contenuto dell'attributo \memberdata{answer}.
	\item{\method{+ setAnswer(answer: String): void}}\\
	Imposta la risposta alla domanda segreta per il recupero della password. Il metodo sovrascrive il contenuto dell'attributo \memberdata{answer} con il valore tipo \texttt{String} passato come parametro d'ingresso.
	\item{\method{+ getName(): String}}\\
	Restituisce il nome dell'utente ritornando il contenuto dell'attributo \memberdata{name}.
	\item{\method{+ setName(name: String): void}}\\
	Imposta il nome dell'utente sovrascrivendo il contenuto dell'attributo \memberdata{name} con il valore tipo \texttt{String} passato al metodo come parametro d'ingresso.
	\item{\method{+ getSurname(): String}}\\
	Restituisce il cognome dell'utente restituendo il contenuto dell'attributo \memberdata{surname}.
	\item{\method{+ setSurname(surname: String): void}}\\
	Imposta il cognome dell'utente sovrascrivendo il contenuto dell'attributo \memberdata{surnamename} con il valore tipo \texttt{String} passato al metodo come parametro d'ingresso.
	\item{\method{+ getPicturePath(): String}}\\
	Restituisce una stringa con il percorso dell'immagine del profilo dell'utente, restituendo il contenuto dell'attributo \memberdata{path}.
	\item{\method{+ setPicturePath(path: String): void}}\\
	Imposta il percorso dell'immagine profilo di un utente, sovrascrivendo il contenuto dell'attributo \memberdata{path} con il valore tipo \texttt{String} passato al metodo come parametro d'ingresso.
	\item{\method{+ getState(): State}}\\
	Restituisce lo stato in cui si trova l'utente, ritornando il contenuto dell'attributo \memberdata{state}.
	\item{\method{+ setState(state: State): void}}\\
	Imposta lo stato in cui si trova l'utente, sovrascrivendo il contenuto dell'attributo \memberdata{state} con il valore ricevuto come parametro d'ingresso.
	\item{\method{+ getAddressBook(): Set<AddressBookEntry>}}\\
	Metodo che ritorna il contenuto di \memberdata{addressBook}.
	\item{\method{+ addAddressBookEntry(entry: AddressBookEntry): void}}\\
	Metodo usato per aggiungere ad \memberdata{addressBook} una nuova \classname{AddressBookEntry} passata come parametro d'ingresso.
	\item{\method{+ toJson(): String}}\\
	Metodo usato per ritornare il contenuto di un istanza di \classname{UserData} sotto forma di stringa formattata in \textit{Json}. La stringa ritornata deve corrispondere al seguente formato:\\\\
	\verb|{name:"mio_nome",surname:"mio_cognome",email:"mia_mail"|\\\verb|,picturePath:"mia_immagine",id:"mio_id"}|\\
	
	dove i valori tra virgolette rappresentano il contenuto dei rispettivi campi dati contenuti nella classe.
\end{description}

\classsection{Group}

\subsubsection*{Funzione}
Implementazione dell'interfaccia \classname{IGroup}.

\subsubsection*{Relazioni d'uso}

Nessuna relazione evidenziata.

\subsubsection*{Classi estese ed interfacce implementate}
\begin{itemize}
	\item \classname{IGroup}: interfaccia d'implementazione.
\end{itemize}


\subsubsection*{Attributi}

\begin{itemize}
	\item{\memberdata{-- id: long}}
	Attributo del codice identificativo del gruppo.
	\item{\memberdata{-- name: String}}:
	Attributo del nome del gruppo.
\end{itemize}

\subsubsection*{Metodi}

\begin{description}
	\item{\method{+ getId(): Long}}\\
	Restituisce l'identificativo univoco di uno gruppo di una rubrica utente.
	\item{\method{+ getName(): String}}\\
	Restituisce il nome di un gruppo di una rubrica utente.
	\item{\method{+ setName(name: String): void}}\\
	Imposta il nome di un gruppo di una rubrica utente.
	\item{\method{+ toJson(): String}}\\
	Metodo usato per ritornare il contenuto di un istanza di \classname{Group} sotto forma di stringa formattata in \textit{Json}. La stringa ritornata deve corrispondere al seguente formato:\\\\
	\verb|{id:"mio_id",name:"mio_nome"}|\\
	
	dove i valori tra virgolette rappresentano il contenuto dei rispettivi campi dati contenuti nella classe.
\end{description}

\classsection{AddressBookEntry}

\subsubsection*{Funzione}
Implementazione dell'interfaccia IAddressBookEntry.

\subsubsection*{Relazioni d'uso}

\begin{itemize}
	\item \classname{IUserData}: usata per definire gli attributi destinati a identificare il possessore dell'istanza di \classname{AddressBookEntry} e il relativo contatto in essa registrato.
\end{itemize}

\subsubsection*{Classi estese ed interfacce implementate}
\begin{itemize}
	\item \classname{IAddressBookEntry}: interfaccia d'implementazione della classe.
\end{itemize}

\subsubsection*{Attributi}

\begin{itemize}
	\item{\memberdata{-- id: long}}
	Attributo del codice identificativo della classe.
	\item{\memberdata{-- group: IGroup}}
	Attributo destinato a identificare il gruppo a cui appartiene il contatto \classname{IUserName} registrato nella classe. Si ricorda che il contatto può anche non appartenere ad alcun gruppo.
	\item{\memberdata{-- contact: IUserData}}
	Attributo destinato ad identificare il contatto registrato nell'istanza di \classname{AddressBookEntry}.
	\item{\memberdata{-- owner: IUserData}}:
	Attributo destinato ad identificare il possessore dell'istanza di \classname{AddressBookEntry}.
	\item{\memberdata{-- blocked: boolean}}
	Attributo booleano necessario per bloccare il contatto. Tale blocco avviene impostando l'attributo a \textit{true}.
\end{itemize}


\subsubsection*{Metodi}

\begin{description}
	\item{\method{+ getId(): Long}}\\
	Restituisce l'identificativo univoco della \inglese{entry} di una rubrica utente del sistema \caName, nello specifico il contenuto dell'attributo \memberdata{id}.
	\item{\method{+ getEntry(): IUserData}}\\
	Restituisce il contenuto dell'attributo \memberdata{contact}.
	\item{\method{+ setEntry(contact: IUserData): void}}\\
	Imposta l'utente \classname{IUserData} (passato come parametro d'ingresso) come contatto della rubrica, nello specifico il contenuto dell'attributo \memberdata{contact}.
	\item{\method{+ getGroup(): IGroup}}\\
	Restituisce il gruppo a cui appartiene lo \classname{IUserData} registrato nella rubrica, nello specifico il contenuto dell'attributo \memberdata{group}.
	\item{\method{+ setGroup(group: IGroup): void}}\\
	Imposta il gruppo di appartenenza dello \classname{IUserData} registrato nella rubrica, nello specifico il contenuto dell'attributo \memberdata{group}.
	\item{\method{+ getOwner(): IUserData}}\\
	Restituisce lo \classname{IUserData} possesore di questa \inglese{entry} della rubrica, nello specifico il contenuto dell'attributo \memberdata{owner}.
	\item{\method{+ setOwner(owner: IUserData ): void}}\\
	Imposta l'utente \classname{IUserData} possessore della \textit{entry} della rubrica, nello specifico il contenuto dell'attributo \memberdata{owner}.
	\item{\method{+ toJson(): String}}\\
	Metodo usato per ritornare il contenuto di un istanza di \classname{AddressBookEntry} sotto forma di stringa formattata in \textit{Json}. La stringa ritornata deve corrispondere al seguente formato:\\\\
	\verb|{id:"mio_id",contact:"mio_contatto",group:"mia_gruppo",blocked:"bloccato"}|\\
	
	dove i valori tra virgolette rappresentano il contenuto dei rispettivi campi dati contenuti nella classe.

\end{description}

\subsection{Package org.softwaresynthesis.mytalk.server.abook.controller}\label{sec:call}

%TODO da verificare

\classsection{AddContactController}

\subsubsection*{Funzione}
\textit{Controller} che ha il compito di aggiungere alla rubrica un nuovo contatto.

\subsubsection*{Relazioni d'uso}

\begin{itemize}
	\item \texttt{java.io.IOException}: eccezione richiamabile dal metodo \method{doAction()}.
	\item \texttt{java.io.PrintWriter}: classe istanziata all'interno del metodo \method{doAction()}. Usata per scrivere l'output del \inglese{controller}.
	\item \texttt{javax.servlet.ServletException}: eccezione sollevabile dai metodi \method{doAction()}.
	\item \texttt{javax.servlet.http.HttpServletRequest}: classe usata per interagire con le richieste \texttt{ajax} inoltrate dal client.
	\item \texttt{javax.servlet.http.HttpServletResponse}: classe usata per interagire con le richieste \texttt{ajax} inoltrate dal client.
	\item \classname{org.softwaresynthesis.mytalk.server.dao.DataPersistenceManager}: classe usata per comunicare tramite Hibernate con la tabella UserData della base di dati.
\end{itemize}

\subsubsection*{Classi estese ed interfacce implementate}
\begin{itemize}
	\item \texttt{server.AbstractController}: classe estesa.
\end{itemize}

\subsubsection*{Attributi}

Nessun attributo evidenziato

\subsubsection*{Metodi}
\begin{description}
	\item{\method{\# doAction(request HttpServletRequest, response HttpServletResponse): void}}\\	
	Metodo usato per eseguire la richiesta di aggiungere un nuovo contatto all'interno della rubrica. Il flusso principale inizia  con la creazione di un oggetto \classname{dao.DataPersistanceManager} avente nome UserDAO. Quindi viene aperto un blocco try catch in cui si salva in un istanza di tipo \texttt{HttpSession}, l'oggetto ritornato da una chiamata \verb|request.getSession(false)|. I passi successivi sono (nell'ordine):
	\begin{itemize}
		\item salvare in attributo di tipo long, contactID, l'id dell'utente. Tale id dovrà essere ottenuto con una chiamata:\\
		\verb|request.getParameter("contactId")|;\\
		
		\item si crea un istanza di tipo \classname{IUserData} a partire da una chiamata a metodo:\\
		\verb|session.getAttribute("user")|;
		
		\item si crea un istanza di tipo \classname{IUserData} che conterrà ``l'utente'' che dovrà essere registrato nella rubrica del chiamante. L'oggetto dovrà essere istanziato in seguito ad una chiamata:\\
		\verb|userDAO.getByID(contactId)|;
	\end{itemize}
	Quindi se il contatto (che dovrò andare a registrare) è stato correttamente istanziato (contatto != null) allora il flusso principale procede creando e impostando un istanza di \classname{AddressBookEntry}. Il metodo termina ``scrivendo'' \inglese{true} all'interno di un istanza di \texttt{PrintWriter} creata a partire da \texttt{response.getWriter()} ed eseguendo il metodo \method{addAddressBookEntry()} (passando l'entry creata in precedenza) a partire dall'istanza che rappresenta l'utente richiedente.
	
	Se invece si è osservato che l'oggetto contenente il contatto da aggiungere, ha valore uguale a null, allora il metodo termina scrivendo \inglese{false} all'intero della dello stesso \texttt{PrintWriter} già citato.
\end{description}

%TODO da verificare

\classsection{DeleteContactsController}

\subsubsection*{Funzione}
\inglese{Controller} richiamata dal \textit{front controller} per eseguire l'eliminazione di un contatto presente nella propria rubrica.

\subsubsection*{Relazioni d'uso}

\begin{itemize}
		\item \texttt{java.io.IOException}: eccezione richiamabile dal metodo \method{doAction()}.
	\item \texttt{java.io.PrintWriter}: classe istanziata all'interno del metodo \method{doAction()}. Usata per scrivere l'output del \inglese{controller}.
	\item \texttt{javax.servlet.ServletException}: eccezione sollevabile dai metodi \method{doAction()}.
	\item \texttt{javax.servlet.http.HttpServletRequest}: classe usata per interagire con le richieste \texttt{ajax} inoltrate dal client.
	\item \texttt{javax.servlet.http.HttpServletResponse}: classe usata per interagire con le richieste \texttt{ajax} inoltrate dal client. 
	\item \classname{org.softwaresynthesis.mytalk.server.dao.DataPersistanceManager}: classe usata per comunicare tramite Hibernate con la tabella UserData della base di dati.
\end{itemize}

\subsubsection*{Classi estese ed interfacce implementate}
\begin{itemize}
	\item \texttt{server.AbstractController}: classe estesa.
\end{itemize}

\subsubsection*{Attributi}

Nessun attributo evidenziato

\subsubsection*{Metodi}
\begin{description}
	\item{\method{\# doAction(request HttpServletRequest, response HttpServletResponse): void}}\\	
	Metodo usato per eseguire la richiesta di eliminazione di un contatto presente all'interno della rubrica. Tale metodo deve creare un istanza di \texttt{AddressBookEntry} con i dati relativi alla \textit{entry} da eliminare, quindi richiamerà il metodo \method{removeAddressBookEntry()} a partire dalle due istanze di \texttt{IUserData} (quella che rappresenta l'utente richiedete e quella che rappresenta l'utente da eliminare) passando come parametro l'\textit{entry} definita in precedenza. Più nello specifico Il flusso principale inizia  con la creazione di un oggetto \classname{dao.DataPersistanceManager} avente nome UserDAO, quindi viene aperto un blocco \textit{try-catch} in cui si salva in un istanza di tipo \texttt{HttpSession}, l'oggetto ritornato da una chiamata \verb|request.getSession(false)|. I passi successivi sono (nell'ordine):
	\begin{itemize}
		\item salvare in attributo di tipo long, contactID, l'id dell'utente. Tale id dovrà essere ottenuto con una chiamata:\\
		\verb|request.getParameter("contactId")|;\\
		
		\item si crea un istanza di tipo \classname{IUserData} a partire da una chiamata a metodo:\\
		\verb|session.getAttribute("user")|;
		
		\item si crea un istanza di tipo \classname{IUserData} che conterrà ``l'utente'' che dovrà essere eliminato dalla rubrica del chiamante. L'oggetto dovrà essere istanziato in seguito ad una chiamata:\\
		\verb|userDAO.getByID(contactId)|;
	\end{itemize}
	Quindi se il contatto (che dovrò andare a cancellare dalla lista) è stato correttamente istanziato (\texttt{contatto != null}) allora il flusso principale procede creando e impostando un istanza di \classname{AddressBookEntry}. Quindi a partire dall'oggetto che rappresenta l'utente richiedente, viene richiamato il metodo \method{removeAddressBookEntry()}. Il metodo termina ``scrivendo'' \inglese{true} all'interno di un istanza di \texttt{PrintWriter} creata a partire da \texttt{response.getWriter()}.
	
	Se invece si è osservato che l'oggetto contenente il contatto da rimuovere, ha valore uguale a \textit{null}, allora il metodo termina impostando \inglese{false} all'intero della dello stesso \texttt{PrintWriter} già citato.
\end{description}

%TODO da verificare

\classsection{AddGroupController}

\subsubsection*{Funzione}
\inglese{Controller} richiamata dal \textit{front controller} per creare un nuovo gruppo nella propria rubrica.

\subsubsection*{Relazioni d'uso}

\begin{itemize}
	\item \texttt{java.io.IOException}: eccezione richiamabile dal metodo \method{doAction()}.
	\item \texttt{java.io.PrintWriter}: classe istanziata all'interno del metodo \method{doAction()}. Usata per scrivere l'output del \inglese{controller}.
	\item \texttt{javax.servlet.ServletException}: eccezione sollevabile dai metodi \method{doAction()}.
	\item \texttt{javax.servlet.http.HttpServletRequest}: classe usata per interagire con le richieste \texttt{ajax} inoltrate dal client.
	\item \texttt{javax.servlet.http.HttpServletResponse}: classe usata per interagire con le richieste \texttt{ajax} inoltrate dal client.
	\item \classname{org.softwaresynthesis.mytalk.server.dao.DataPersistanceManager}: classe usata per comunicare tramite Hibernate con la tabella UserData della base di dati.
\end{itemize}

\subsubsection*{Classi estese ed interfacce implementate}
\begin{itemize}
	\item \texttt{server.AbstractController}: classe estesa.
\end{itemize}

\subsubsection*{Attributi}

Nessun attributo evidenziato

\subsubsection*{Metodi}
\begin{description}
	\item{\method{\# doAction(request HttpServletRequest, response HttpServletResponse): void}}\\	
	Metodo usato per eseguire la richiesta di creazione di un gruppo. Nell'ordine proposto, devono essere eseguite le seguenti operazioni:
	\begin{itemize}
		\item creazione di un oggetto \classname{dao.DataPersistanceManager} avente nome groupDAO;
		\item creazione di un oggetto \texttt{HttpSession} avente nome session;
		\item creazione di un oggetto \classname{abook.IGroup} avente nome group;
		\item creazione di un oggetto \classname{abook.IUserData} avente nome user;
		\item creazione di un oggetto \texttt{PrinterWriter} avente nome writer;
		\item creazione di due stringhe, name e result;
	\end{itemize}
	A questo punto viene creato un blocco \textit{try-catch}. Il blocco \textit{catch} imposta la stringa \textit{result} a \inglese{false}, il blocco \textit{try} dovrà invece seguire il seguente iter:
	\begin{itemize}
		\item si imposta \textit{session} con il valore ritornato da una chiamata:\\
		\verb|request.getSession(false)|;
		\item si imposta la variabile user con il valore presente nell'attributo \textit{user} di sessione (usare l'istruzione \verb|session.getAttribute("user")|;
		\item impostare name con il contenuto del parametro \textit{groupName} presente nell'oggetto request;
		\item verificare se \inglese{name} è \textit{null} o vuoto. Nel caso procedere impostando \textit{result} a \textit{false} e uscendo dal blocco \textit{try}.
		\item altrimenti, se \inglese{name} contiene un nome valido si crea un nuovo gruppo (salvandolo nella variabile group), e impostando il nome e l'utente proprietario.
		\item richiamare \method{insert()} a partire dall'istanza groupDAO, passando come parametro group (che ora conterrà in gruppo da inserire nel database). Viene impostato \textit{result} a \inglese{true},
	il metodo termina scrivendo su \inglese{write} (usando l'omonimo metodo) il valore contenuto in \textit{result}. 		
	\end{itemize}
	
\end{description}

%TODO da verificare

\classsection{DeleteGroupController}

\subsubsection*{Funzione}
\inglese{Controller} richiamata dal \textit{front controller} per eliminare un gruppo da una rubrica.

\subsubsection*{Relazioni d'uso}

\begin{itemize}
\item \texttt{java.io.IOException}: eccezione richiamabile dal metodo \method{doAction()}.
	\item \texttt{java.io.PrintWriter}: classe istanziata all'interno del metodo \method{doAction()}. Usata per scrivere l'output del \inglese{controller}.
	\item \texttt{javax.servlet.ServletException}: eccezione sollevabile dai metodi \method{doAction()}.
	\item \texttt{javax.servlet.http.HttpServletRequest}: classe usata per interagire con le richieste \texttt{ajax} inoltrate dal client.
	\item \texttt{javax.servlet.http.HttpServletResponse}: classe usata per interagire con le richieste \texttt{ajax} inoltrate dal client.
	\item \classname{org.softwaresynthesis.mytalk.server.dao.DataPersistanceManager}: classe usata per comunicare tramite Hibernate con la tabella UserData della base di dati.
\end{itemize}

\subsubsection*{Classi estese ed interfacce implementate}
\begin{itemize}
	\item \texttt{server.AbstractController}: classe estesa.
\end{itemize}

\subsubsection*{Attributi}

Nessun attributo evidenziato

\subsubsection*{Metodi}
\begin{description}
	\item{\method{\# doAction(request HttpServletRequest, response HttpServletResponse): void}}\\	
	Metodo usato per eseguire la richiesta di eliminazione di un gruppo. Nell'ordine proposto, devono essere eseguite le seguenti operazioni:
	\begin{itemize}
		\item creazione di un'istanza di \classname{dao.AddressBookEntryDAO} denominata entryDAO;
		\item creazione di un \classname{dao.DataPersistanceManager} denominata groupDAO.
		\item creazione di un'istanza di \classname{abook.IGroup} denominata group;
		\item creazione di un istanza di \classname{abook.IAddressBookEntry} denominata entry;
		\item creazione di un istanza di \texttt{Iterator<IAddressBookEntry>} denominata iterator;
		\item creazione di un istanza di \texttt{PrintWriter} denominata writer;
		\item creazione di un \texttt{Set<IAddressBookEntry>} denominato entrys;
		\item creazione di una stringa denominata result (utilizzata per registrare il messaggio da stampare sul writer);
	\end{itemize} 
	A questo punto deve essere definito un blocco \textit{try-catch}. Dentro il blocco \textit{catch} si  imposta \textit{result} con il valore \inglese{false}. Per quanto riguarda il blocco \textit{try}, il programmatore dovrà definire il seguente iter:
	\begin{itemize}
		\item salvare in group il valore ritornato da una chiamata \method{getByID} a cui passo l'identificativo del gruppo (\texttt{request.getParameter("groupId")}) a partire dall'oggetto groupDAO.
		\item quindi si esegue una verifica sul contenuto di group. Se \texttt{group == null} allora si imposta \textit{result} a \inglese{false}. Altrimenti si carica in \textit{entry} l'oggetto ritornato da una chiamata \method{getAddressBook()} a partire da group.
		\item a tal punto (sempre dentro al costrutto condizionale \texttt{if(group != null))} si crea un iteratore tramite \texttt{entrys.iterator()} e ciclando con tale iteratore si va a modificare la voce ``gruppo'' di tutti i contati che nella mia rubrica appartengono a tale gruppo (cosi facendo l'eliminazione del gruppo non porterà all'eliminazione dei contatti presenti in tale gruppo).
		\item dunque si esegue la chiamata:
		\verb|groupDAO.delete(group)|\\
		e si imposta \textit{result} a \inglese{true}.
	\end{itemize}
	Il metodo termina scrivendo su \textit{writer} il contenuto di \textit{result}.
	
\end{description}

%TODO da verificare

\classsection{AddInGroupController}

\subsubsection*{Funzione}
\inglese{Controller} richiamata dal \textit{front controller} per inserire un utente nella propria rubrica, all'interno di un gruppo ben determinato.

\subsubsection*{Relazioni d'uso}

\begin{itemize}
	\item \texttt{java.io.IOException}: eccezione richiamabile dal metodo \method{doAction()}.
	\item \texttt{java.io.PrintWriter}: classe istanziata all'interno del metodo \method{doAction()}. Usata per scrivere l'output del \inglese{controller}.
	\item \texttt{javax.servlet.ServletException}: eccezione sollevabile dai metodi \method{doAction()}.
	\item \texttt{javax.servlet.http.HttpServletRequest}: classe usata per interagire con le richieste \texttt{ajax} inoltrate dal client.
	\item \texttt{javax.servlet.http.HttpServletResponse}: classe usata per interagire con le richieste \texttt{ajax} inoltrate dal client.
	\item \classname{abook.IAddressBookEntry}:Interfaccia che definisce il comportamento di una generica \inglese{entry} della rubrica utente. La classe che verrà descritta crea istanze di tipo \classname{abook.IAddressBookEntry}.
	\item \classname{abook.IGroup}:
Interfaccia che definisce il comportamento di un gruppo generico. La classe che qui si descrive crea istanze di tipo \classname{abook.IGroup}.
	\item \classname{abook.IUserData}:Interfaccia che definisce il comportamento di un utente generico. La classe che qui si descrive crea istanze di tipo \classname{abook.IUserData}.
	\item \classname{abook.AddressBookEntry}: classe usata per comunicare tramite \textit{Hibernate} con la tabella AddressBookEntry della base di dati.
	\item \classname{dao.DataPersistanceManager}: classe usata per comunicare tramite \textit{Hibernate} con la tabella Group della base di dati.
	\item \classname{dao.DataPersistanceManager}: classe usata per comunicare tramite \textit{Hibernate} con la tabella UserData della base di dati.
\end{itemize}

\subsubsection*{Classi estese ed interfacce implementate}
\begin{itemize}
	\item \texttt{server.AbstractController}: classe estesa.
\end{itemize}

\subsubsection*{Attributi}

Nessun attributo evidenziato

\subsubsection*{Metodi}

\begin{description}
	\item{\method{\# doAction(request HttpServletRequest, response HttpServletResponse): void}}\\	
	Metodo usato per eseguire la richiesta di aggiungere un contatto in un gruppo. Nell'ordine proposto, devono essere eseguite le seguenti operazioni:
	\begin{itemize}
		\item creazione di un istanza di \classname{DataPersistanceManager} denominata groupDAO;
		\item creazione di una sessione \texttt{HttpSession} avente nome session;
		\item creazione degli oggetti rappresentanti la realtà della rubrica su cui opero: un \texttt{IAddressBookEntry} entry, \texttt{IGroup} group, \texttt{IUserData} user e \texttt{IUserData} friend;
		\item creazione di due identificativi di tipo long: contactId e groupId.
		\item creazione di un \texttt{PrintWriter} denominato \textit{writer} e di una stringa \textit{result} usata con lo scopo di memorizzare il contenuto testuale da scrivere sul writer come messaggio di notifica della \inglese{servlet};
		\item creazione di un istanza \texttt{DataPersistanceManager} userDAO.
	\end{itemize}
	Dopo questa fase di creazione delle variabili il metodo deve definire un costrutto \textit{try-catch}. All'interno del blocco catch dovrà essere predisposta la memorizzazione della parola ``\inglese{false}''. Passando invece alla definizione del blocco \textit{try}, al suo interno si dovranno predisporre le seguenti istruzioni:
	\begin{itemize}
		\item inizializzazione di session al valore ottenuto da una chiamata:\\
		\verb|request.getSession(false)|
		\item istanziazione di user a partire dal valore contenuto nella \textit{request}. Per la precisione si dovrà usare un istruzione del tipo:\\
		\verb|session.getAttribute("user")|
		\item con una procedura analoga alla precedente definisco il contenuto di contactID. Usare \verb|getParameter("contactId")| a partire dall'oggetto \textit{request};
		\item istruzione analoga alla precedente usata per istanziare il contenuto di groupId (nome del parametro da ottenere: groupId);
		\item inizializzare friend con una chiamata \method{getById} passando come parametro contactId.
		\item inizializzare group con una chiamata \method{getById} passando come parametro groupId.
		\item il metodo controlla se \texttt{group!=null}. Nel caso il flusso principale prosegue come segue: inizializzazione di \textit{entry} e modifica dei dati stessi di \textit{entry} mediante le chiamate a metodo dei vari ``set'' che la costituiscono. Nello specifico si intende impostare l'istanza in modo che definisca un contatto in rubrica non bloccato e registrato nel gruppo \textit{group}. Il possessore sarà user e il contatto registrato friend. Quindi viene eseguito l'\inglese{update} tramite una chiamata:\\
		\verb|userDAO.update(user)|
		result viene impostato a \textit{true} e il programma esce dal costrutto condizionale.
		\item se \textit{group} non era diverso da \textit{null} allora si entra nel ramo else del costrutto condizionale già citato. Result viene impostato a false e il metodo esce dal ramo else.
	\end{itemize}
	Il metodo termina scrivendo su \textit{writer} il contenuto di \textit{result}.
	
\end{description}

%TODO da verificare

\classsection{DeleteFromGroupController}

\subsubsection*{Funzione}
\inglese{Controller} richiamata dal \textit{front controller} per inserire un utente nella propria rubrica, all'interno di un gruppo ben determinato.

\subsubsection*{Relazioni d'uso}

\begin{itemize}
	\item \texttt{java.io.IOException}: eccezione richiamabile dal metodo \method{doAction()}.
	\item \texttt{java.io.PrintWriter}: classe istanziata all'interno del metodo \method{doAction()}. Usata per scrivere l'output del \inglese{controller}.
	\item \texttt{javax.servlet.ServletException}: eccezione sollevabile dai metodi \method{doAction()}.
	\item \texttt{javax.servlet.http.HttpServletRequest}: classe usata per interagire con le richieste \texttt{ajax} inoltrate dal client.
	\item \texttt{javax.servlet.http.HttpServletResponse}: classe usata per interagire con le richieste \texttt{ajax} inoltrate dal client. \classname{abook.IAddressBookEntry}.
	\item \classname{abook.IGroup}:
Interfaccia che definisce il comportamento di un gruppo generico. La classe che verrà descritta crea istanze di tipo \classname{abook.IGroup}.
	\item \classname{abook.IUserData}:Interfaccia che definisce il comportamento di un utente generico. La classe che verrà descritta crea istanze di tipo \classname{abook.IUserData}.
	\item \classname{abook.AddressBookEntry}: classe usata per comunicare tramite \textit{Hibernate} con la tabella AddressBookEntry della base di dati.
	\item \classname{dao.DataPersistanceManager}: classe usata per comunicare tramite \textit{Hibernate} con la tabella Group e UserData della base di dati.
\end{itemize}

\subsubsection*{Classi estese ed interfacce implementate}
\begin{itemize}
	\item \texttt{server.AbstractController}: classe estesa.
\end{itemize}

\subsubsection*{Attributi}

Nessun attributo evidenziato


\subsubsection*{Metodi}

\begin{description}
	\item{\method{\# doAction(request HttpServletRequest, response HttpServletResponse): void}}\\	
	Metodo usato per eseguire la richiesta di rimuovere un contatto da un gruppo. Nell'ordine proposto, devono essere eseguite le seguenti operazioni:
	\begin{itemize}
		\item creazione degli oggetti necessari al completamento dell'operazione:
		\begin{itemize}
			\item \classname{DataPersistanceManager} groupDAO;
			\item \texttt{HttpSession} session;
			\item \classname{IAddressBookEntry} entry;
			\item \classname{IGroup} group;
			\item \classname{IUserData} friend;
			\item \classname{IUserData} user;
			Long contactIdl;
			Long groupId;
			\item \classname{PrintWriter} writer;
			String result;
			\item \classname{DataPersistanceManager} userDAO;
		\end{itemize}
		
		\item quindi il metodo procede definendo un blocco \textit{try-catch}. Nel ramo \textit{catch} si imposta \textit{result} a \inglese{false}. Passando invece alla definizione del ramo \textit{try}, in esso devono essere definiti i seguenti punti:
		\begin{itemize}
			\item session viene impostato mediante una chiamata a metodo:
			\verb|request.getSession(false)|;
			\item user e groupID vengono impostate ottenendo l'omonimo attributo da \textit{request} (usare il metodo \texttt{getParameter(nome\_parametro)};
			\item impostare friend e group a partire dai relativi tipi di istanze DAO e facendosi restituire i dati presenti nel database tramite una chiamata a metodo \method{getByID} a cui passa i relativi id (contactId e groupid);
			\item avvia un costrutto condizionale dotato di ramo \textit{else}, se \texttt{group!=null} allora il metodo procede nel ramo \textit{if} andando ad impostare \textit{entry} con i relativi parametri, in modo da ricreare l'istanza \classname{IAddressBookEntry} da rimuovere dal database per mezzo del metodo \method{removeAddressBookEntry()} a cui passa \textit{entry}. 
			\item altrimenti, se \texttt{group == null} allora entra nel ramo else del costrutto condizionale e procede impostando \textit{result} a \inglese{false}.
		\end{itemize}
	\end{itemize}
	Il metodo termina scrivendo su \textit{writer} il contenuto di \textit{result}.
	
\end{description}

%TODO da verificare

\classsection{BlockContactController}

\subsubsection*{Funzione}
\inglese{Controller} richiamata dal \textit{front controller} per bloccare un contatto presente nella propria rubrica.

\subsubsection*{Relazioni d'uso}

\begin{itemize}
	\item \texttt{java.io.IOException}: eccezione richiamabile dal metodo \method{doAction()}.
	\item \texttt{java.io.PrintWriter}: classe istanziata all'interno del metodo \method{doAction()}. Usata per scrivere l'output del \inglese{controller}.
	\item \texttt{javax.servlet.ServletException}: eccezione sollevabile dai metodi \method{doAction()}.
	\item \texttt{javax.servlet.http.HttpServletRequest}: classe usata per interagire con le richieste \texttt{ajax} inoltrate dal client.
	\item \texttt{javax.servlet.http.HttpServletResponse}: classe usata per interagire con le richieste \texttt{ajax} inoltrate dal client.
	\item \classname{abook.IAddressBookEntry}:Interfaccia che definisce il comportamento di una generica \inglese{entry} della rubrica utente. La classe che verrà descritta crea istanze di tipo \classname{abook.IAddressBookEntry}.
	\item \classname{abook.IUserData}:Interfaccia che definisce il comportamento di un utente generico. La classe che verrà descritta crea istanze di tipo \classname{abook.IUserData}.
	\item \classname{dao.DataPersistanceManager}: classe usata per comunicare tramite \textit{Hibernate} con la tabella UserData della base di dati.
\end{itemize}

\subsubsection*{Classi estese ed interfacce implementate}
\begin{itemize}
	\item \texttt{server.AbstractController}: classe estesa.
\end{itemize}

\subsubsection*{Attributi}

Nessun attributo evidenziato


\subsubsection*{Metodi}

\begin{description}
	\item{\method{\# doAction(request HttpServletRequest, response HttpServletResponse): void}}\\	
	Metodo usato per eseguire la richiesta di bloccare un contatto presente in una rubrica utente. Nell'ordine proposto, devono essere eseguite le seguenti operazioni:
	\begin{itemize}
		\item creazione degli oggetti necessari al completamento dell'operazione:
		\begin{itemize}
			\item \classname{DataPersistanceManager} groupDAO;
			\item \texttt{HttpSession} session;
			\item \classname{IAddressBookEntry} entry;
			\item \classname{IUserData} friend;
			\item \classname{IUserData} user;
			Long contactIdl;
			\item \texttt{Iterator<IAddressBookEntry>} iterator;
			\item \texttt{Set<IAddressBookEntry>} entrys;
			\item \classname{PrintWriter} writer;
			String result;
			\item \classname{DataPersistanceManager} userDAO;
		\end{itemize}
		
		\item quindi il metodo procede definendo un blocco \textit{try-catch}. Nel ramo \textit{catch} si imposta \textit{result} a \inglese{false}. Passando invece alla definizione del ramo \textit{try} devono essere definiti i seguenti punti:
		\begin{itemize}
			\item \textit{session} viene impostato mediante una chiamata a metodo:
			\verb|request.getSession(false)|;
			\item \textit{user} viene impostato ottenendo l'omonimo attributo da \textit{request} (usare il metodo \texttt{getParameter(``user'')};
			\item impostare \textit{friend} a partire dal relativo tipo di istanza DAO e facendo restituire i dati presenti nel database tramite una chiamata a metodo \method{getByID} a cui passa il parametro \textit{contactId};
			\item avvia un costrutto condizionale dotato di ramo else. Se \texttt{friend!=null} il metodo procede nel ramo \textit{if} andando a scaricare l'elenco delle \textit{entry} a partire da \textit{user}. Quindi utilizzando l'iteratore \textit{iterator} va a scorrere tale set di \textit{entry} e quando verifica la presenza del contatto \textit{friend} all'interno, modifica l'\textit{entry} attualmente selezionata richiamando il metodo \method{setBlocked(true)}. Quindi prima di uscire dal costrutto \textit{if}, viene eseguita un operazione di \textit{update} a partire da user e si memorizza nella variabile \textit{result} il valore \inglese{true}.
			\item altrimenti, se \texttt{friend == null} allora entra nel ramo else del costrutto condizionale e procede impostando \textit{result} a \inglese{false}.
		\end{itemize}
	\end{itemize}
	Il metodo termina scrivendo su \textit{writer} il contenuto di \textit{result}.
	
\end{description}

%TODO da verificare

\classsection{UnblockContactController}

\subsubsection*{Funzione}
\inglese{Controller} richiamata dal \textit{front controller} per sbloccare un contatto presente nella propria rubrica.

\subsubsection*{Relazioni d'uso}

\begin{itemize}
	\item \texttt{java.io.IOException}: eccezione richiamabile dal metodo \method{doAction()}.
	\item \texttt{java.io.PrintWriter}: classe istanziata all'interno del metodo \method{doAction()}. Usata per scrivere l'output del \inglese{controller}.
	\item \texttt{javax.servlet.ServletException}: eccezione sollevabile dai metodi \method{doAction()}.
	\item \texttt{javax.servlet.http.HttpServletRequest}: classe usata per interagire con le richieste \texttt{ajax} inoltrate dal client.
	\item \texttt{javax.servlet.http.HttpServletResponse}: classe usata per interagire con le richieste \texttt{ajax} inoltrate dal client.
	\item \classname{abook.IAddressBookEntry}:Interfaccia che definisce il comportamento di una generica \inglese{entry} della rubrica utente. La classe che verrà descritta crea istanze di tipo \classname{abook.IAddressBookEntry}.
	\item \classname{abook.IUserData}:Interfaccia che definisce il comportamento di un utente generico. La classe che verrà descritta crea istanze di tipo \classname{abook.IUserData}.
	\item \classname{dao.DataPersistanceManager}: classe usata per comunicare tramite \textit{Hibernate} con la tabella UserData della base di dati.
\end{itemize}

\subsubsection*{Classi estese ed interfacce implementate}
\begin{itemize}
	\item \texttt{server.AbstractController}: classe estesa.
\end{itemize}

\subsubsection*{Attributi}

Nessun attributo evidenziato

\subsubsection*{Metodi}

\begin{description}
	\item{\method{\# doAction(request HttpServletRequest, response HttpServletResponse): void}}\\	
	Metodo usato per eseguire la richiesta di sbloccare un contatto presente in una rubrica utente. Nell'ordine proposto, devono essere eseguite le seguenti operazioni:
	\begin{itemize}
		\item creazione degli oggetti necessari al completamento dell'operazione:
		\begin{itemize}
			\item \classname{DataPersistanceManager} groupDAO;
			\item \texttt{HttpSession} session;
			\item \classname{IAddressBookEntry} entry;
			\item \classname{IUserData} friend;
			\item \classname{IUserData} user;
			Long contactIdl;
			\item \texttt{Iterator<IAddressBookEntry>} iterator;
			\item \texttt{Set<IAddressBookEntry>} entrys;
			\item \classname{PrintWriter} writer;
			String result;
			\item \classname{DataPersistanceManager} userDAO;
		\end{itemize}
		
		\item quindi il metodo procede definendo un blocco \textit{try-catch}. Nel ramo \textit{catch} si imposta \textit{result} a \inglese{false}. Passando invece alla definizione del ramo \textit{try}, in esso devono essere definiti i seguenti punti:
		\begin{itemize}
			\item \textit{session} viene impostato mediante una chiamata a metodo:
			\verb|request.getSession(false)|;
			\item user viene impostata ottenendo l'omonimo attributo da \textit{request} (usare il metodo \texttt{getParameter(``user'')};
			\item impostare\textit{friend} a partire dal relativo tipo di istanza DAO e facendo restituire i dati presenti nel database tramite una chiamata a metodo \method{getByID} a cui passa il parametro \textit{contactId};
			\item avvia un costrutto condizionale dotato di ramo else. Se \texttt{friend!=null} il metodo procede nel ramo \textit{if} andando a scaricare l'elenco delle \textit{entry} a partire da \textit{user}. Quindi utilizzando l'iteratore \textit{iterator} va a scorrere tale set di \textit{entry} e quando verifica la presenza del contatto \textit{friend\textit} all'interno, modifica l'\textit{entry} attualmente selezionata richiamando il metodo \method{setBlocked(false)}. Quindi prima di uscire dal costrutto \textit{if}, viene eseguita un operazione di \textit{update} a partire da \textit{user} e si memorizza nella variabile \textit{result} il valore \inglese{true}.
			\item altrimenti, se \texttt{friend == null} allora entra nel ramo else del costrutto condizionale e procede impostando \textit{result} a \inglese{false}.
		\end{itemize}
	\end{itemize}
	Il metodo termina scrivendo su \textit{writer} il contenuto di \textit{result}.
	
\end{description}

%TODO da verificare

\classsection{SearchController}

\subsubsection*{Funzione}
\inglese{Controller} richiamata dal \textit{front controller} per effettuare una ricerca sulla propria rubrica. La ricerca consiste nell'individuare tutti i contatti che contengono nei campi \textit{name}, \textit{surname} o \textit{email} la parola ricercata. Per eseguire questa procedura di ricerca ci si avvale del metodo \method{searchGeneric(nome\_parametro)} della classe \classname{DataPersistanceManager}.

\subsubsection*{Relazioni d'uso}

\begin{itemize}
	\item \texttt{java.io.IOException}: eccezione richiamabile dal metodo \method{doAction()}.
	\item \texttt{java.io.PrintWriter}: classe istanziata all'interno del metodo \method{doAction()}. Usata per scrivere l'output del \inglese{controller}.
	\item \texttt{javax.servlet.ServletException}: eccezione sollevabile dai metodi \method{doAction()}.
	\item \texttt{javax.servlet.http.HttpServletRequest}: classe usata per interagire con le richieste \texttt{ajax} inoltrate dal client.
	\item \texttt{javax.servlet.http.HttpServletResponse}: classe usata per interagire con le richieste \texttt{ajax} inoltrate dal client.
	\item \classname{abook.IUserData}:Interfaccia che definisce il comportamento di un utente generico. La classe che verrà descritta crea istanze di tipo \classname{abook.IUserData}.
	\item \classname{dao.DataPersistanceManager}: classe usata per comunicare tramite \textit{Hibernate} con la tabella UserData della base di dati.
\end{itemize}

\subsubsection*{Classi estese ed interfacce implementate}
\begin{itemize}
	\item \texttt{server.AbstractController}: classe estesa.
\end{itemize}

\subsubsection*{Attributi}

Nessun attributo evidenziato

\subsubsection*{Metodi}

\begin{description}
	\item{\method{\# doAction(request HttpServletRequest, response HttpServletResponse): void}}\\	
	Metodo usato per eseguire una ricerca generica sui contatti presenti nella rubrica del \textit{client}. Nell'ordine proposto, devono essere eseguite le seguenti operazioni:
	\begin{itemize}
		\item creazione degli oggetti necessari al completamento dell'operazione:
		\begin{itemize}
			\item \classname{IUserData} entry;
			\item \texttt{List<IUserData>} users;
			\item \texttt{Iterator<IUserData>} iterator;
			\item \texttt{Set<IAddressBookEntry>} entrys;
			\item \classname{PrintWriter} writer;
			\item String result e String parameter;
			\item \classname{DataPersistanceManager} userDAO;
		\end{itemize}
		
		\item quindi il metodo procede definendo un blocco \textit{try-catch}. Nel ramo \textit{catch} si imposta \textit{result} a \inglese{false}. Passando invece alla definizione del ramo \textit{try}, in esso devono essere definiti i seguenti punti:
		\begin{itemize}
			\item parameter viene impostato con il valore del parametro da ricercare. Tale operazione è da eseguirsi con l'istruzione:\\
			\verb|request.getParameter("param")|;
			\item quindi il metodo procede effettuando la ricerca del parametro. Tale operazione si svolge sfruttando le specifiche della classe \classname{DataPersistanceManager}. Nello specifico deve essere eseguita l'istruzione:\\
			\verb|users = userDAO.searchGeneric(parameter)|;
			\item a questo punto si dovrà predisporre un iteratore per scorrere la lista \textit{users};
			\item si imposta \textit{result} al valore ``{'' e si entra in un ciclo \inglese{while} che potrà terminare solo quando l'iteratore già citato avrà raggiunto l'ultimo elemento ispezionabile;
			\item dentro al ciclo \inglese{while} si procederà con la formattazione della stringa \textit{result} al fine di restituire al \textit{client} un formato sul quale possa operare. Nello specifico la stringa \textit{result} dovrà contenere (per ogni istanza di IUserData presente in \textit{users}):\\
			
			\verb|\"ID_della_entry":{"name":"NOME_UTENTE",|\\
			\verb|"surname":"COGNOME_UTENTE,"email":EMAIL_UTENTE",|\\
			\verb|"id":"ID_UTENTE","picturePath":"PATH_IMG",|\\
			\verb|"state":"STATO","block":"BLOCCATO/SBLOCCATO"},|\\
			
			Si osservi che in quanto è riportato in precedenza, ciò che è scritto in maiuscolo corrisponde al valore effettivo di quel parametro, quindi per esempio la parola NOME\_UTENTE sarà di fatto sostituita dall'effettivo nome dell'istanza \classname{IUserData} attualmente sotto esame. Per ottenere tali valori ricorrerà ai metodi: \texttt{getName()}, \texttt{getSurname()}, \texttt{getMail()}, \texttt{getId()}, \texttt{getPath()} richiamabili a partire dall'\textit{entry} attualmente sotto esame.
			\item al termine di tale procedura la stringa \textit{result} è pronta per essere restituita al \textit{client}, quindi si aggiunge a \textit{result} il valore ``}'', usato come carattere di terminazione.
		\end{itemize}
	\end{itemize}
	Il metodo termina scrivendo su \textit{writer} il contenuto di \textit{result}.
	
\end{description}

%TODO da verificare

\classsection{GetContactsController}

\subsubsection*{Funzione}
\inglese{Controller} richiamata dal \textit{front controller} per scaricare la lista dei contatti presenti nella propria rubrica.

\subsubsection*{Relazioni d'uso}

\begin{itemize}
	\item \texttt{java.io.IOException}: eccezione richiamabile dal metodo \method{doAction()}.
	\item \texttt{java.io.PrintWriter}: classe istanziata all'interno del metodo \method{doAction()}. Usata per scrivere l'output del \inglese{controller}.
	\item \texttt{javax.servlet.ServletException}: eccezione sollevabile dai metodi \method{doAction()}.
	\item \texttt{javax.servlet.http.HttpServletRequest}: classe usata per interagire con le richieste \texttt{ajax} inoltrate dal client.
	\item \texttt{javax.servlet.http.HttpServletResponse}: classe usata per interagire con le richieste \texttt{ajax} inoltrate dal client.
	\item \classname{abook.IAddressBookEntry}:Interfaccia che definisce il comportamento di una \inglese{entry} della rubrica utente.
	\item \classname{abook.IUserData}:Interfaccia che definisce il comportamento di un utente generico. La classe che verrà descritta crea istanze di tipo \classname{abook.IUserData}.
	\item \classname{dao.DataPersistanceManager}: classe usata per comunicare tramite Hibernate con la tabella UserData della base di dati.
\end{itemize}

\subsubsection*{Classi estese ed interfacce implementate}
\begin{itemize}
	\item \texttt{server.AbstractController}: classe estesa.
\end{itemize}

\subsubsection*{Attributi}

Nessun attributo evidenziato

\subsubsection*{Metodi}

\begin{description}
	\item{\method{\# doAction(request HttpServletRequest, response HttpServletResponse): void}}\\	
	Metodo usato per eseguire la richiesta di download della lista dei contatti. Nell'ordine proposto, devono essere eseguite le seguenti operazioni:
	\begin{itemize}
		\item creazione degli oggetti necessari al completamento dell'operazione:
		\begin{itemize}
			\item \texttt{HttpSession} session;
			\item \classname{IUserData} user;
			\item \classname{IUserData} friend;
			\item \texttt{Iterator<IAddressBookEntry>} iterator;
			\item \texttt{Set<IAddressBookEntry>} entrys;
			\item \classname{PrintWriter} writer;
			\item String result;
			\item \classname{IAddressBookEntry} entry;
		\end{itemize}
		
		\item quindi il metodo procede definendo un blocco \textit{try-catch}. Nel ramo \textit{catch} si imposta \textit{result} a \inglese{null}. Passando invece alla definizione del ramo \textit{try}, devono essere definiti i seguenti punti:
		\begin{itemize}
			\item session viene impostata tramite una chiamata a metodo del tipo:\\
			\verb|session = request.getSession(false)|;
			\item quindi il metodo procede impostando user a partire dalla sessione precedentemente creata. Tale operazione richiede l'uso del metodo \method{getAttribute(``user''};
			\item viene inizializzato l'insieme \textit{contacts} tramite una chiamata \method{getAddressBook()} richiamata a partire dalla variabile \textit{user};
			\item a questo punto si dovrà predisporre un iteratore per scorrere l'insieme \textit{contacts};
			\item si imposta \textit{result} al valore ``{'' e si entra in un ciclo \inglese{while} che potrà terminare solo quando l'iteratore già citato avrà raggiunto l'ultimo elemento ispezionabile;
			\item dentro al ciclo \inglese{while}, dopo aver impostato \textit{entry} al valore \texttt{next()} dell'iteratore, viene salvato in friend il contatto ritornato da una chiamata a \textit{entry}.\method{getContact()}. Si procederà  poi con la formattazione della stringa result al fine di restituire al \textit{client} un formato sul quale possa operare. Nello specifico la stringa result dovrà contenere (per ogni istanza di IUserData presente in users):\\
			
			\verb|\"ID_della_entry":{"name":"NOME_UTENTE",|\\
			\verb|"surname":"COGNOME_UTENTE,"email":EMAIL_UTENTE",|\\
			\verb|"id":"ID_UTENTE","picturePath":"PATH_IMG",|\\
			\verb|"state":"STATO","block":"BLOCCATO/SBLOCCATO"},|\\
			
			Si osservi che in quanto è riportato in precedenza, ciò che è scritto in maiuscolo corrisponde al valore effettivo di quel parametro. Quindi per esempio la parola NOME\_UTENTE sarà di fatto sostituita dall'effettivo nome dell'istanza \classname{IUserData} attualmente sotto esame. Per ottenere tali valori si ricorre ai metodi: \texttt{getName()}, \texttt{getSurname()}, \texttt{getMail()}, \texttt{getId()}, \texttt{getPath()} richiamabili a partire da friend.
			\item al termine di tale procedura la stringa \textit{result} è pronta per essere restituita al \textit{client}. Quindi si aggiunge a \textit{result} il valore ``}'', usato come carattere di terminazione.
		\end{itemize}
	\end{itemize}
	Il metodo termina scrivendo su \textit{writer} il contenuto di \textit{result}.
	
\end{description}

%TODO da verificare

\classsection{GetGroupsController}

\subsubsection*{Funzione}
\inglese{Controller} richiamata dal \textit{front controller} per scaricare la propria rubrica, provvista di gruppi e contatti in essi presenti.

\subsubsection*{Relazioni d'uso}

\begin{itemize}
	\item \texttt{java.io.IOException}: eccezione richiamabile dal metodo \method{doAction()}.
	\item \texttt{java.io.PrintWriter}: classe istanziata all'interno del metodo \method{doAction()}. Usata per scrivere l'output del \inglese{controller}.
	\item \texttt{javax.servlet.ServletException}: eccezione sollevabile dai metodi \method{doAction()}.
	\item \texttt{javax.servlet.http.HttpServletRequest}: classe usata per interagire con le richieste \texttt{ajax} inoltrate dal client.
	\item \texttt{javax.servlet.http.HttpServletResponse}: classe usata per interagire con le richieste \texttt{ajax} inoltrate dal client.
	\item \classname{abook.IAddressBookEntry}:Interfaccia che definisce il comportamento di una \inglese{entry} della rubrica utente.
	\item \classname{abook.IUserData}:Interfaccia che definisce il comportamento di un utente generico. La classe che verrà descritta crea istanze di tipo \classname{abook.IUserData}.
		\item \classname{abook.IGroup}:Interfaccia che definisce il comportamento di un gruppo generico. La classe che verrà descritta crea istanze di tipo \classname{abook.IGroup}.
	\item \classname{dao.DataPersistanceManager}: classe usata per comunicare tramite \textit{Hibernate} con la tabella Group della base di dati.
\end{itemize}

\subsubsection*{Classi estese ed interfacce implementate}
\begin{itemize}
	\item \texttt{server.AbstractController}: classe estesa.
\end{itemize}

\subsubsection*{Attributi}

Nessun attributo evidenziato

\subsubsection*{Metodi}

\begin{description}
\item{\method{\# doAction(request HttpServletRequest, response HttpServletResponse): void}}\\	
	Metodo usato per eseguire la richiesta di download della lista dei contatti utente). Nell'ordine proposto, devono essere eseguite le seguenti operazioni:
	\begin{itemize}
		\item creazione degli oggetti necessari al completamento dell'operazione:
		\begin{itemize}
			\item \texttt{HttpSession} session;
			\item \classname{DataPersistanceManager} groupDAO;
			\item \classname{IAddressBookEntry} entry;
			\item \classname{IGroup} group;
			\item \classname{IUserData} user;
			\item \texttt{Iterator<IAddressBookEntry>} entryIter;
			\item \texttt{Iterator<IGroup>} groupIter;
			\item \texttt{Set<IAddressBookEntry>} addentrys;
			\item \classname{PrintWriter} writer;
			\item String result;
		\end{itemize}
		
		\item quindi il metodo procede definendo un blocco \textit{try-catch}. Nel ramo catch si imposta \textit{result} a \inglese{false}. Passando invece alla definizione del ramo \textit{try}, in esso devono essere definiti i seguenti punti:
		\begin{itemize}
			\item session viene impostata tramite una chiamata a metodo del tipo:\\
			\verb|session = request.getSession(false)|;
			\item quindi il metodo procede impostando user a partire dalla sessione precedentemente creata. Tale operazione richiede l'uso del metodo \method{getAttribute(``user'')};
			\item viene inizializzato l'oggetto groupDAO;
			\item si inizializza la lista groups dei gruppi. Tale operazione deve essere eseguita mediante l'istruzione:\\
			\verb|groupDAO.getByOwner(user.getId())|
			\item il metodo definisce un costrutto condizionale \textit{if-else} basato sulla condizione \texttt{groups != null}. Nel ramo \textit{else} il metodo non fa altro che impostare \textit{result} a \inglese{false}. Nel ramo \textit{if} saranno definite le seguenti istruzioni:	
			\begin{itemize}
				\item predisporre un iteratore per scorrere l'insieme groups (inizializzazione di groupIter);
				\item si imposta \textit{result} al valore ``\{'' e si entra in un ciclo \inglese{while} che potrà terminare solo quando l'iteratore già citato avrà raggiunto l'ultimo elemento ispezionabile;
				\item all'interno del ciclo si estrae dall'iteratore il gruppo attualmente in esame e si salva sull'oggetto group;
				\item a partire da group, tramite una chiamata \method{getAddressBook()}, si ottiene la lista di entry di quel gruppo e la si slava in addEntry;
				\item si definisce un nuovo iteratore addEntry.\method{iterator()} e lo si memorizza in entryIter;
				\item si ``scrive'' dentro a \textit{result}, i dati inerenti il gruppo in esame (name, id) e si aggiunge la stringa ``contacts:['';
				\item si entra in un nuovo ciclo \inglese{while} che cicla sulle \textit{entry} dell'iteratore entryIter. All'interno di tale ciclo si estraggono le \textit{entry} che costituiscono i contatti del gruppo in esame e si riporta in result l'id di tali contatti (usare metodo \method{getId()});
				\item all'uscita del ciclo \inglese{while} più interno si concatena al contenuto di \textit{result} la stringa ``]\}'';
				\item all'uscita del ciclo \inglese{while} più esterno si concatena al contenuto di \textit{result} la stringa ``\}'';
			\end{itemize}
		\end{itemize}
	\end{itemize}
	Il metodo termina scrivendo su \textit{writer} il contenuto di \textit{result}.	
	
\end{description}

\subsection{Package org.softwaresynthesis.mytalk.server.call}\label{sec:call}

\classsection{ICall}

\subsubsection*{Funzione}
Interfaccia che rappresenta una chiamata effettuata dal sistema \caName. Gli oggetti che implementano tale interfaccia vengono usati per rappresentare lo storico delle chiamate di un utente.

\subsubsection*{Relazioni d'uso}
\begin{itemize}
	\item \texttt{java.util.Date}: tipo utilizzato per definire la data d'inizio e fine di una chiamata.
\end{itemize}

\subsubsection*{Classi estese ed interfacce implementate}

\begin{itemize}
	\item \texttt{org.softwaresynthesis.mytalk.server.IMyTalkObject}: interfaccia estesa. Utilizzata per dire che la classe è da considerarsi come un \texttt{transfer object} gestito da \texttt{Hibernate}.
\end{itemize}

\subsubsection*{Metodi}
\begin{description}
	\item{\method{+ getId(): Long}}\\
	Restituisce l'identificativo univoco della chiamata.
	\item{\method{+ getStart(): Date}}\\
	Restituisce un istanza di \texttt{java.util.Date} di avvio della chiamata.
	\item{\method{+ setStart(startDate: Date): void}}\\
	Imposta la data di avvio della chiamata.
	\item{\method{+ getEnd(): Date}}\\
	Restituisce la data in cui termina la chiamata
	\item{\method{+ setEnd(endDate: Date): void}}\\
	Imposta la data in cui termina la chiamata.
	\item{\method{+ getCall(): Set<ICallList>}}\\
	Metodo usato per ottenere la lista dei partecipanti alla chiamata.
	\item{\method{+ setCall(callList: Set<ICallList>): void}}\\
	Metodo usato per impostare la lista dei partecipanti alla chiamata.
	\item{\method{+ addCall(call: ICallList): boolean}}\\
	Metodo usato per aggiungere un partecipante alla chiamata.
	\item{\method{+ removeCall(call: ICallList): boolean}}\\
	Metodo usato per rimuovere un partecipate alla chiamata 
\end{description}

\classsection{Call}

\subsubsection*{Funzione}
Classe che implementa l'interfaccia \classname{ICall}.

\subsubsection*{Relazioni d'uso}
\begin{itemize}
	\item \texttt{java.util.Date}: tipo utilizzato per definire la data d'inizio e fine di una chiamata.
\end{itemize}

\subsubsection*{Classi estese ed interfacce implementate}
\begin{itemize}
	\item \classname{ICall} interfacci d'implementazione.
\end{itemize}

\subsubsection*{Attributi}

\begin{itemize}
	\item{\memberdata{-- id: long}}
	Attributo contenente il codice identificativo della chiamata.
	\item{\memberdata{-- startDate: Date}}
	Attributo di tipo \texttt{java.util.Date} contenente la data (compresa l'ora) d'inizio della chiamata.
	\item{\memberdata{-- endDate: Date}}
	Attributo di tipo \texttt{java.util.Date} contenente la data (compresa l'ora) in cui la chiamata è terminata.
	\item{\memberdata{-- calls: Set<ICallList>}}
	Attributo usato per definire la lista dei partecipanti alla chiamata.
\end{itemize}

\subsubsection*{Metodi}
\begin{description}
	\item{\method{+ getId(): Long}}\\
	Restituisce l'identificativo univoco della chiamata ritornando il contenuto dell'attributo \memberdata{id}.
	\item{\method{+ getStart(): Date}}\\
	Restituisce il contenuto del campo \memberdata{startDate}.
	\item{\method{+ setStart(startDate: Date): void}}\\
	Imposta la data di avvio della chiamata, sovrascrivendo il contenuto \memberdata{startData}.
	\item{\method{+ getEnd(): Date}}\\
	Restituisce il contenuto del campo \memberdata{endDate}.
	\item{\method{+ setEnd(endDate: Date): void}}\\
	Imposta la data di avvio della chiamata, sovrascrivendo il contenuto \memberdata{endData}.
	\item{\method{+ getCall(): Set<ICallList>}}\\
	Restituisce il contenuto del campo \memberdata{calls}.
	\item{\method{+ setCall(callList: Set<ICallList>): void}}\\
	Metodo usato per impostare la lista dei partecipanti alla chiamata. Sovrascrive il contenuto di \memberdata{calls}.
	\item{\method{+ addCall(call: ICallList): boolean}}\\
	Metodo usato per aggiungere un partecipante alla chiamata. Aggiunge al \memberdata{calls} un nuovo partecipante (call). Il metodo termina ritornando \inglese{true} se l'operazione va a buon fine. \inglese{false} altrimenti.
	\item{\method{+ removeCall(call: ICallList): boolean}}\\
	Metodo usato per rimuovere un partecipate alla chiamata. La rimozione avviene invocando il metodo \texttt{remove(call)} dall'attributo \memberdata{calls}. Il metodo ritorna un valore booleano per indicare l'esisto dell'operazione: \inglese{false} in presenza di errori, \inglese{true} altrimenti.
\end{description}

\classsection{ICallList}

\subsubsection*{Funzione}
Interfaccia rappresentante una \textit{entry} di uno storico chiamate di un utente del sistema mytalk.

\subsubsection*{Relazioni d'uso}
\begin{itemize}
	\item \classname{IUserData}: l'interfaccia \classname{ICallList} definisce più metodi che restituiscono oggetti aventi tipo di ritorno \classname{IUserData}, essi sono i metodi \inglese{get} per ottenere l'utente che ha effettuato la chiamta. Analogamente \classname{IUserData} viene usato come parametro d'ingresso per i metodi \inglese{set} collegati ai metodi già citati.
	\item \classname{ICall}: l'interfaccia \classname{ICallList} definiscono metodi che restituiscono oggetti aventi tipo di ritorno \classname{ICall}.
\end{itemize}

\subsubsection*{Classi estese ed interfacce implementate}
\begin{itemize}
	\item \texttt{org.softwaresynthesis.mytalk.server.IMyTalkObject}: interfaccia da estendere. Ogni oggetto che implementerà l'interfaccia \classname{ICallList} dovrà essere in grado di convertire il proprio contenuto informativo in formato \textit{Json}.
\end{itemize}

\subsubsection*{Metodi}
\begin{description}
	\item{\method{+ getId(): Long}}\\
	Restituisce l'identificativo univoco di una \inglese{entry} di uno storico chiamate.
	
	\item{\method{+ setIdCall(call: ICall): void}}\\
	Imposta l'id della chiamta \classname{ICall} (passato come parametro d'ingresso).
	\item{\method{+ getIdCall(): IUserData}}\\
	Restituisce un id di un oggetto avente tipo \classname{ICall} rappresentante una chiamata effettuata dall'utente.
	
	\item{\method{+ setIdUser(contact: IUserData): void}}\\
	Imposta l'id dell'utente \classname{IUserData} (passato come parametro d'ingresso) come l'utente che ha effettuato la chiamata.
	\item{\method{+ getIdUser(): IUserData}}\\
	Restituisce l'id di oggetto avente tipo \classname{IUserData} rappresentante l'utente che ha effettuato la chiamata.
	
	\item{\method{+ setCaller(caller: boolean): void}}\\
	Imposta l'utente che ha effettuato la chiamata come chiamante se il parametro ricevuto ha valore true.
	\item{\method{+ getCaller(): boolean}}\\
	Ritorna un valore di tipo bool se l'utente della chiamata è colui che l'ha iniziata o meno.
\end{description}




\classsection{CallList}

\subsubsection*{Funzione}
Classe che implementa l'interfaccia \classname{ICallList}.

\subsubsection*{Relazioni d'uso}

Nessuna relazione evidenziata.

\subsubsection*{Classi estese ed interfacce implementate}
\begin{itemize}
	\item \texttt{call.ICallList}: interfaccia implementata dalla classe.
\end{itemize}

\subsubsection*{Attributi}

\begin{itemize}
	\item{\memberdata{-- id: long}}
	Attributo contenente il codice identificativo dello storico chiamata.
	\item{\memberdata{-- idCall: long}}
	Attributo contenente il codice identificativo della chiamata.
	\item{\memberdata{-- idUser: long}}
	Attributo contenente il codice identificativo dell'utente partecipante alla chiamata.
	\item{\memberdata{-- caller: boolean}}
	Attributo che identifica se l'utente è il chiamante o meno.
\end{itemize}

\subsubsection*{Metodi}
\begin{description}
	\item{\method{+ getId(): Long}}\\
	Restituisce il valore dell'attributo \memberdata{id}.
	
	\item{\method{+ setIdCall(call: Long): void}}\\
	Imposta il valore dell'attributo \memberdata{idCall} con il valore ricevuto da parametro ``call''.
	\item{\method{+ getIdCall(): Long}}\\
	Restituisce il valore dell'attributo \memberdata{idCall}.
	
	\item{\method{+ setIdUser(contact: Long): void}}\\
	Imposta il valore dell'attributo \memberdata{idUser} con il valore ricevuto da parametro ``contact''.
	\item{\method{+ getIdUser(): Long}}\\
	Restituisce il valore dell'attributo \memberdata{idUser}.
	
	\item{\method{+ setCaller(caller: boolean): void}}\\
	Imposta il valore dell'attributo \memberdata{caller} con il valore ricevuto da parametro ``caller''.
	\item{\method{+ getCaller(): boolean}}\\
	Restituisce il valore dell'attributo \memberdata{caller}.
\end{description}

\subsection{Package org.softwaresynthesis.mytalk.server.call.controller}\label{sec:callServlet}

\classsection{GetCallsController}

\subsubsection*{Funzione}
\texttt{Controller} da richiamare per effettuare il download della lista che rappresenta lo storico delle chiamate effettuate e ricevute dell'utente. I dati vengono ritornati sotto forma di stringa in formato JSON.

\subsubsection*{Relazioni d'uso}
\begin{itemize}
	\item \texttt{java.io.IOException}: eccezione richiamabile dal metodo\method{doAction()}.
	\item \texttt{java.io.PrintWriter}: classe istanziata all'interno del metodo \method{doAction()}. Usata per scrivere l'output del \inglese{controller}.
	\item \texttt{javax.servlet.ServletException}: eccezione sollevabile dai metodi \method{doAction()}.
	\item \texttt{javax.servlet.http.HttpServletRequest}: classe usata per interagire con le richieste \texttt{ajax} inoltrate dal client.
	\item \texttt{javax.servlet.http.HttpServletResponse}: classe usata per interagire con le richieste \texttt{ajax} inoltrate dal client.
	\item \classname{call.ICall}: usata per definire un oggetto rappresentante una chiamata.
	\item \classname{abook.IUserData}: usata per definire un utente generico, che ha partecipato alla chiamata.
\end{itemize}

\subsubsection*{Classi estese ed interfacce implementate}
\begin{itemize}
	\item \texttt{server.AbstractController}: classe estesa.
\end{itemize}

\subsubsection*{Attributi}

Nessun attributo evidenziato

\subsubsection*{Metodi}
\begin{description}
	
	\item{\method{\# doAction(request: HttpServletRequest, response: HttpServletResponse): void}}\\
	Metodo usato per accedere alla funzionalità della \texttt{controller}. Il procedimento da definire passa attraverso 3 step essenziali:
	\begin{itemize}
		\item \texttt{Inizializzazione}: il metodo definisce ed apre una nuova sessione a partire dalla richiesta request passata come parametro;
		\item \texttt{Elaborazione dati}: il metodo interroga il database attraverso la classe \classname{DataPersistanceManager}. Per ogni oggetto \classname{ICall} in cui si evidenzia che uno dei due partecipanti è il client (\classname{IUserData}) che ha inoltrato la richiesta al \texttt{front controller}, si aggiunge alla stringa result da ritornare i dati di tale chiamata. Nello specifico la formattazione della stringa di ritorno, dovrà descrivere la seguente logica:\\
		
		\verb|{name:"NOME_UTENTE",start="DATA_INIZIO",end="DATA_FINE"}|\\
		
		Si consideri che la dove vi è un nome in maiuscolo, vi dovrà essere il dato inerente estrapolato dall'istanza \classname{ICall} presa in considerazione.
		\item \texttt{Restituzione risultato}: dopo aver elaborato i dati, il \texttt{controller} dovrà scrivere nel \texttt{PrintWriter} associato alla request:
			\begin{itemize}
				\item una stringa formattata come sopra, se l'utente ha effettuato o ricevuto almeno una chiamata;
				\item null, altrimenti.
			\end{itemize}
	\end{itemize}

\end{description}

\subsection{Package org.softwaresynthesis.mytalk.server.message}\label{sec:message}

\classsection{IMessage}

\subsubsection*{Funzione}
Interfaccia che rappresenta un messaggio di segreteria del sistema \caName.

\subsubsection*{Relazioni d'uso}
\begin{itemize}
	\item \texttt{java.util.Date}: tipo utilizzato per definire la data in cui è stato registrato un messaggio.
	\item \classname{abook.IUserData}: usata per rappresentare il mittente e il destinatario del messaggio di segreteria.
\end{itemize}

\subsubsection*{Classi estese ed interfacce implementate}

\begin{itemize}
	\item \texttt{org.softwaresynthesis.mytalk.server.IMyTalkObject}: interfaccia estesa. Utilizzata per dire che la classe è da considerarsi come un \texttt{transfer object} gestito da \texttt{Hibernate}.
\end{itemize}

\subsubsection*{Metodi}
\begin{description}
	\item{\method{+ getId(): Long}}\\
	Restituisce l'identificativo univoco del messaggio.
	\item{\method{+ getSender(): IUserData}}\\
	Restituisce un istanza di tipo classname{abook.IUserData} che rappresenta il mittente del messaggio.
	\item{\method{+ setSender(sender: IUserData): void}}\\
	Imposta il mittente del messaggio.
	\item{\method{+ getReceiver(): IUserData}}\\
	Restituisce un istanza di tipo classname{abook.IUserData} che rappresenta il destinatario del messaggio.
	\item{\method{+ setReceiver(receiver: IUserData): void}}\\
	Imposta il destinatario del messaggio.
	\item{\method{+ getNewer(): boolean}}\\
	Restituisce un valore booleano che identifica lo stato del messaggio (``già letto'' se il valore ritornato è \textit{true}, ``da leggere'' se il valore ritornato è \textit{false}).
	\item{\method{+ setNewer(status: boolean): void}}\\
	Imposta il messaggio come ``già ascoltato'' o come ``da ascoltare''.
	\item{\method{+ getVideo(): boolean}}\\
	Restituisce un booleano che determina se si tratta di un messaggio audio oppure audio-video.
	\item{\method{+ setVideo(video: boolean): void }}\\	
	Imposta il messaggio come messaggio audio-video o come messaggio contenente solamente una traccia audio.
	\item{\method{+ getDate(): Date}}\\
	Restituisce la data in cui il mittente ha lasciato il messaggio nella segreteria del destinatario.
	\item{\method{+ setDate(date: Date): void}}\\
	Imposta la data in cui il mittente ha lasciato il messaggio nella segreteria del destinatario.
\end{description}

\classsection{Message}

\subsubsection*{Funzione}
Classe che implementa l'interfaccia \classname{IMessage}.

\subsubsection*{Relazioni d'uso}
\begin{itemize}
	\item \texttt{java.util.Date}: tipo utilizzato per definire la data in cui è stato registrato un messaggio.
	\item \classname{abook.IUserData}: usata per rappresentare il mittente e il destinatario del messaggio di segreteria.
\end{itemize}

\subsubsection*{Classi estese ed interfacce implementate}
\begin{itemize}
	\item \classname{message.IMessage}: interfaccia d'implementazione.
\end{itemize}

\subsubsection*{Attributi}

\begin{itemize}
	\item{\memberdata{-- id: long}}
	Attributo contenente il codice identificativo del messaggio.
	\item{\memberdata{-- sender: IUserData}}
	Attributo contenente lo \classname{IUserData} che rappresenta il destinatario del messaggio.
	\item{\memberdata{-- receiver: IUserData}}
	Attributo contenente lo \classname{IUserData} che rappresenta il destinatario del messaggio.
	\item{\memberdata{-- status: boolean}}
	Attributo contenente un valore booleano che identifica se il messaggio è stato visionato/ascoltato o meno. L'attributo se impostato a \textit{true}, designa il messaggio come \textit{ascoltato}. Invece se impostato a \textit{false} identifica il messaggio come \textit{ancora da ascoltare}.
	\item{\memberdata{-- video: boolean}}
	Attributo che stabilisce se il messaggio contiene o meno una traccia video. Si ricorda che ogni messaggio contiene (di \inglese{default}) una traccia audio. Se tale attributo è impostato a \textit{true} allora il messaggio contiene una traccia video.
	\item{\memberdata{-- date: Date}}
	Attributo che definisce l'orario di invio del messaggio.
\end{itemize}

\subsubsection*{Metodi}
\begin{description}
	\item{\method{+ getId(): Long}}\\
	Restituisce l'identificativo univoco del messaggio, ritornando il contenuto di \memberdata{id}.
	\item{\method{+ getSender(): IUserData}}\\
	Restituisce un istanza di tipo classname{abook.IUserData} che rappresenta il mittente del messaggio. Nello specifico il metodo ritorna l'attributo \memberdata{sender}.
	\item{\method{+ setSender(sender: IUserData): void}}\\
	Imposta il mittente del messaggio, sovrascrivendo il contenuto dell'attributo \memberdata{sender}.
	\item{\method{+ getReceiver(): IUserData}}\\
	Restituisce un istanza di tipo classname{abook.IUserData} che rappresenta il destinatario del messaggio. Nello specifico il metodo ritorna l'attributo \memberdata{receiver}.
	\item{\method{+ setReceiver(receiver: IUserData): void}}\\
	Imposta il destinatario del messaggio, sovrascrivendo il contenuto dell'attributo \memberdata{receiver}.
	\item{\method{+ getNewer(): boolean}}\\
	Metodo che ritorna il contenuto dell'attributo \memberdata{status}.
	\item{\method{+ setNewer(status: boolean): void}}\\
	Imposta il messaggio come ``già ascoltato'' o come ``da ascoltare'', sovrascrivendo il contenuto dell'attributo \memberdata{status}.
	\item{\method{+ getVideo(): boolean}}\\
	Restituisce il contenuto dell'attributo \memberdata{video}.
	\item{\method{+ setVideo(video: boolean): void }}\\	
	Metodo usato per impostare la ``natura'' del messaggio, impostando il contenuto dell'attributo \memberdata{video} mettendolo a \textit{true} se il messaggio contiene una traccia video, oppure a \textit{false} se non la contiene.
	\item{\method{+ getDate(): Date}}\\
	Restituisce la data in cui il mittente ha lasciato il messaggio nella segreteria del destinatario, ritornando il contenuto dell'attributo \memberdata{date}.
	\item{\method{+ setDate(date: Date): void}}\\
	Imposta la data in cui il mittente ha lasciato il messaggio nella segreteria del destinatario, sovrascrivendo il contenuto dell'attributo \memberdata{date}.

\end{description}

\subsection{Package org.softwaresynthesis.mytalk.server.message.controller}\label{sec:messageServlet}

\classsection{AddMessageController}

\subsubsection*{Funzione}
\inglese{Controller} da richiamare per inserire un messaggio nella segreteria di un utente.

\subsubsection*{Relazioni d'uso}
\begin{itemize}
	\item \texttt{java.io.IOException}: eccezione richiamabile dal metodo\method{doAction()}.
	\item \texttt{java.io.PrintWriter}: classe istanziata all'interno del metodo \method{doAction()}. Usata per scrivere l'output della del \inglese{controller}.
	\item \texttt{javax.servlet.ServletException}: eccezione sollevabile dai metodi \method{doAction()}.
	\item \texttt{javax.servlet.http.HttpServletRequest}: classe usata per interagire con le richieste \texttt{ajax} inoltrate dal client.
	\item \texttt{javax.servlet.http.HttpServletResponse}: classe usata per interagire con le richieste \texttt{ajax} inoltrate dal client.
	\item \classname{message.IMessage}: usata per definire un oggetto rappresentante un messaggio in segreteria.
	\item \classname{abook.IUserData}: usata per definire un utente generico, che ha partecipato alla chiamata.
\end{itemize}

\subsubsection*{Classi estese ed interfacce implementate}
\begin{itemize}
	\item \texttt{server.AbstractController}: classe estesa.
\end{itemize}

\subsubsection*{Attributi}

Nessun attributo evidenziato

\subsubsection*{Metodi}
\begin{description}
	
	\item{\method{\# doAction(request: HttpServletRequest, response: HttpServletResponse): void}}\\
	Metodo usato per accedere alla funzionalità del \inglese{controller}. Il procedimento da definire passa attraverso 3 step essenziali:
	\begin{itemize}
		\item \texttt{Inizializzazione}: il metodo definisce ed apre una nuova sessione a partire dalla richiesta request passata come parametro;
		\item \texttt{Elaborazione dati}: il metodo procede creando un \classname{IMessage} con i dati forniti dalla request e dopo aver aperto una \texttt{transaction} verso il database, esegue l'operazione \method{save()} sul relativo oggetto DAO. Al termine sarà necessario procedere con un operazione di \texttt{commit}. Nel realizzare il metodo, il programmatore incaricato dovrà tenere in considerazione la possibilità di fallimento durante l'esecuzione delle operazioni mediante la \texttt{transaction}. Ciò andrà gestito mediante un operazione di \texttt{rollback} atta ad eliminare una possibile inconsistenza dei dati.
		\item \texttt{Restituzione risultato}: dopo aver elaborato i dati, la \inglese{servlet} dovrà scrivere nel \texttt{PrintWriter} associato alla request:
			\begin{itemize}
				\item ``\inglese{true}'' se l'operazione è andata a buon fine;
				\item ``\inglese{false}'' altrimenti.
			\end{itemize}
	\end{itemize}

\end{description}

\classsection{DeleteMessageController}

\subsubsection*{Funzione}
\inglese{Controller} da richiamare per eliminare un messaggio.

\subsubsection*{Relazioni d'uso}
\begin{itemize}
	\item \texttt{java.io.IOException}: eccezione richiamabile dal metodo\method{doAction()}.
	\item \texttt{java.io.PrintWriter}: classe istanziata all'interno del metodo \method{doAction()}. Usata per scrivere l'output della del \inglese{controller}.
	\item \texttt{javax.servlet.ServletException}: eccezione sollevabile dai metodi \method{doAction()}.
	\item \texttt{javax.servlet.http.HttpServletRequest}: classe usata per interagire con le richieste \texttt{ajax} inoltrate dal client.
	\item \texttt{javax.servlet.http.HttpServletResponse}: classe usata per interagire con le richieste \texttt{ajax} inoltrate dal client.
	\item \classname{message.IMessage}: usata per definire un oggetto rappresentante un messaggio in segreteria.
	\item \classname{abook.IUserData}: usata per definire un utente generico, che ha partecipato alla chiamata.
\end{itemize}

\subsubsection*{Classi estese ed interfacce implementate}
\begin{itemize}
	\item \texttt{server.AbstractController}: classe estesa.
\end{itemize}

\subsubsection*{Attributi}

Nessun attributo evidenziato

\subsubsection*{Metodi}
\begin{description}

\item{\method{\# doAction(request: HttpServletRequest, response: HttpServletResponse): void}}\\
	Metodo usato per accedere alla funzionalità della \inglese{controller}. Il procedimento da definire passa attraverso 3 step essenziali:
	\begin{itemize}
		\item \texttt{Inizializzazione}: il metodo definisce ed apre una nuova sessione a partire dalla richiesta request passata come parametro;
		\item \texttt{Elaborazione dati}: il metodo procede creando un \classname{IMessage} con i dati forniti dalla request e dopo aver aperto una \texttt{transaction} verso il database, esegue un operazione di ricerca dell'instanza appena creata all'interno della lista di \classname{IMessage} ottenuta dal db. Quindi se si verifica che tale istanza è effettivamente presente nel database, il metodo esegue l'operazione \method{delete()} sul relativo oggetto DAO. Al termine sarà necessario procedere con un operazione di \texttt{commit}. Nel realizzare il metodo, il programmatore incaricato dovrà tenere in considerazione la possibilità di fallimento durante l'esecuzione delle operazioni mediante la \texttt{transaction}. Ciò andrà gestito mediante un operazione di \texttt{rollback} atta ad eliminare una possibile inconsistenza dei dati.
		\item \texttt{Restituzione risultato}: dopo aver elaborato i dati, il \texttt{controller} dovrà scrivere nel \texttt{PrintWriter} associato alla request:
			\begin{itemize}
				\item ``\inglese{true}'' se l'operazione è andata a buon fine;
				\item ``\inglese{false}'' altrimenti (compreso non solo il caso di errore dovuto ad un problema di connessione verso il db, ma anche quello relativo al mancato match del elemento ricercato con quelli presenti nel database).
			\end{itemize}
	\end{itemize}

\end{description}

\classsection{UpdateStatusMessageController}

\subsubsection*{Funzione}
\inglese{Controller} da richiamare per modificare lo stato di un messaggio. L'idea alla base è che un messaggio può trovarsi in uno dei seguenti due stati:
\begin{itemize}
	\item letto;
	\item non letto;
\end{itemize}
Il \inglese{controller} permette di effettuare una transizione di stato da ``non letto'' a ``letto''.

\subsubsection*{Relazioni d'uso}
\begin{itemize}
	\item \texttt{java.io.IOException}: eccezione richiamabile dal metodo\method{doAction()}.
	\item \texttt{java.io.PrintWriter}: classe istanziata all'interno del metodo \method{doAction()}. Usata per scrivere l'output della del \inglese{controller}.
	\item \texttt{javax.servlet.ServletException}: eccezione sollevabile dai metodi \method{doAction()}.
	\item \texttt{javax.servlet.http.HttpServletRequest}: classe usata per interagire con le richieste \texttt{ajax} inoltrate dal client.
	\item \texttt{javax.servlet.http.HttpServletResponse}: classe usata per interagire con le richieste \texttt{ajax} inoltrate dal client.
	\item \classname{message.IMessage}: usata per definire un oggetto rappresentante un messaggio in segreteria.
	\item \classname{abook.IUserData}: usata per definire un utente generico, che ha partecipato alla chiamata.
\end{itemize}

\subsubsection*{Classi estese ed interfacce implementate}
\begin{itemize}
	\item \texttt{server.AbstractController}: classe estesa.
\end{itemize}

\subsubsection*{Attributi}

Nessun attributo evidenziato


\subsubsection*{Metodi}
\begin{description}
	\item{\method{\# doAction(request: HttpServletRequest, response: HttpServletResponse): void}}\\
	Il procedimento da definire passa attraverso 3 step essenziali:
	\begin{itemize}
		\item \texttt{Inizializzazione}: il metodo definisce ed apre una nuova sessione a partire dalla richiesta request passata come parametro;
		\item \texttt{Elaborazione dati}: il metodo procede creando un \classname{IMessage} con i dati forniti dalla request e dopo aver aperto una \texttt{transaction} verso il database, esegue un operazione di ricerca dell'instanza appena creata all'interno della lista di \classname{IMessage} ottenuta dal db. Quindi se si verifica che tale istanza è effettivamente presente nel database, il metodo esegue l'operazione \method{update()} sul relativo oggetto DAO, andando a modificare lo stato in di lettura del messaggio. Al termine sarà necessario procedere con un operazione di \texttt{commit}. Nel realizzare il metodo il programmatore incaricato dovrà tenere in considerazione la possibilità di fallimento durante l'esecuzione delle operazioni mediante la \texttt{transaction}. Ciò andrà gestito mediante un operazione di \texttt{rollback} atta ad eliminare una possibile inconsistenza dei dati.
		\item \texttt{Restituzione risultato}: dopo aver elaborato i dati, il \inglese{controller} dovrà scrivere nel \texttt{PrintWriter} associato alla request:
			\begin{itemize}
				\item ``\inglese{true}'' se l'operazione è andata a buon fine;
				\item ``\inglese{false}'' altrimenti.
			\end{itemize}
	\end{itemize}

\end{description}

\classsection{GetMessagesController}

\subsubsection*{Funzione}
\inglese{Controller} da richiamare per scaricare la lista di messaggi attualmente presenti nella propria segreteria.

\subsubsection*{Relazioni d'uso}
\begin{itemize}
	\item \texttt{java.io.IOException}: eccezione richiamabile dal metodo\method{doAction()}.
	\item \texttt{java.io.PrintWriter}: classe istanziata all'interno del metodo \method{doAction()}. Usata per scrivere l'output della del \inglese{controller}.
	\item \texttt{javax.servlet.ServletException}: eccezione sollevabile dai metodi \method{doAction()}.
	\item \texttt{javax.servlet.http.HttpServletRequest}: classe usata per interagire con le richieste \texttt{ajax} inoltrate dal client.
	\item \texttt{javax.servlet.http.HttpServletResponse}: classe usata per interagire con le richieste \texttt{ajax} inoltrate dal client.
	\item \classname{message.IMessage}: usata per definire un oggetto rappresentante un messaggio in segreteria.
	\item \classname{abook.IUserData}: usata per definire un utente generico, che ha partecipato alla chiamata.
\end{itemize}

\subsubsection*{Classi estese ed interfacce implementate}
\begin{itemize}
	\item \texttt{server.AbstractController}: classe estesa.
\end{itemize}

\subsubsection*{Attributi}

Nessun attributo evidenziato

\subsubsection*{Metodi}
\begin{description}
	\item{\method{\# doAction(request: HttpServletRequest, response: HttpServletResponse): void}}\\
	Il procedimento da definire passa attraverso 3 step essenziali:
	\begin{itemize}
		\item \texttt{Inizializzazione}: il metodo definisce ed apre una nuova sessione a partire dalla richiesta request passata come parametro;
		\item \texttt{Elaborazione dati}: quindi esegue un operazione di download degli oggetti \classname{IMessage} il cui possessore è l'utente che ha fatto richiesta di \inglese{download}. in seguito a tale operazione, sarà necessario popolare e formattare la stringa di ritorno come segue: per ogni messaggio presente nella lista si dovrà riportare la dicitura:\\
		
		\verb|{name:"NOME_UTENTE",state="LETTO/NON LETTO"|\\
		\verb|,url="INDIRIZZO_MESSAGGIO",data="DATA_REGISTRAZIONE_MESSAGGIO"}|\\
	
		\item \texttt{Restituzione risultato}: dopo aver elaborato i dati, il \inglese{controller} dovrà scrivere nel \texttt{PrintWriter} associato alla request:
			\begin{itemize}
				\item la sequenza
				\item null, altrimenti.
			\end{itemize}
	\end{itemize}

\end{description}

\subsection{Package org.softwaresynthesis.mytalk.server.authentication.security}\label{sec:authentication}

\classsection{ISecurityStrategy}

\subsubsection*{Funzione}
Interfaccia che identifica il comportamento di una strategia generica di crittografia dei dati.

\subsubsection*{Relazioni d'uso}
Nessuna relazione d'uso evidenziata.

\subsubsection*{Classi estese ed interfacce implementate}

Nessuna relazione d'uso evidenziata.

\subsubsection*{Metodi}
\begin{description}
	\item{\method{+ encrypt(plainText: String): String}}\\
	Cripta la stringa di testo ricevuta come parametro e la restituisce al chiamante.\\\\
	Il metodo può solleva eccezioni:
	\begin{itemize}
		\item \exception{IOException}: l'implementazione di tale metodo dovrà predisporre il lancio di un eccezione qualora il messaggio passato al metodo (plainText) sia vuoto.
	\end{itemize}
	\item{\method{+ decrypt(encryptedText: String)}}\\
	Decripta la stringa di testo ricevuta come parametro e la restituisce al chiamante.\\\\
	Il metodo può sollevare eccezioni:
	\begin{itemize}
		\item \exception{IOException}: l'implementazione di tale metodo dovrà predisporre il lancio di un eccezione qualora il messaggio passato al metodo (plainText) sia vuoto.
	\end{itemize}
\end{description}

\classsection{AESAlgorithm}

\subsubsection*{Funzione}
Implementazione della strategia di codifica/decodifica con l'uso dell'algoritmo AES a 128bit.

\subsubsection*{Relazioni d'uso}
\begin{itemize}
	\item \texttt{java.security.Key}: usata per creare un istanza di una chiave durante il processo di criptazione.
	\item \texttt{javax.crypto.Cipher}: usata per creare un istanza di un oggetto di criptazione che implementa l'algoritmo AES a 128 bit.
	\item \texttt{javax.crypto.spec.SecretKeySpec}: usata dall'algoritmo di per costruire una chiave segreta a partire da un array di byte.
	\item \texttt{sun.misc.BASE64Encoder}: utilizzata per eseguire una trasformazione da stringa a byte.
	\item \texttt{sun.misc.BASE64Decoder}: utilizzata per eseguire una trasformazione da byte in stringa.
\end{itemize}

\subsubsection*{Classi estese ed interfacce implementate}
\begin{itemize}
	\item \classname{ISecurityStrategy}: interfaccia d'implementazione.
\end{itemize}

\subsubsection*{Attributi}
\begin{description}
  \item{\memberdata{-- \{frozen\} \underline{keyValue}: byte[]}}\\
  Array di byte usato per definire il valore della chiave su cui si basa l'algoritmo AES di criptazione. La chiave effettiva sarà creata a partire da tale attributo, per mezzo della classe \texttt{javax.crypto.spec.SecretKeySpec}.
  \item{\memberdata{-- \{frozen\} \underline{algorithm}: String}}\\
  Stringa costante che identifica il nominativo dell'algoritmo usato, e che dovrà essere specificato durante l'uso di \texttt{javax.crypto.Cipher}. l'attributo avrà valore ``AES''.
\end{description}

\subsubsection*{Metodi}
\begin{description}
	\item{\method{- generateKey(): Key}}\\
	Metodo che restituisce una chiave di tipo \texttt{java.security.Key}, a partire dall'array \memberdata{keyValue}. Per fare ciò, il metodo usa il costruttore di \texttt{javax.crypto.spec.SecretKeySpec}.
	Il metodo può sollevare eccezioni:
	\begin{itemize}
		\item \exception{Exception}: il metodo può sollevare un'eccezione generica.
	\end{itemize}
	\item{\method{+ encrypt(plainText: String): String}}\\
	Metodo usato per criptare un testo di tipo \texttt{String} passato come parametro d'ingresso. Il metodo usa una chiave ottenuta a partire dal metodo \method{generateKey()} in associazione alla classe \texttt{javax.crypto.Cipher} per criptare il testo tramite l'algoritmo AES.
	\begin{itemize}
		\item \exception{Exception}: il metodo solleva un'eccezione.
	\end{itemize}
	\item{\method{+ decrypt(encryptedText: String)}}\\
	Decripta la stringa di testo ricevuta come parametro, attuando una procedura inversa a quella presentata nel metodo \method{encrypt(plainText: String)}\\\\
	Il metodo può sollevare eccezioni:
	\begin{itemize}
		\item \exception{Exception}: il metodo può sollevare un'eccezione generica.
	\end{itemize}
\end{description}

\classsection{AESTemplate}

\subsubsection*{Funzione}
Rappresenta la struttura dell'algoritmo di crittografia AES a 128 bit. Implementa \texttt{template method}.

\subsubsection*{Relazioni d'uso}
\begin{itemize}
	\item \texttt{java.security.Key};
	\item \texttt{javax.crypto.Cipher};
	\item \texttt{javax.crypto.spec.SecretKeySpec};
\end{itemize}

\subsubsection*{Classi estese ed interfacce implementate}

Nessuna relazione evidenziata.

\subsubsection*{Attributi}
\begin{description}
  \item{\memberdata{-- key: byte[]}}\\
  Array di byte usato per definire il valore della chiave su cui si basa l'algoritmo AES di criptazione. La chiave effettiva sarà creata a partire da tale attributo, per mezzo della classe \texttt{javax.crypto.spec.SecretKeySpec}.
  \item{\memberdata{-- cipher: Cipher}}\\
  Oggetto rappresentante l'algoritmo di criptazione.
  \item{\memberdata{-- generated: Key}}\\
  Oggetto rappresentante la chiave generata dall'algoritmo di criptazione.
\end{description}

\subsubsection*{Metodi}
\begin{description}

	\item{\method{+ AESTemplate()}}\\
	Costruttore pubblico. imposta key con il seguente contenuto:\\
	\verb|new byte[]{'C', 'p', '2', 'Q', 'j', 'w', 'M', 'F', '7', 'e', 'j', 'N', 't', 'd', 'b', '1'};|.

	\item{\method{+ \{frozen\} execute(text: String): String}}\\
	Metodo che definisce l'iter di esecuzione di una generica procedura riguardante il sistema di sicurezza. Il metodo si occupa di inizializzare il cifratore \memberdata{cipher} mediante il metodo \texttt{getIstance(''AES'')}. Quindi si imposta il contenuto di \memberdata{generated} con ciò che viene restituito da \texttt{SecretKeySpec(this.key, ''AES'')}. Il metodo invoca \method{compleateAlgorithm}, e salva il suo valore di ritorno in una Stringa denominata result. Il metodo termina restituendo result

	\item{\method{\# \{abstract\} compleateAlgotihm(text: String): String}}\\
	Metodo astratto. Costituisce l'implementazione del pattern \texttt{template method}.

	\item{\method{\# getCipher(): Cipher}}\\
	Metodo che restituisce il cifratore \memberdata{cipher} alle sottoclassi in modo che possano criptare/decriptare il testo.

	\item{\method{\# getGenerateKey(): Key}}\\
	Metodo che restituisce una chiave di tipo \texttt{java.security.Key}, a partire dall'attributo \memberdata{generated}.
	
\end{description}

\classsection{AESEncode}

\subsubsection*{Funzione}
Rappresenta la struttura dell'algoritmo di criptazione usato nel sistema. Ridefinisce il metodo astratto \method{compleateAlgorithm} della classe padre \classname{AESTemplate}.

\subsubsection*{Relazioni d'uso}
\begin{itemize}
	\item \texttt{javax.crypto.Cipher}: usato come cifratore durante la procedura di criptazione;
	\item \texttt{sun.misc.BASE64Encoder}: encoder di criptazione.
\end{itemize}

\subsubsection*{Classi estese ed interfacce implementate}

\begin{itemize}
	\item \classname{AESTemplate}: classe estesa.
\end{itemize}

\subsubsection*{Attributi}

Nessun attributo evidenziato.

\subsubsection*{Metodi}
\begin{description}

	\item{\method{\# \{final\} compleateAlgotihm(text: String): String}}\\
	Metodo usato per definire l'algoritmo di criptazione. Il metodo definisce un oggetto di tipo \texttt{BASE64Encoder} denominato encoder e lo inizializza con il suo costruttore di default. Quindi si definisce un \texttt{array} di byte (buffer) in cui si dovrà caricare il contenuto del testo da criptare. Si osservi che per tale operazione è richiesto l'utilizzo del metodo \texttt{getBytes()} eseguito a partire da text. Quindi il metodo prosegue definendo  un cifratore di tipo \texttt{Cipher}. Una stringa result sarà usata per restituire il risultato dell'operazione di criptazione. Il prossimo passaggio definisce il cuore della procedura di criptazione, ed è importante che lo sviluppatore si attenga scrupolosamente a quanto qui riportato. Il metodo definisce le istruzioni:\\
	
	\verb|cipher.init(Cipher.ENCRYPT_MODE, super.getGeneratedKey());|\\
	\verb|buffer = cipher.doFinal(buffer);|\\
	\verb|result = encoder.encode(buffer);|\\
	
	Il metodo termina ritornando il contenuto di result.
	
\end{description}

\classsection{AESDecode}

\subsubsection*{Funzione}
Rappresenta la struttura dell'algoritmo di criptazione usato nel sistema. Ridefinisce il metodo astratto \method{compleateAlgorithm} della classe padre \classname{AESTemplate}.

\subsubsection*{Relazioni d'uso}
\begin{itemize}
	\item \texttt{javax.crypto.Cipher}: usato come cifratore durante la procedura di decriptazione;
	\item \texttt{sun.misc.BASE64Decoder}: decoder di decriptazione.
\end{itemize}

\subsubsection*{Classi estese ed interfacce implementate}

\begin{itemize}
	\item \classname{AESTemplate}: classe estesa.
\end{itemize}

\subsubsection*{Attributi}

Nessun attributo evidenziato.

\subsubsection*{Metodi}
\begin{description}

	\item{\method{\# \{final\} compleateAlgotihm(text: String): String}}\\
	Metodo usato per definire l'algoritmo di decriptazione. Il metodo definisce un oggetto di tipo \texttt{BASE64Decoder} denominato decoder e lo inizializza con il suo costruttore di default. Quindi si definisce un \texttt{array} di byte (buffer) in cui si dovrà caricare il contenuto del testo da decriptare. Si osservi che per tale operazione è richiesto l'utilizzo del metodo \texttt{decodeBuffer(text)} eseguito a partire da decoder. Quindi il metodo prosegue definendo  un cifratore di tipo \texttt{Cipher}. Una stringa result sarà usata per restituire il risultato dell'operazione . Il prossimo passaggio definisce il cuore della procedura di decriptazione, ed è importante che lo sviluppatore si attenga scrupolosamente a quanto qui riportato. Il metodo definisce le istruzioni:\\
	
	\verb|cipher.init(Cipher.DECRYPT_MODE, super.getGeneratedKey());|\\
	\verb|buffer = cipher.doFinal(buffer);|\\
	\verb|result = new String(buffer);|\\
	
	Il metodo termina ritornando il contenuto di result.
	
\end{description}

\subsection{Package org.softwaresynthesis.mytalk.server.authentication}\label{sec:authentication}

\classsection{PrincipalImpl}

\subsubsection*{Funzione}
Oggetto che permette di identificate univocamente uno \texttt{IUserData} memorizzato nel database del sistema \caName.

\subsubsection*{Relazioni d'uso}
\begin{itemize}
	\item \texttt{java.io.Serializable}: interfaccia d'implementazione usata per rendere serializzabili le istanze della classe.
	\item \texttt{java.security.Principal}: interfaccia d'implementazione usata per rendere caratterizzate le istanze della classe.
\end{itemize}

\subsubsection*{Classi estese ed interfacce implementate}

Nessuna relazione evidenziata.

\subsubsection*{Attributi}
\begin{description}
  \item{\memberdata{-- \{frozen\} \underline{serialVersionUID}: long}}\\
  Identificativo univoco per la classe, usato al fine di rendere l'oggetto serializzabile.
  \item{\memberdata{-- mail: String}}\\
  Attributo che rappresenta l'indirizzo e-mail dell'utente.
\end{description}

\subsubsection*{Metodi}
\begin{description}
	\item{\method{+ PrincipalImpl(mail: String)}}\\
	Classe costruttore. crea un oggetto \texttt{PrincipalImpl} che permette di determinare univocamente lo \texttt{IUserData} che ha effettuato il login.
	\item{\method{+ getName(): String}}\\
	Restituisce l'elemento identificativo (mail dell'utente) dello \texttt{IUserData} che ha effettuato la procedura di login.
	\item{\method{+ equals(obj: Object): boolean}}\\
	Verifica l'uguaglianza di due oggetti \texttt{PrincipalImpl} sulla base di un confronto tra gli indirizzi mail degli utenti confrontati.
	\item{\method{+ hashCode(): int}}\\
	Restituisce il codice hash dell'oggetto di invocazione.				\item{\method{+ toString(): String}}\\
	Restituisce l'istanza dell'oggetto sotto forma di stringa. In particolare evidenziando l'indirizzo e-mail dell'utente.
\end{description}

\classsection{AuthenticationModule}

\subsubsection*{Funzione}
Modulo di autenticazione utilizzato dal sistema \caName.

\subsubsection*{Relazioni d'uso}
\begin{itemize}
	\item \texttt{java.io.IOException}: la classe è in grado di sollevare eccezioni relative a operazioni di IO.
	\item \texttt{java.security.Principal}: cinterfaccia che definisce i dati di autenticazione di un utente.
	\item \texttt{java.util.Map}: parametri passati ad \method{iniziaile()}.
	\item \texttt{java.util.Set}: la classe può definire istanze di Set, per memorizzare i dati trattati.
	\item \texttt{javax.security.auth.callback.Callback}: vettore utilizzato per contenere i dati di autenicazione.
	\item \texttt{javax.security.auth.callback.CallbackHandler}: parametri passati ad \method{iniziaile()}.
	\item \texttt{javax.security.auth.callback.NameCallback}: oggetto che contiene il name dell'utente da autenticare.
	\item \texttt{javax.security.auth.callback.PasswordCallback}: oggetto che contiene la password dell'utente da autenticare
	\item \texttt{javax.security.auth.callback.UnsupportedCallbackException}: la classe è in grado di sollevare eccezioni relative a operazioni di UnsupportedCallback.
	\item \texttt{javax.security.auth.login.FailedLoginException}: la classe è in grado di sollevare eccezioni relative a operazioni di FailedLoginException.
	\item \texttt{javax.security.auth.login.LoginException}: la classe è in grado di sollevare eccezioni relative a operazioni di LoginException.
	\item \texttt{javax.security.auth.Subject}: definisce il soggetto da autenticare.
	\item \classname{abook.IUserData}: la classe interagisce con istanze di oggetti identificabili come utenti.
	\item \classname{dao.UserDataDAO}: utilizzata per modificare il contenuto del database. Nello psecifico viene usata per interagire con la tabella UserData.
\end{itemize}

\subsubsection*{Classi estese ed interfacce implementate}
\begin{itemize}
	\item \texttt{javax.security.auth.spi.LoginModule}: interfaccia d'implementazione.
\end{itemize}

\subsubsection*{Attributi}
\begin{description}
  \item{\memberdata{-- login: boolean}}\\
  Attributo che determina se il login è avvenuto oppure no.
  \item{\memberdata{-- commit: boolean}}\\
  Attributo che determina se il commit è già stato eseguito.
  \item{\memberdata{-- handler: CallbackHandler}}\\
  Oggetto utilizzato per il caricamento delle credenzialie.
  \item{\memberdata{-- password: char[]}}\\
  Attributo che che contiene la password immessa dall'utente.
  \item{\memberdata{-- username: String}}\\
   Attributo che che contiene lo username immesso dall'utente.
  \item{\memberdata{-- principal: Principal}}\\
  Contiene la caratteristica autenticativa del subject. Nel sistema MyTalk tale proprietà è riservata al campo mail.
  \item{\memberdata{-- subject: Subject}}\\
  Attributo che definisce il soggetto da autenticare.
\end{description}

\subsubsection*{Metodi}
\begin{description}
	\item{\method{+ initialize(subject: Subject, handler: CallbackHandler, sharedState: Map, option: Map): void }}\\
	Il metodo ha il compito di inizializzare gli attributi privati della classe con i parametri d'input.

	\item{\method{+ login(): boolean throws LoginException}}\\
	Metodo richiamato per verificare un riscontro tra i dati passati per l'autenticazione e i dati presenti nel database. Il metodo deve appoggiarsi alla classe \classname{UserDataDAO} per le operazione di estrazione dati dal database.
	
	\item{\method{+ commit(): boolean throws LoginException}}\\
	Metodo richiamato direttamente dal framework se il login va a buon fine. Il suo scopo è quellodi inizializzare il subject con i relativi principle.
	
	\item{\method{+ abort(): boolean throws LoginException}}\\
	Metodo usato per bloccare la procedura di login.
	
	\item{\method{+ logout(): boolean throws LoginException}}\\
	Metodo per il logout del sistema. Il suo scopo è quello di eliminare il principle e il subject.

\end{description}

\classsection{CredentialLoader}

\subsubsection*{Funzione}
Permette di caricare le credenziali di autenticazione, fornite dall'utente, per preparare la fase di login.

\subsubsection*{Relazioni d'uso}
\begin{itemize}
	\item \texttt{javax.security.auth.callback.Callback}: usato durante la procedura di inserimento dati.
	\item \texttt{javax.security.auth.callback.NameCallback}: tipo di dato richiesto in ingresso per completare la parte di login.
	\item \texttt{javax.security.auth.callback.PasswordCallback}: tipo di dato richiesto in ingresso per completare la parte di login.
	\item \texttt{java.io.IOException}: eccezione che può essere sollevata dal metodo \texttt{handle} definito dall'interfaccia \texttt{javax.security.auth.callback.CallbackHandler}.
	\item \texttt{javax.security.auth.callback.UnsupportedCallbackException}: eccezione che può essere sollevata dal metodo \texttt{handle} definito dall'interfaccia \texttt{javax.security.auth.callback.CallbackHandler}.
\end{itemize}

\subsubsection*{Classi estese ed interfacce implementate}
\begin{itemize}
	\item \texttt{javax.security.auth.callback.CallbackHandler}: Interfaccia implementata dalla classe.
\end{itemize}

\subsubsection*{Attributi}
\begin{description}
  \item{\memberdata{-- credential: AuthenticationData}}\\
  Attributo contenente i dati da autenticare per la login dell'utente.
  \item{\memberdata{-- security: ISecurityStrategy}}\\
  Attributo che contiene un oggetto che definisce un algoritmo di criptazione per i dati. Necessario per criptare i dati di autenticazione.
\end{description}

\subsubsection*{Metodi}
\begin{description}
	\item{\method{+ CredentialLoader(credential: AuthenticationData, security: ISecurityStrategy)}}\\
	Costruttore pubblico. Crea un istanza con le credenziali fornite dall'utente (fornite in fase di login).

	\item{\method{+ handle(callbacks: Callback[]): void}}\\
	Effettua il caricamento e crittografa delle credenziali fornite dall'utente per la fase di login\\\\
	Il metodo può sollevare eccezioni:
	\begin{itemize}
		\item \exception{IOException}
		\item \exception{UnsupportedCallbackException}
	\end{itemize}

\end{description}

\classsection{Loader}

\subsubsection*{Funzione}
Oggetto usato per il caricamento delle credenziali di autenticazione.

\subsubsection*{Relazioni d'uso}
\begin{itemize}
	\item \texttt{java.io.IOException};
	\item \texttt{javax.servlet.http.HttpServletRequest};
	\item \texttt{org.softwaresynthesis.mytalk.server.authentication.security.ISecurityStrategy}: usata per definire l'algoritmo di criptazione.
\end{itemize}

\subsubsection*{Classi estese ed interfacce implementate}

\begin{itemize}
	\item \texttt{javax.security.auth.callback.Callback}: interfaccia d'implementazione.
\end{itemize}

\subsubsection*{Attributi}
\begin{description}
  \item{\memberdata{-- callback: Callback}}\\
  \item{\memberdata{-- strategy: ISecurityStrategy}}\\
\end{description}

\subsubsection*{Metodi}
\begin{description}
	\item{\method{+ Loader(callback: Callback)}}\\
	Costruttore pubblico. Crea il ''caricatore'' con la giusta istanza di callback.
	
	\item{\method{\# getCallback(): callback}}\\
	Metodo usato per ritornare l'attributo \memberdata{callback}.
	
	\item{\method{\# getISecurityStrategy(): ISecurityStrategy}}\\
	Metodo usato per ritornare l'attributo \memberdata{strategy}.
	
\end{description}

\classsection{NameLoader}

\subsubsection*{Funzione}
Caricatore per lo username. Ha il compito di caricare e preparare lo username fornito dall'utente per la successiva fase di autenticazione.

\subsubsection*{Relazioni d'uso}
\begin{itemize}
	\item \texttt{java.io.IOException};
	\item \texttt{javax.servlet.http.HttpServletRequest};
	\item \texttt{javax.security.auth.callback.NameCallback)};
\end{itemize}

\subsubsection*{Classi estese ed interfacce implementate}

\begin{itemize}
	\item \texttt{Loader}: classe estesa.
\end{itemize}

\subsubsection*{Attributi}

Nessun attributo evidenziato.

\subsubsection*{Metodi}
\begin{description}
	\item{\method{+ NameLoader()}}\\
	Costruttore pubblico che richiama il costruttore della super classe passando come parametro: \verb|new NameCallback("username")|.
	
	\item{\method{+ load(request: HttpServletRequest): void}}\\
	Il metodo crea un istanza di \texttt{NameCallback} a partire dalla callback ottenuta dalla chiamata \verb|super.getCallback()|. Quindi carica in una variabile locale di tipo String il nome dell'utente. Tale parametro si ottiene mediante una chiamata a metodo getParameter(''username'') invocata dall'istanza request. Quindi si imposta tale nome sull'istanza callback (usare \verb|callback.setName(username)|).
	
	\item{\method{+ getData(): String}}\\
	Metodo usato per ottenere il nome dell'utente. Il metodo ritorna il contenuto di una chiamata \verb|super.getCallback().getName()|.
	
\end{description}

\classsection{PasswordLoader}

\subsubsection*{Funzione}
Caricatore per la password. Ha il compito di caricare e preparare la password fornita dall'utente per la successiva fase di autenticazione.

\subsubsection*{Relazioni d'uso}
\begin{itemize}
	\item \texttt{java.io.IOException};
	\item \texttt{javax.servlet.http.HttpServletRequest};
	\item \texttt{javax.security.auth.callback.PasswordCallback)};
\end{itemize}

\subsubsection*{Classi estese ed interfacce implementate}

\begin{itemize}
	\item \texttt{Loader}: classe estesa.
\end{itemize}

\subsubsection*{Attributi}

Nessun attributo evidenziato.

\subsubsection*{Metodi}
\begin{description}
	\item{\method{+ PasswordLoader()}}\\
	Costruttore pubblico che richiama il costruttore della super classe passando come parametro: \verb|new PasswordCallback("password")|.
	
	\item{\method{+ load(request: HttpServletRequest): void}}\\
	Il metodo crea un istanza di \texttt{PasswordCallback} a partire dalla callback ottenuta dalla chiamata \verb|super.getCallback()|. Quindi carica in una variabile locale di tipo String la password dell'utente, ed usando l'algoritmo di sicurezza definito dall'istanza di tipo ISecurityStrategy, cripta la password. Tale parametro si ottiene mediante una chiamata a metodo getParameter(''password'') invocata dall'istanza request. Quindi si imposta tale nome sull'istanza callback (usare \verb|callback.setName(PASSWORD_CRIPTATA)|).
	
	\item{\method{+ getData(): String}}\\
	Metodo usato per ottenere la password dell'utente. Il metodo salva il contenuto di una chiamata \verb|super.getCallback().getPassword()|. Poi, prima di ritornare il valore, esegue la decriptazione mediante l'algoritmo ISecurityStrategy utilizzato.
	
\end{description}

\subsection{Package org.softwaresynthesis.mytalk.server.authentication.controller}\label{sec:autservlet}

%TODO da rivedere nella definizione dei metodi! 
\classsection{LogoutController}

\subsubsection*{Funzione}
\inglese{Controller} da richiamare per effettuare il logout dal sistema.

\subsubsection*{Relazioni d'uso}
\begin{itemize}
	\item \texttt{java.io.IOException}: eccezione richiamabile dal metodo\method{doAction()}.
	\item \texttt{java.io.PrintWriter}: classe istanziata all'interno del metodo \method{doAction()}. Usata per scrivere l'output della \inglese{servlet}.
	\item \texttt{javax.servlet.ServletException}: eccezione sollevabile dai metodi \method{doAction()}.
	\item \texttt{javax.servlet.http.HttpServletRequest}: classe usata per interagire con le richieste \texttt{ajax} inoltrate dal client.
	\item \texttt{javax.servlet.http.HttpServletResponse}: classe usata per interagire con le richieste \texttt{ajax} inoltrate dal client.
	\item \classname{abook.IUserData}: usata per definire un utente.
	\item \classname{dao.DataPersistanceManager}: usata per comunicare tramite Hibernate con la tabella UserData della base di dati.
\end{itemize}

\subsubsection*{Classi estese ed interfacce implementate}
\begin{itemize}
	\item \texttt{server.AbstractController}: classe estesa.
\end{itemize}

\subsubsection*{Attributi}

Nessun attributo evidenziato

\subsubsection*{Metodi}
\begin{description}
	\item{\method{\# doAction(request: HttpServletRequest, response: HttpServletResponse): void}}\\
	Metodo che costituisce il kernel logico di risposta della \inglese{servlet}. Il metodo inizia caricando i parametri ricevuti mediante \texttt{HttpServletRequest}. Una chiamata a tale \inglese{servlet} corrisponde ad una richiesta di logout e si attua impostando a false la sessione contenuta nell'oggetto HttopServletRequest request. Il metodo procede con la creazione di un istanza  di \texttt{LoginContext} denominata context. Quindi viene effettuata la memorizzazione dell'oggetto ritornato da una chiamata:
\\
\verb|session.getAttribute("LoginContext");|
\\

L'oggetto cosi ottenuto dovrà essere controllato, ovvero si dovrà accertare che il tipo dinamico è conforme al tipo \texttt{LoginContext} (il programmatore dovrà obbligatoriamente usare   la primitiva \texttt{istanceof}). Quindi si sovrascrive il contenuto di context con quanto ottenuto dalla chiamata a metodo sopracitata, e si richiama il metodo \texttt{logout()} a partire dall'oggetto context. Il metodo termina invalidando la sessione richiamando il metodo \texttt{invalidate()} di HttpSession.

\end{description}

\classsection{LoginController}

\subsubsection*{Funzione}
\inglese{Controller} da richiamare per effettuare l'autenticazione dell'utente.

\subsubsection*{Relazioni d'uso}
\begin{itemize}
\item \texttt{java.io.IOException}: eccezione richiamabile dal metodo\method{doAction()}.
	\item \texttt{java.io.PrintWriter}: classe istanziata all'interno del metodo \method{doAction()}. Usata per scrivere l'output della \inglese{servlet}.
	\item \texttt{javax.servlet.ServletException}: eccezione sollevabile dai metodi \method{doAction()}.
	\item \texttt{javax.servlet.http.HttpServletRequest}: classe usata per interagire con le richieste \texttt{ajax} inoltrate dal client.
	\item \texttt{javax.servlet.http.HttpServletResponse}: classe usata per interagire con le richieste \texttt{ajax} inoltrate dal client.
	\item \classname{abook.IUserData}: usata per definire un utente.
	\item \classname{dao.DataPersistanceManager}: usata per comunicare tramite Hibernate con la tabella \texttt{UserData} della base di dati.
\end{itemize}

\subsubsection*{Classi estese ed interfacce implementate}
\begin{itemize}
	\item \texttt{server.AbstractController}: classe estesa.
\end{itemize}

\subsubsection*{Attributi}

Nessun attributo evidenziato

\subsubsection*{Metodi}
\begin{description}
	\item{\method{\# doAction(request: HttpServletRequest, response: HttpServletResponse): void}}\\

Il metodo inizia caricando i parametri ricevuti mediante \texttt{HttpServletRequest}. Una chiamata a tale metodo corrisponde ad una richiesta di login e si attua impostando due campi String con i dati ottenuti da una chiamata a \texttt{getParameter(``username'')} e \texttt{getParameter(``password'')} . Quindi si crea un istanza di \texttt{HttpSession} tramite il metodo \texttt{getSession(true)} del parametro request. Il metodo prosegue controllando se l'username e la password precedentemente ottenute hanno valore diverso da null. Nel caso il flusso principale continua creando un'istanza di \classname{AutenthicationData} (denominata credential) a partire dai parametri username e password. Quindi si carica la path del file di configurazione in un apposita stringa \texttt{pathFileConfig}, tramite la chiamata a metodo:\\
	
\verb|System.getenv("MyTalkConfiguration") + "\\LoginConfiguration.conf"|.\\

Il flusso principale prosegue all'interno di un blocco \texttt{try-catch} creando:
	\begin{itemize}
		\item[•] un \classname{CredentialLoader} (loader);
		\item[•] un \classname{LoginContext} (context);
		\item[•] un \classname{dao.UserDataDao} (user);
	\end{itemize}
Quindi tramite context si esegue il login e si impostano gli attributi di sessione come segue: 

\texttt{session.setAttribute("LoginContext", context);}

Infine si carica in user un istanza di \classname{dao.UserDataDAO} ottenuta dalla chiamata a metodo \method{UserDataDAO.getByEmail()}, e restituendo user in un formato di formattazione \texttt{Json} (\texttt{user.toJson()}).

\end{description}

\classsection{RegisterController}

\subsubsection*{Funzione}
\inglese{Conroller} da richiamare per effettuare la registrazione al sistema MyTalk.

\subsubsection*{Relazioni d'uso}
\begin{itemize}
	\item \texttt{java.io.IOException}: eccezione richiamabile dal metodo\method{doAction()}.
	\item \texttt{java.io.PrintWriter}: classe istanziata all'interno del metodo \method{doAction()}. Usata per scrivere l'output della \inglese{servlet}.
	\item \texttt{javax.servlet.ServletException}: eccezione sollevabile dai metodi \method{doAction()}.
	\item \texttt{javax.servlet.http.HttpServletRequest}: classe usata per interagire con le richieste \texttt{ajax} inoltrate dal client.
	\item \texttt{javax.servlet.http.HttpServletResponse}: classe usata per interagire con le richieste \texttt{ajax} inoltrate dal client.
	\item \classname{abook.IUserData}: usata per definire un utente.
	\item \classname{dao.DataPersistanceManager}: usata per comunicare tramite Hibernate con la tabella \texttt{UserData} della base di dati.
	\item \classname{authentication.security.ISecurityStrategy}: usata per criptare/decriptare i dati da inviare/ricevere
	\item \classname{authentication.security.AESAlgorithm}: implementazione di \classname{authentication.ISecurityStrategy} usata dalla classe descritta.
	
\end{itemize}

\subsubsection*{Classi estese ed interfacce implementate}
\begin{itemize}
	\item \texttt{server.AbstractController}: classe estesa.
\end{itemize}

\subsubsection*{Attributi}

Nessun attributo evidenziato

\subsubsection*{Metodi}
\begin{description}
	\item{\method{\# doAction(request: HttpServletRequest, response: HttpServletResponse): void}}\\

Il metodo inizia caricando i parametri ricevuti mediante \texttt{HttpServletRequest}. Una chiamata a tale metodo corrisponde ad una richiesta di registrazione al sistema e si attua impostando i dati di registrazione:
	\begin{itemize}
		\item name;
		\item username;
		\item mail;
		\item password;
		\item path immagine;
		\item domanda segreta;
		\item risposta alla domanda segreta;
	\end{itemize}
	Per eseguire tali operazioni, il metodo predispone delle opportune variabili di tipo String e crea uno \classname{abook.IUserData} user. Quindi all'interno di un blocco try catch viene impostata una strategia di criptaggio dei dati (creando un istanza di \classname{authentication.security.AESAlgorithm}). Quindi vengono caricati i dati di registrazione nelle variabili precedentemente create. Per eseguire tali operazioni è necessario usare l'istruzione:\\
	
	\verb|request.getParameter("NOME_PARAMTERO")|\\
	
	dove i nomi dei parametri sono:
	\begin{itemize}
		\item username (usata per l'indirizzo mail);
		\item name;
		\item surname;
		\item password;
		\item answer;
		\item question;
		\item picturePath;
	\end{itemize}
	Il passo successivo (sempre dentro il blocco try) consiste nel creare un istanza di \classname{dao.UserDataDAO} userDAO, e nell'eseguire un operazione di userDAO.\method{insert()} passando come parametro user (l'utente precedentemente creato). Il blocco try termina impostando result a ``\inglese{true}''.
	
	Per quanto riguarda il blocco catch, in esso viene impostato result a ``\inglese{false}''.
	
	Il metodo termina scrivendo in un \texttt{PrintWriter} il valore di result.
\end{description}


\subsection{Package org.softwaresynthesis.mytalk.server.connection}\label{sec:connection}

\classsection{PushInbound}

\subsubsection*{Funzione}
Classe per la definizione di un apposito canale di comunicazione client-server. La classe è un'estensione di \texttt{org.apache.catalina.websocket.MessageInbound}. Si osservi che l'associazione tra un istanza di \classname{PushInbound} e un utente del sistema è univoca: fintanto che connesso un utente ha un proprio \classname{PushInbound} presente sul server.

\subsubsection*{Relazioni d'uso}
\begin{itemize}
	\item \texttt{java.io.IOException}: Eccezione che può sollevare il metodo \method{OnTextMessage()}.
	\item \texttt{java.nio.CharBuffer}: tipologia di \textit{buffer} usata per memorizzare il messaggio ricevuto come parametro d'ingresso.
	\item \texttt{java.util.Iterator}: usata per scorrere il contenuto dell'insieme di \classname{AddressBookEntry}
	\item \texttt{java.util.Set}: struttura dati usata per memorizzare le \inglese{entry} (\classname{AddressBookEntry}) che costituiscono la rubrica dell'utente.
	\item \texttt{com.google.gson.*}: converte i dati interni in una stringa formato json.
	\item \classname{State}: classe interna utilizzata per rappresentare lo stato dell'utente a cui è associata l'istanza di \classname{PushInbound}. Un utente ha uno stato fintanto che è connesso (quindi fintanto che esiste sul server un'istanza di \classname{PushInbound}  ad esso associata). Gli stati possibili sono: \texttt{available} (utente connesso e libero, quindi disposto a ricevere chiamate) e \texttt{occupied} (utente connesso ma attualmente occupato in un'altra conversazione).
	\item \classname{abook.AddressBookEntry}: usata per definire la rubrica dell'utente, richiamata dal metodo \method{onTextMessage} nel momento in cui si presenta la necessità di aggiornare lo stato dell'utente e renderlo visibile ai suoi contatti.
	\item \classname{abook.IUserData}: usata per riferire un istanza della classe ad un particolare utente.
	\item \classname{dao.DataPersistanceManger}: usata per riferire un istanza della classe ad un particolare utente.

\end{itemize}

\subsubsection*{Classi estese ed interfacce implementate}
\begin{itemize}
	\item \texttt{org.apache.catalina.websocket.MessageInbound}: classe da estendere, aggiunge alla classe \classname{PushInbound} le funzionalità necessarie per renderla un ``canale di comunicazione'' utilizzabile dai client.
\end{itemize}

\subsubsection*{Attributi}
\begin{description}
  \item{\memberdata{-- id: Long}}\\
  Identificativo di tipo \textit{Long} del canale di comunicazione associato ad un client univoco.
  \item{\memberdata{-- state: State}}\\
  Attributo usato per memorizzare lo stato dell'utente ``proprietario'' del \classname{PushInbound}.
\end{description}

%TODO da rivedere.

\subsubsection*{Metodi}
\begin{description}

	\item{\method{+ setId(n: Long): void}}\\
	Metodo per impostare il valore contenuto nell'attributo \memberdata{id}.
	
	\item{\method{+ getId(): Long}}\\
	Metodo che ritorna il contenuto dell'attributo \memberdata{id}.
	
	\item{\method{+ onTextMessage(message: CharBuffer): void}}\\
	Metodo invocato al momento della ricezione di un messaggio da parte del client. Il metodo riceve in input un oggetto di tipo \texttt{CharBuffer} contenente il messaggio inviato dal client.
Inizialmente il metodo crea un istanza per ognuno dei seguenti oggetti:
	\begin{itemize}
		\item[•]\texttt{Gson}: tipo di \textit{Json} definito da \textit{Google};
		\item[•]\texttt{JsonParser}: \textit{parser} usato per scorrere il contenuto di una stringa \textit{Json};
		\item[•]\texttt{JsonArray}: array popolato a partire dal \inglese{parsing} della stringa di messaggio data in input.
	\end{itemize}
	
	Prima di procedere si voglia considerare quanto segue: il metodo dopo aver ``segmentato'' il messaggio ricevuto in input, si occupa di esaminarne il contenuto che può essere di 5 tipologie, ciascuna identificata tramite un valore intero positivo da 1 a 5. Tale valore deve essere salvato nell'istanza di \texttt{JsonArray}
	Detto ciò, tornando a definire il flusso principale del metodo, si osservi che:

	\begin{itemize}
		\item[•]Se la richiesta inoltrata è del tipo 1: il metodo prende in lettura il messaggio e imposta il contenuto di \memberdata{id} con il valore letto mediante procedura di \inglese{parsing} analoga a quella definita al passo precedente. Per farlo utilizza il metodo \texttt{fromJson} richiamato dall'istanza \texttt{gsonObj} di tipo \texttt{Gson} creata al passo precedente. A tale metodo passa il contenuto dell'array e in particolare ciò che è salvato nella posizione 1 (utilizzo di metodo \texttt{get(int i)}).
		\item[•]Se la richiesta inoltrata è del tipo 2: il metodo procede con le istruzioni necessarie a scambiare i dati per la chiamata. Nello specifico viene salvato in un attributo di tipo \texttt{Long}, l'\texttt{id} del client che desidero chiamare, quindi ricerco l'oggetto \texttt{PushInbound} associato al client che desidero contattare, e gli inoltro il messaggio ricevuto come parametro d'ingresso.
		\item[•]Se la richiesta inoltrata è del tipo 3: il metodo comunica al client ``destinatario'' della chiamata, l'\texttt{id} del client chiamante. La procedura è analoga a quella identificata nel punto precedente, con la specifica che il messaggio inoltrato è l'identificativo del cliente che desidera avviare la chiamata.
		\item[•]Se la richiesta inoltrata è del tipo 4: il metodo si occupa dell'eliminazione della \textit{webSocket}.
		\item[•]Se la richiesta inoltrata è del tipo 5: il metodo viene usato per modificare lo stato dell'utente, con il valore passato tramite messaggio. Dopo la ricezione del messaggio, il metodo cambia il valore del campo \texttt{state} con il valore ricevuto come parametro d'ingresso. Quindi procede ricavando la lista degli utenti nella rubrica dell'utente che ha cambiato stato e comunica loro che è avvenuto un cambiamento di stato.
	\end{itemize}
\end{description}

\clearpage

\section{Specifica sotto-architettura \texttt{clientpresenter}}\label{sec:clientpresenterarchitecture}

La sotto-architettura \texttt{clientpresenter}\footnote{%
  Considerazioni in tutto e per tutto analoghe si applicano anche alla sotto-architettura \texttt{clientview} descritta nella sezione \vref{sec:clientviewarchitecture} del presente documento.
}
richiede una trattazione speciale a causa del particolare dominio tecnologico coinvolto. Essa è infatti definita con il linguaggio JavaScript che pur essendo definito come linguaggio orientato agli oggetti non permette di definire classi ed è, inoltre, debolmente tipizzato.

Al fine di facilitare al programmatore la comprensione del progetto (pensando quindi alle varie entità come a classi) senza però confonderlo in fase di stesura del codice, si è stabilito di adoperare la seguente terminologia:

\begin{description}
	\item{\scshape\bfseries Attributi}: saranno definiti con una sintassi simile a quella già usata per la parte server, pertanto al nome dell'attributo sarà associato il tipo ``logico'' che idealmente rappresenta. Dal momento che in JavaScript non esiste il controllo dei tipi, si è stabilito di attenersi a una notazione simile a quella utilizzata per lo strumento di generazione automatica della documentazione JSDoc. In particolare i tipi legali consentiti sono:
	\begin{itemize}
	  \item \texttt{HTMLElement} che corrispondono a nodi DOM già esistenti nella pagina HTML che definisce l'interfaccia utente oppure modificati (a \textit{run-time}) sulla base di informazioni ottenute dal server e che saranno indicati come:
	  	\begin{verbatim}
				(+, -) attributo: HTMLElement
			\end{verbatim}
	
    \item \texttt{String} se l'attributo è destinato a contenere esclusivamente valori di tipo \textit{String} che sarà segnato nel documento come:
			
			\begin{verbatim}
				(+, -) attributo: String
			\end{verbatim}
			
			\item \texttt{Array}, se l'attributo è un \inglese{array}, che in JavaScript non chiede di essere definito per tipo di valori contenibili che nel presente documento sarà segnalato come:
			
			\begin{verbatim}
				(+, -) attributo: Array
			\end{verbatim}
			
			\item \texttt{Number}, che è da considerarsi come il supertipo di tutti i tipi numerici e che sarà menzionato nel presente documento come:
			\begin{verbatim}
			  (+, -) attributo: Number
			\end{verbatim}
			
			\item \texttt{Boolean}, che rappresenta un valore booleano e che sarà indicato nel presente documento come:
			\begin{verbatim}
			  (+, -) attributo: Boolean
			\end{verbatim}
			
			\item \texttt{Object}, che è da intendersi come un oggetto generico JavaScript, dotato una serie di proprietà per ciascuna delle quali è definito un valore e che sarà rappresentato come:
			\begin{verbatim}
			  (+, -) attributo: Object
			\end{verbatim}
	\end{itemize}
		
Per quanto riguarda l'accessibilità degli attributi, invece, saranno distinti i due casi di accessibilità pubblica e accessibilità privata, indicati nel presente documento con un formalismo analogo a quello utilizzato per la parte server.

Nell'attività di codifica, i programmatori dovranno attenersi alle seguenti pratiche:
\begin{itemize}
  \item gli attributi pubblici saranno assimilati a proprietà dell'oggetto creato tramite la funzione costruttore accessibili all'esterno, e dovranno pertanto essere specificati come:
  \begin{verbatim}
    this.attributo = new ClassName();
  \end{verbatim}
  
  \item gli attributi privati invece, dal momento che saranno visibili solo all'interno del costruttore e di tutte le funzioni annidate all'interno di esso -- vale a dire i metodi della classe -- dovranno essere dichiarati come variabili locali secondo la consueta sintassi JavaScript:
  \begin{verbatim}
    var attributo = new ClassName();
  \end{verbatim}
\end{itemize}
		
	\item{\scshape\bfseries Metodi}: la sintassi usata dal programmatore per definire un metodo in JavaScript dipende dalla natura pubblica o privata del metodo stesso, in particolare:
	\begin{itemize}
	
\item i metodi pubblici possono essere assimilati in tutto e per tutto a delle proprietà che l'oggetto espone verso l'esterno, pertanto dovranno essere specificati come:
\begin{verbatim}
  this.doSomething = function(arg0, arg1, ...) {
          ...corpo del metodo...
  }
\end{verbatim}

\item i metodi privati, invece, essendo utilizzabili solo all'interno di altri metodi della classe, possono essere definiti come funzioni annidate all'interno del costruttore della classe stessa, secondo la consueta sintassi
\begin{verbatim}
  function doSomething(arg0, arg1, ...) {
        ...corpo del metodo...
  }
\end{verbatim}

La sintassi utilizzata all'interno del presente documento, in ogni caso, rimane quella che è stata definita in precedenza e prevede i marcatori \verb|+| per i metodi pubblici e \verb|-| per i metodi privati.

Per quanto concerne il tipo di ritorno, qualora non fosse previsto alcun valore da restituire, esso sarà indicato nel presente documento con la parola chiave \texttt{void}, analogamente a come sono stati trattati i metodi nella sotto-architettura \texttt{server}.
	\end{itemize}
\end{description}

\subsection{Package \texttt{org.softwaresynthesis.mytalk.clientpresenter.guicontrol}}\label{sec:guicontrol}

\classsection{AccountSettingsPresenter}
\subsubsection*{Funzione}
\inglese{Presenter} incaricato di gestire il pannello delle impostazioni e dei dati personali dell'utente visualizzando i dati inseriti in fase di registrazione e permettendo di modificarli in un secondo momento.

\subsubsection*{Relazioni d'uso}
Nessuna relazione d'uso evidenziata.

\subsubsection*{Classi estese ed interfacce implementate}
Nessuna relazione evidenziata.

\subsubsection*{Attributi}
\begin{description}
\item{\memberdata{-- thisPanel: HTMLElement}}\\
Attibuto che rappresenta un riferimento all'elemento HTML, da intendersi come nodo DOM, che rappresenta la radice del pannello associato a questo \inglese{presenter} ed è costituito da un \verb'<div>' avente come identificativo la stringa \verb'AccountSettingsPanel'.
\end{description}

\subsubsection*{Metodi}
\begin{description}

\item{\method{+ onShowAccountSettingPanel(): void}}\\
Questo gestore di eventi viene attivato al verificarsi dell'evento \verb'showAccountSettingPanel' e determina il caricamento del \inglese{template} che corriponde al pannello delle impostazioni utente, demandando in un secondo momento alla vista il compito di inizializzare il comportamento del pannello mediante il metodo \method{display()} che essa mette a disposizione.

\item{\method{+ sendUserData(data: Object): Boolean}}\\
Il metodo ha il compito di gestire la modifica dell'utente ai propri dati verificando innanzitutto se essi hanno subito qualche variazione rispetto a quanto noto in precedenza, mediante una chiamata al metodo \method{hasSomethingChanged}. In caso affermativo, dovrà essere inviata una richiesta all'indirizzo del \inglese{front controller} per trasferire la modifica sui dati salvati sul database del server.

A tale scopo dovrà essere creata una nuova istanza di \verb'XMLHttpRequest' per effettuare una richiesta AJAX sincrona, con particolare attenzione a utilizzare un'istanza di \inglese{form data} dal momento che dovranno essere spediti anche dati di natura non testuale, con una sintassi del tipo:
\begin{verbatim}
  var form = new FormData();
  formData.append("operation", "accountSettings");
        ...altri valori da inviare...
  var request = new XMLHttpRequest();
  request.open("POST", controllerURL, false);
  request.send(form);
\end{verbatim}

Da ultimo, il metodo dovrà restituire il valore ottenuto deserializzando il testo della risposta ottenuta dal server, che è \verb'true' se l'operazione è andata a buon fine e \verb'false' altrimenti.

\item{\method{-- hasSomethingChanged(data: Object): Boolean}}\\
L'oggetto passato come parametro in ingresso a questo metodo è costituito dalle proprietà \verb'name', \verb'surname' e \verb'picturePath' che rappresentano gli unici dati utente che sono suscettibili di possibili modifiche. Il metodo dovrà restituire \verb'true' se almeno una di queste proprietà ha un valore diverso rispetto alle informazioni che sono memorizzate nel client  e \verb'false' altrimenti.

\end{description}


% Qui dovrà andarci ADDRESSBOOKPRESENTER


\classsection{CallHistoryPresenter}

\subsubsection*{Funzione}
\inglese{Presenter} incaricato alla visualizzazione e alla gestione dello storico chiamate dell'utente.

\subsubsection*{Relazioni d'uso}
Nessuna relazione d'uso evidenziata.

\subsubsection*{Classi estese ed interfacce implementate}
Nessuna relazione evidenziata.

\subsubsection*{Attributi}
\begin{description}
  \item{\memberdata{-- calls: Array}}\\
  Array enumerativo che conterrà tutte le chiamate effettuate dall'utente recuperate dal server.
  \item{\memberdata{-- thisPanel: HTMLElement}}\\
  Elemento HTML, da intendersi come nodo DOM, che rappresenta la vista controllata da questo \inglese{presenter} e che corrisponde a un elemento \verb'<div>' avente come nome identificativo la stringa \verb'CallHistroryPanel'.
\end{description}

\subsubsection*{Metodi}
\begin{description}

  \item{\method{+ displayList(): void}}\\
  Metodo che ha il compito di popolare la lista delle chiamate in cui è coinvolto un determinato utente. Il metodo ha il compito di procurarsi la lista delle chiamate attraverso il server al fine di inizializzare il campo dati privato \verb'calls'.
  
  Quindi, per ognuna delle chiamate che sono state scaricate, ha il compito di aggiornare la vista corrispondente provocando la visualizzazione dei dati della chiamata.
  
  \item{\method{+ onShowCallHistoryPanel(): void}}\\
  Questo gestore di eventi ha il compito di provocare la visualizzazione del pannello contenente lo storico delle chiamate demandando alla vista il compito di costruirsi e caricare il relativo \inglese{template}.
  
    \item{\method{-- getCalls(): Array}}\\
  Metodo che ha il compito di recuperare tutta la lista delle chiamate effettuate e tracciate nel server, inviando una richiesta AJAX sincrona al \inglese{front controller} che nell'intestazione contiene il parametro \verb'operation' impostato a \verb'getCalls'.
  
  In particolare, il metodo ha il compito di creare una nuova istanza di \verb'XMLHttpRequest' con la \inglese{query string} specificata in precedenza e di inviarla al server quindi ritornare il testo della risposta ottenuta dopo averlo deserializzato e costruito l'array di chiamate corrispondente.

\end{description}


% Qui dovrà andarci COMMUNICATIONPRESENTER


% Qui dovrà andarci CONTACTPRESENTER


\classsection{GroupPresenter}

\subsubsection*{Funzione}
\inglese{Presenter} incaricato di gestire il pannello per la visualizzazione e l'amministrazione dei gruppi presenti nella rubrica.

\subsubsection*{Relazioni d'uso}
Nessuna relazione d'uso evidenziata.

\subsubsection*{Classi estese ed interfacce implementate}
Nessuna relazione evidenziata.

\subsubsection*{Attributi}
\begin{description}

  \item{\memberdata{-- contacts: Object}}\\
  Array associativo in cui sono memorizzati i contatti della rubrica, indicizzato in base al numero identificativo di ciascun contatto.
  
  \item{\memberdata{-- groups: Object}}\\
  Array associativo in cui sono memorizzati i gruppi presenti nella rubrica, indicizzato in base al numero identificativo di ciascun gruppo.
  
  \item{\memberdata{-- thisPanel: Object}}\\
  Riferimento all'elemento HTML, da intendersi come nodo DOM, che rappresenta il pannello controllato dal \inglese{presenter} ed è costituito da un \verb'<div>' con l'identificativo \verb'GroupPanel'.
  
\end{description}

\subsubsection*{Metodi}
\begin{description}

  \item{\method{+ displayList(): void}}\\
  Metodo che visualizza all'interno del pannello la lista dei gruppi presenti in rubrica. A tale scopo è necessario innanzitutto inizializzare i campi dati \memberdata{contacts} e \memberdata{groups}, quindi è possibile popolare la vista mediante una chiamata al metodo \method{addListItem} messo a disposizione dalla vista per ognuno dei gruppi che sono stati precedentemente recuperati.
  
  \item{\method{+ onShowGroupPanel(): void}}\\
  Questo gestore di eventi è associato all'evento con l'identificativo \verb'showGroupPanel' e determina la visualizzazione del pannello per l'amministrazione dei gruppi della rubrica caricando il relativo \inglese{template} e demandando in un secondo momento alla vista il compito di inizializzare le proprie componenti grafiche.
  
  \item{\method{+ selectCandidates(group: Object): Object}}\\
  Il metodo ha il compito di estrarre dall'insieme di tutti i contatti memorizzati nella rubrica quelli che non appartengono al gruppo passato come parametro e che, pertanto, sono dei candidati per l'aggiunta allo stesso.
  
  L'algoritmo mediante il quale dovrà essere raggiunto questo scopo consiste nell'effettuare una copia \textit{profonda} dell'array associativo dei contatti presenti in rubrica, che costituirà l'insieme dei candidati da restituire al chiamante.
  
  In seguito, si dovrà scorrere la proprietà \verb'contacts' del gruppo (che è un array enumerativo dei numeri identificativi dei contatti che vi appartengono) e rimuovere gli oggetti corrispondenti dall'array dei candidati con l'istruzione \verb'delete'.
  
  Al termine del \inglese{loop}, l'array associativo così ottenuto conterrà solo i contatti che non appartengono al gruppo come richiesto e potrà quindi essere restituito al chiamante.

\end{description}


\classsection{LoginPresenter}

\subsubsection*{Funzione}
\inglese{Presenter} incaricato di gestire il pannello di \textit{login}, ha il compito di recuperare i dati inseriti dall'utente (username e password) dall'interfaccia grafica, controllare se assumono dei valori legali e in seguito passarli al server per eseguire l'autenticazione.

Il \inglese{presenter} si occupa inoltre di gestire il caso in cui l'autenticazione non può essere portata a termine per errori nei dati di input e fornisce infine le funzionalità per il recupero della password.

\subsubsection*{Relazioni d'uso}

Nessuna relazione evidenziata.

\subsubsection*{Classi estese ed interfacce implementate}

Nessuna relazione evidenziata.

\subsubsection*{Attributi}
\begin{description}

\item{\memberdata{-- thisPanel: HTMLElement}}\\
Attributo che definisce il contenuto del nodo DOM inerente alla schermata di \inglese{login}, nello specifico tale nodo corrisponde ad un \verb+<div>+ il cui valore identificativo è uguale a \texttt{LoginPanel}.

\end{description}

\subsubsection*{Metodi}
\begin{description}

\item{\method{+ getUsername(): String}}\\
Metodo che ha il compito di recuperare tramite l'interfaccia della vista il nome utente inserito dall'utente, quindi effettuare un controllo preliminare per assicurare che il valore rientri all'interno del dominio dei valori legali.

Il controllo dovrà essere effettuato mediante un'opportuna espressione regolare del tipo:
\begin{verbatim}
  ^[A-Z0-9._%+-]+@[A-Z0-9.-]+\.[A-Z]{2,4}$
\end{verbatim}
e, nel caso in cui il valore non corrispondesse, il metodo dovrà sollevare un'eccezione contenente il testo ``indirizzo email non valido''. Dal momento che il metodo restituisce una stringa, in caso di valore accettabile quest'ultimo dovrà essere restituito al chiamante.

\item{\method{+ getPassword(): String}}\\
Metodo che ha il compito di recuperare tramite l'interfaccia della vista la password inserita dall'utente. Prima di restituire il valore ottenuto al chiamante, il metodo ha il compito di verificare che il valore sia valido e che, in particolare, non sia \verb+undefined+ oppure uguale alla stringa vuota. In simili circostanze, il metodo solleva un'eccezione contenente esattamente il testo ``password non specificata''.

\item{\method{+ getSecretQuestion(username: String): String}}\\
Il metodo a il compito di richiedere al server tramite una chiamata AJAX sincrona la domanda segreta associata all'utente che ha come nome identificativo la stringa passata come parametro, e restituirla al chiamante.

Al suo interno dovrà essere creata pertanto una nuova istanza di \verb+XMLHttpRequest+ destinata all'indirizzo del \inglese{front controller} presente sul server, e dovrà essere inviata  passando al metodo \verb+send+ la stringa:
\begin{verbatim}
  "operation=question&username=" + encodeURIComponent(username)
\end{verbatim}

Dal momento che il testo della risposta ricevuta dal server contiene la domanda segreta, sarà infine sufficiente restituire il valore ottenuto con \verb+request.responseText+.

\item{\method{+ login(data: Object): String}}\\
Il metodo prende come parametro un oggetto dotato delle due proprietà \verb+usename+ e \verb+password+ sulla base delle quali avrà il compito di costruire la \inglese{query string} da inviare al server per la verifica delle credenziali.

Il metodo dovrà pertanto costruire una \verb+XMLHttpRequest+ da inviare all'indirizzo del \inglese{front controller} presente sul server con la stringa
\begin{verbatim}
    operation=login&username=" + encodeURIComponent(data.username) +
                  "&password=" + encodeURIComponent(data.password)
\end{verbatim}
quindi invocare il già visto metodo \verb+testCredentials+ con il testo della risposta. Il valore da restituire al chiamante è proprio la stringa di \inglese{query} che è stata trasmessa al server.

\item{\method{+ onLogin(user: Object): void}}\\
Questo gestore di evento incapsula il comportamento del \inglese{presenter} nel momento in cui è sollevato un evento avente come identificativo la stringa \verb'login'.

Il metodo ha il compito di memorizzare all'interno del client i dati di autenticazione (in particolare, il valore numerico ID necessario all'identificazione del client da parte del server e degli altri client), avviare la connessione al server e provocare la visualizzazione della \inglese{home screen} dell'applicativo.

\item{\method{+ onRemoveLoginPanel() : void}}\\
Questo gestore di evento ha il compito di provocare la distruzione del pannello di autenticazione, richiamando il metodo opportuno della vista corrispondente, vale a dire \verb+destroy()+.

\item{\method{+ onShowLoginPanel(): void}}\\
Questo gestore di evento ha il compito di provocare la visualizzazione del pannello di login. Questo dovrà essere ottenuto sollevando un nuovo evento tramite il quale tutti i pannelli presenti sono rimossi (\verb'removeAllPanel') quindi demandando alla vita il compito di costruirsi e caricare il relativo \inglese{template}.

\item{\method{-- hasAnsweredCorrectly(username: String, answer: String): void}}\\
Il metodo ha il compito di inviare la risposta ottenuta tramite l'interfaccia grafica al server per verificare se corrisponde al valore memorizzato nel database del sistema in fase di registrazione iniziale. In particolare, il metodo ha il compito di inviare una richiesta AJAX sincrona all'indirizzo del \inglese{front controller} la cui intestazione è la stringa
\begin{verbatim}
  "operation=answer&username=" + username + "&answer=" + answer
\end{verbatim}

Una volta ricevuta la risposta, poiché questa già corrisponde alla stringa \verb+true+ o \verb+false+ che rappresentano le costanti booleane in formato JSON, sarà sufficiente restituirne la valutazione ottenuta con la consueta funzione \verb+parse+.

\item{\method{-- testCredentials(data: String): void}}\\
Il metodo prende come parametro una stringa JSON ricevuta dal server che, in caso di dati di autenticazione andata a buon fine, corrisponde a una forma serializzata dell'oggetto che rappresenta l'utente autenticato ed è invece uguale alla stringa \verb+null+ in caso contrario.

Il metodo ha il compito di deserializzare i dati dell'utente tramite la funzione di libreria \verb+JSON.parse(data)+ e, nel caso in cui questa sia valutata positivamente -- vale a dire non è né \verb+null+ né \verb+undefined+ -- sollevare un evento avente come nome la stringa \verb'login' e come proprietà \verb+user+ un riferimento al nuovo utente creato.

In caso di autenticazione non andata a buon fine, invece, dovrà essere aggiornata in modo opportuno la vista per segnalare la condizione di errore.

\end{description}


\classsection{MainPresenter}

\subsubsection*{Funzione}
\inglese{Presenter} incaricato di gestire il pannello principale, ossia quello che occupa la regione centrale.

\subsubsection*{Relazioni d'uso}
Nessuna relazione d'uso evidenziata.

\subsubsection*{Classi estese ed interfacce implementate}
Nessuna relazione evidenziata.

\subsubsection*{Attributi}
\begin{description}
\item{\memberdata{-- thisPanel: HTMLElement}}\\
  Riferimento all'elemento controllato dal \inglese{presenter} vale a dire un \verb'<div>' avente come identificativo la stringa \verb'MainPanel'.
\end{description}

\subsubsection*{Metodi}
\begin{description}

  \item{\method{+ onRemoveMainPanel(): void}}\\
  Questo gestore di eventi è richiamato nel momento in cui è sollevato un evento con nome identificativo \texttt{removeMainPanel} e ha il compito di richiamare il metodo \method{destroy()} della vista corrispondente.
  
  \item{\method{+ onShowMainPanel(): void}}\\
  Questo gestore di eventi è richiamato nel momento in cui è sollevato un evento avente come nome identificativo la stringa \texttt{showMainPanel} e ha il compito di provocare il caricamento del \inglese{template} del pannello e, in seguito, di demandare al metodo \method{display()} della vista la sua inizializzazione.

\end{description}


\classsection{MessagePresenter}

\subsubsection*{Funzione}
\inglese{Presenter} incaricato di gestire i messaggi in segreteria che permette di scaricarli dal server, eliminarli e impostarne lo stato.

\subsubsection*{Relazioni d'uso}
Nessuna relazione d'uso evidenziata.

\subsubsection*{Classi estese ed interfacce implementate}
Nessuna relazione evidenziata.

\subsubsection*{Attributi}
\begin{description}
 \item{\memberdata{-- messages: Array}}\\
 Array destinato a contenere la lista dei messaggi della segreteria il cui destinatario è l'utente del client.

 \item{\memberdata{-- thisPanel: HTMLElement}}\\
  Riferimento all'elemento HTML, da intendersi come nodo DOM, che rappresenta il pannello della segreteria ed è concretamente un \verb'<div>' con identificativo uguale a \verb'MessagePanel'.
\end{description}

\subsubsection*{Metodi}
\begin{description}

  \item{\method{+ deleteMessage(message: Object) Boolean}}\\
  Elimina un messaggio dalla segreteria inviando all'indirizzo del \inglese{front controller} del server una richiesta AJAX sincrona tramite un'istanza di \verb'XMLHttpRequest' che contiene la \inglese{query string}
  \begin{verbatim}
    "operation=deleteMessage&idMessage=" + message.id
  \end{verbatim}
  
  Poiché il server restituisce le stringhe \verb'"true"' o \verb'"false"' che corrispondono alla rappresentazione in formato JSON delle omonime costanti letterali del linguaggio, sarà sufficiente che il metodo restituisca il risultato dell'espressione \verb'JSON.parse(request.responseText)'. In caso di errore, inoltre, dovrà comparire un messaggio d'errore indicante il motivo.
  
    \item{\method{+ displayList(): void}}\\
  Tale metodo ha il compito di popolare la vista con i messaggi della segreteria telefonica provvedendo in primo luogo a scaricarli dal server mediante il metodo \method{getMessages()} al fine di inizializzare l'attributo \memberdata{messages}, quindi aggiornare la vista in maniera corrispondente invocando per ognuno dei messaggi scaricati in precedenza il metodo \method{addListItem} messo a disposizione dalla vista.
  
    \item{\method{+ onShowMessagePanel(): void}}\\
  Questo gestore di eventi viene attivato al verificarsi dell'evento \verb'showMessagePanel' e determina il caricamento del \inglese{template} che corrisponde al pannello della segreteria telefonica, demandando in un secondo momento alla vista il compito di inizializzare il comportamento del pannello chiamando il metodo \method{display()} messo a disposizione dalla vista.
  
    \item{\method{+ setStatus(message: Object, valueToSet: Boolean): Boolean}}\\
  Imposta lo stato di un messaggio come ``nuovo'' oppure ``già letto'' a seconda del parametro ricevuto in input. Tale metodo ha il duplice compito di modificare lo stato del messaggio ricevuto come parametro in base al valore di \verb'valueToSet' e quindi inviare una richiesta AJAX sincrona al server che contiene la \inglese{query string}
  \begin{verbatim}
    "operation=updateMessage&idMessage=" + message.id +
               "&read=" + valueToSet
  \end{verbatim}
  
  Poiché il server restituisce le stringhe \verb'"true"' o \verb'"false"' che corrispondono alla rappresentazione in formato JSON delle omonime costanti letterali del linguaggio, sarà sufficiente che il metodo restituisca il risultato dell'espressione \verb'JSON.parse(request.responseText)'.
  
    \item{\method{-- getMessages(): Array}}\\
  Ottiene i messaggi di segreteria salvati nel server creando una nuova \verb'XMLHttpRequest' da inviare all'indirizzo del \inglese{front controller} del server per realizzare una richiesta AJAX sincrona nella cui \inglese{query string} dovrà obbligatoriamente essere presente il parametro \verb'operation', con una sintassi del tipo:
  \begin{verbatim}
    request.send("operation=getMessages");
  \end{verbatim}
  
  Dal momento che il server risponde con la versione serializzata in JSON della lista di messaggi desiderata, il metodo dovrà limitarsi a restituire il risultato dell'invocazione di \verb'JSON.parse(request.responseText)'.
  
\end{description}


% Qui dovrà andarci PRESENTERMEDIATOR


\classsection{RegisterPresenter}

\subsubsection*{Funzione}
\inglese{Presenter} incaricato di gestire il \inglese{form} di registrazione al sistema.

\subsubsection*{Relazioni d'uso}

Nessuna relazione d'uso evidenziata.

\subsubsection*{Classi estese ed interfacce implementate}

Nessuna relazione evidenziata.

\subsubsection*{Attributi}
\begin{description}
  \item{\memberdata{-- thisPanel: HTMLElement}}\\
  Elemento HTML (da intendersi come nodo DOM) controllato da questo Presenter, che corrisponde a un elemento \verb'<div>' avente identificativo uguale a \verb'RegisterPanel'.
\end{description}

\subsubsection*{Metodi}
\begin{description}

	\item{\method{+ getAnswer(): String}}\\
	Metodo che estrae dal \inglese{form} la risposta alla domanda segreta associata al nuovo utente e la restituisce se specificata, mentre in caso contrario solleva un'eccezione avente come messaggio esattamente la stringa ``risposta alla domanda segreta non specificata''.
	
	\item{\method{+ getName(): String}}\\
	Metodo che estrae dalla \inglese{form} il valore del nome del nuovo utente e lo restituisce.
	
		\item{\method{+ getPassword(): String}}\\
	Metodo che estrae dalla \inglese{form} la \textit{password} associata al nuovo utente e la restituisce se presente, mentre in caso contrario solleva un'eccezione avente come messaggio esattamente la stringa ``password non specificata''.
	
		\item{\method{+ getPicturePath(): String}}\\
	Metodo che estrae dal \inglese{form} il percorso dell'immagine del nuovo profilo utente e lo restituisce.
	
		\item{\method{+ getQuestion(): String}}\\
	Metodo che estrae dalla \inglese{form} la domanda segreta associata al nuovo utente e la restituisce se presente, mentre in caso contrario solleva un'eccezione avente come messaggio esattamente la stringa ``domanda segreta non specificata''.

	\item{\method{+ getSurname(): String}}\\
	Metodo che estrae dalla \inglese{form} il valore del cognome del nuovo utente e lo restituisce.
	

	\item{\method{+ getUsername(): String}}\\
	Metodo che estrae dalla \inglese{form} lo \textit{username} del nuovo utente (per \textit{username} si intende l'indirizzo email dell'utente). Il metodo dovrà verificare che la mail restituita corrisponda all'espressione regolare:
\begin{verbatim}
	  ^[A-Z0-9._%+-]+@[A-Z0-9.-]+\.[A-Z]{2,4}$
\end{verbatim}
	
Se la mail è valida il metodo termina ritornandone il valore, mentre nel caso in cui non sia un indirizzo ben formato o non sia stato specificato il metodo solleva delle eccezioni aventi come messaggio, rispettivamente, ``indirizzo email non valido'' e ``indirizzo email non specificato''.

	\item{\method{+ onRemoveRegistrationPanel(): void}}\\
	Questo gestore di evento ha il compito di provocare la distruzione del pannello di registrazione, richiamando il metodo opportuno della vista corrispondente, vale a dire \verb+destroy()+.
	
	\item{\method{+ onShowRegistrationPanel(): void}}\\
	Questo gestore di evento ha il compito di provocare la visualizzazione del pannello di registrazione. Ciò dovrà essere ottenuto sollevando un nuovo evento tramite il quale tutti i pannelli presenti sono rimossi (\verb'removeAllPanel') quindi demandando alla vista il compito di costruirsi e caricare il relativo \inglese{template}.
	
	\item{\method{+ register(userData: Object): void}}\\
	Metodo che invia i dati ricevuti al \inglese{front controller} del server per la creazione di un nuovo account utente. Tale metodo inizia creando un istanza di \texttt{XMLHttpRequest} sincrona da inviare all'indirizzo del server. Il metodo dovrà estrarre le proprietà obbligatorie \verb'username', \verb'password', \verb'question', \verb'answer' e quelle facoltative \verb'name', \verb'surname' e \verb'picturePath'.
	
	Dopo averle opportunamente codificate per per la trasmissione con la funzione di libreria \verb'encodeURIComponent' queste ultime dovranno essere inserite in un oggetto \verb'FormData' da utilizzare per inviare la richiesta al server. Una volta ricevuta la riposta dal server, il metodo dovrà effettuare la deserializzazione del testo ad essa associato e, se questo contiene un oggetto utente legale, sollevare un evento con nome identificativo \verb'login'.

\end{description}


\classsection{SearchResultPresenter}

\subsubsection*{Funzione}
\inglese{Presenter} incaricato di gestire il pannello che visualizza i risultati di ricerca.

\subsubsection*{Relazioni d'uso}
Nessuna relazione d'uso evidenziata.

\subsubsection*{Classi estese ed interfacce implementate}
Nessuna relazione evidenziata.

\subsubsection*{Attributi}
\begin{description}
\item{\memberdata{-- thisPanel: HTMLElement}}\\
Riferimento al pannello controllato da questo \inglese{presenter}, vale a dire un elemento \verb'<div>' avente come identificativo la stringa \verb'SearchResultPanel'.
\end{description}

\subsubsection*{Metodi}
\begin{description}
  
    \item{\method{+ handleResult(contacts: Object): void}}\\
  Il metodo riceve in input in array associativo indicizzato in base agli identificativi dei contatti che in ogni posizione contiene l'oggetto contatto corrispondente alla chiave.
  
  Il compito del metodo è azzerare il contenuto della lista di utenti visualizzati (se era non vuota), quindi richiamare per ognuno dei contatti contenuti nell'array richiamare il metodo \method{addListItem} fornito dalla vista e passando il contatto come parametro.
  
  \item{\method{+ onShowSearchResultPanel(): void}}\\
  Questo gestore di eventi ha il compito di provocare la visualizzazione del pannello di ricerca caricando il \inglese{template} corrispondente e demandando alla \classname{SearchResultView} di inizializzarsi tramite il metodo \method{display()} preposto a tale scopo.
  
  \item{\method{+ sendSearchRequest(pattern: String): Object}}\\
  Il metodo ha il compito di avviare una ricerca fra gli utenti registrati al sistema a partire dalla stringa \verb'param' come criterio di ricerca (che può corrispondere a parte del nome, del cognome o del nome utente).
  
  A tale scopo il \inglese{presenter} dovrà contattare il server inviando all'indirizzo del\inglese{front controller} una richiesta AJAX sincrona con il metodo POST e avente come \inglese{query string} esattamente la stringa
  \begin{verbatim}
    "operation=search&param=" + pattern
  \end{verbatim}
  
  Dal momento che il server risponde inviando in forma serializzata l'array di contatti desiderato, sarà sufficiente che il metodo restituisca la forma deserializzata del contenuto della risposta, con un'istruzione del tipo:
  \begin{verbatim}
    return JSON.parse(request.responseText);
  \end{verbatim}
  
\end{description}


\classsection{ToolsPresenter}

\subsubsection*{Funzione}
\inglese{Presenter} che ha il compito di gestire il pannello degli strumenti.

\subsubsection*{Relazioni d'uso}
Nessuna relazione d'uso evidenziata.

\subsubsection*{Classi estese ed interfacce implementate}
Nessuna relazione evidenziata.

\subsubsection*{Attributi}
\begin{description}
\item{\memberdata{-- thisPanel: HTMLElement}}\\
Riferimento all'elemento HTML, da intendersi come nodo DOM, che corrisponde al pannello degli strumenti e che nello specifico è rappresentato da un elemento \verb'<div>' avente come identificativo la stringa \verb'ToolsPanel'.
\end{description}

\subsubsection*{Metodi}
\begin{description}

  \item{\method{+ onLogout(): void}}\\
  Questo gestore di eventi viene richiamato al verificarsi dell'evento \verb'logout' e ha il compito di effettuare l'uscita dal sistema. In particolare, nel corpo del metodo dovrà essere creata una nuova \verb'XMLHttpRequest' da inviare all'indirizzo del \inglese{front controller} presente sul server, nella cui \inglese{query string} sarà presente il parametro \verb'operation=logout' per permettere il corretto smistamento della richiesta.
  
  Infine, dal momento che il server risponde con una stringa \verb'true' o \verb'false', il metodo ha il compito di effettuare la deserializzazione con \verb'JSON.parse(request.responseText)' e nel caso in cui il risultato fosse valutato negativamente, provocare la visualizzazione di un generico messaggio d'errore.
  
  \item{\method{+ onRemoveToolsPanel(): void}}\\
  Questo gestore di eventi viene associato come risposta del \inglese{presenter} nel momento in cui viene sollevato un event con nome identificativo uguale a \verb'removeToolsPanel'. Il metodo dovrà provocare l'eliminazione dall'interfaccia grafica del pannello degli strumenti, demandando alla vista la propria rimozione dall'albero DOM tramite il metodo \method{destroy()} messo a disposizione da quest'ultima.
  
  \item{\method{+ onShowReturnToCommunicationPanelButton(): void}}\\
  Questo gestore di eventi è invocato nel momento in cui si attiva una comunicazione di qualsiasi natura (audio, audio/video o testuale) e provoca l'aggiunta di un nuovo pulsante nel pannello grafico tramite il metodo \method{addCommunicationFunction()} fornito dalla vista. Lo scopo di tale pulsante è permettere di tornare al pannello della comunicazione in un secondo momento, dopo che si è passati alla visualizzazione di un altro pannello.
  
  \item{\method{+ onShowToolsPanel(): void}}\\
  Questo gestore di eventi è associato all'evento con identificativo \texttt{showToolsPanel} e ha il compito di provocare la visualizzazione del pannello degli strumenti sulla parte sinistra dell'interfaccia grafica. Il metodo dovrà avere il compito di caricare il \inglese{template} della vista dal server. L'inizializzazione del pannello sarà quindi demandata in un secondo momento alla vista, tramite il metodo \method{initialize} fornito da quest'ultima.

\end{description}


%***********************************************************************
%*                QUESTO È IL CONFINE FRA IL BENE E IL MALE            *
%***********************************************************************
% TODO IBI ERIT FLETUS ET STRIDOR DENTIUM
\classsection{AddressBookPanelPresenter}

\subsubsection*{Funzione}
Presenter incaricato di gestire il pannello della rubrica, contiene le funzioni associate ai \inglese{widget} grafici della vista relativi a quest'ultima e ha la responsabilità di aggiornare la vista sulla base dei dati ricevuti dal \inglese{server}.

\subsubsection*{Relazioni d'uso}
Nessuna relazione d'uso evidenziata.

\subsubsection*{Classi estese ed interfacce implementate}
Nessuna relazione evidenziata.

\subsubsection*{Attributi}
\begin{description}

  \item{\memberdata{-- element: DOM\_node}}\\
  Attributo che definisce il contenuto del nodo DOM inerente alla rappresentazione della rubrica. Nello specifico tale nodo corrisponde ad un \texttt{div} con id \texttt{AddressBookPanel}.
  \item{\memberdata{-- urlServlet: String}}\\
  Attributo che definisce l'URL dal quale è richiamabile la \textit{servlet} usata per interagire con l'apparato server del sistema.
  \item{\memberdata{-- suffix: String}}\\
Suffisso da concatenare all'url per richiamare le \inglese{servlet} appropriate.  
  \item{\memberdata{-- contacts: Array()}}\\
  Attributo che definisce la lista degli utenti presenti nella rubrica dell'utente che ha effettuato l'autenticazione.
   \item{\memberdata{-- groups: Array()}}\\
  Attributo che definisce la lista dei gruppi presenti nella rubrica dell'utente che ha effettuato l'autenticazione. Pero ogni gruppo sono  elencati i contatti contenuti in esso.
   \item{\memberdata{-- operations: String[]}}\\ 
   \inglese{Array} di stringhe usato per memorizzare i nomi delle istruzioni eseguibili, essi rappresentano una parte del nome della \inglese{servlet} incaricata a risolvere l'operazione descritta (e.g. GetContacts rappresenta l'operazione svolta dalla \inglese{servlet} avente nome: \memberdata{urlServlet} concatenato a ``GetContactsServlet''). I valori con i quali si dovrà caricare \memberdata{operations}) sono:
   \begin{itemize}
   	\item GetContacts;
    \item DoAddContact;
    \item DoDeleteContact;
    \item DoInsertInGroup;
    \item DoRemoveFromGroup;
    \item DoCreateGroup;
    \item DoDeleteGroup;
    \item GetGroups;
    \item DoBlock;
    \item DoUnblock;
    \item DoSearch;
   \end{itemize}

\end{description}

\subsubsection*{Metodi}
\begin{description}	
	
	\item{\method{+ inizialize(): void}}\\
	Metodo per inizializzare  \texttt{AddressBookPanel} e popolarlo con i contatti della rubrica, modifica il DOM creando tre div: divSearch, divSort e divList. Quindi sequenzialmente ne crea il contenuto:
		\begin{itemize}
			\item divSearch: definisce un campo per l'\inglese{input} testuale e un bottone per ricercare nella lista un utente avente avete tra le parole chiave, una con valore uguale a quello inserito nel campo d'\inglese{input}.
			\item divSort: crea delle \inglese{select} per specificare le tipologie di ordinamento attuabili sulla lista degli utenti.
			\item divList: \textit{div} che viene popolato con i nomi degli utenti presenti nella rubrica dell'utente che l'ha richiamata.
		\end{itemize}

Quindi dopo aver creato questi elementi il metodo procede inserendo nel DOM il ``sotto-albero'' cosi creato.

	\item{\method{- getAddressBookContacts(): void}}\\
Metodo che recupera i contatti e i gruppi della propria rubrica dal server usando la tecnologia AJAX. Il metodo crea una variabile contenente un istanza \texttt{XMLHttpRequest()}. Quindi con metodo \inglese{POST} inoltra la richiesta alla \inglese{servlet} il cui indirizzo è generato dalla concatenazione di \memberdata{urlServlet} con \memberdata{operation[1]}. Elabora quindi i dati ottenuti mediante il metodo \texttt{JSON.parse()} e li salva negli \textit{array} \memberdata{contacts} e \memberdata{groups}.
	
	\item{\method{+ setup(): void}}\\
	Inserisce i contatti estratti dal \textit{server} all'interno della lista \texttt{AddressBookList} all'interno di \texttt{AddressBookPanel}.
	
	\item{\method{+ displayContactList(contact: String[]): void}}\\
	Metodo usato per visualizzare la lista dei contatti utente, esso prende il nodo DOM (già esistente) e destinato a contenere la lista dei contatti, popolandolo richiamando il metodo \method{addListItem()}.
	
	\item{\method{-- addListItem(list: DOM\_node, contact: Object(contact): void}}\\

Metodo richiamato da \method{displayContactList()} per visualizzare i contatti presenti della rubrica, crea le variabili contenenti i dati del contatto da attribuire al \textit{tag} di \textit{markup} \texttt{li}. Tali variabili sono quelle tipiche di un oggetto Object(contact) (nome, cognome, email, status e immagine).


\item{\method{-- addOptionToSelect(select: DOM\_element, value: String, text: String): void}}\\
Metodo utilizzato per aggiungere un'opzione all'elemento \memberdata{select} ricevuto come parametro dal metodo. Ha il compito di creare un \texttt{TextNode} ed appenderlo alla \memberdata{select}.

	
	\item{\method{+ addContact(contact: Object(contact)): boolean}}\\
	Metodo per l'aggiunta di un contatto alla rubrica, riceve come parametro un oggetto di tipo contatto da inserire, appunto, nella propria rubrica. 
Per assicurarsi di non introdurre ridondanza ed aggiungere un contatto già presente il metodo esegue due controlli: uno lato \inglese{client} e uno lato \textit{server}. 
Il controllo lato \inglese{client}, si scorre la lista \memberdata{contact} alla ricerca di un \texttt{contact} che cerca un \inglese{match} con il nuovo contatto. Nel caso sia già presente il metodo termina ritornando \inglese{false}. Per eseguire il controllo lato \inglese{server} il metodo contatta il \inglese{server} mediante la \inglese{servlet} \texttt{AddressBookDoAddContatctServlet}, questa verificherà se l'utente avente \texttt{id} \textit{contact} è già presente nella rubrica, nel caso non lo sia si occuperà di creare la nuova istanza di \texttt{AddressBookEntry} e la registrerà nel database, in caso contrario non esegue operazioni. 
La chiamata a \inglese{servlet} ritorna un valore booleano che sarà \textit{true} se l'inserimento ha avuto successo, \textit{false} altrimenti. 
Se l'inserimento ha avuto successo il flusso principale prosegue eseguendo l'aggiornamento della rubrica in locale tramite la chiamata a metodo \method{setup()}. Il metodo termina infine ritornando \textit{true} se è stato aggiunto l'utente, altrimenti solleva un'eccezione.

	\item{\method{+ removeContact(contact: Obj): boolean}}\\
Metodo il cui funzionamento è analogo a quello precedentemente definito, si occupa di eliminare un contatto conoscendone l'\texttt{id} passato per parametro d'ingresso. Il flusso principale esegue due controlli per verificare la presenza dell'utente nella rubrica (analogamente a quanto visto nel caso precedente). 
Dopo tale verifica, se il contatto è realmente presente nella rubrica, la \inglese{servlet} \texttt{AddressBookDoRemoveContactServlet} avrà eseguito la cancellazione come conseguenza al match effettuato in fase di verifica, nel caso il metodo richiama \method{setup()} per rieseguire l'aggiornamento della rubrica. Il metodo ritorna infine \textit{true} se effettivamente il contatto era presente ed è stato eliminato, altrimenti solleva un'eccezione.

	\item{\method{+ addGroup(name: String): boolean}}\\
	Metodo per l'aggiunta di un nuovo gruppo nella rubrica, permette l'aggiunta di gruppi con nomi duplicati (ossia nomi di gruppi già presenti nella rubrica). Il flusso principale richiama la \inglese{servlet} \texttt{AddressBookDoCreateGroupServlet} con il nome del gruppo da creare. Successivamente, se la \inglese{servlet} ha dato \method{setup()} esito positivo, viene richiamata il metodo \inglese{setup} che aggiorna la rubrica locale, altrimenti solleva un'eccezione.
	
	\item{\method{+ deleteGroup(idGroup: int): boolean}}\\
	Metodo per l'eliminazione di un nuovo gruppo presente nella rubrica, riceve come parametro d'ingresso l'\texttt{id} del gruppo da eliminare. Il flusso principale richiama la \inglese{servlet} \texttt{AddressBookDoRemoveGroupServlet} passandogli l'\texttt{id} del gruppo candidato all'eliminazione. Quindi, se la \inglese{servlet} notifica l'effettiva esistenza di un gruppo con tale \texttt{id}, essa si occupa dell'eliminazione e il flusso principale prosegue richiamando il metodo \inglese{setup} se la eliminazione ha avuto buon fine, altrimenti solleva un'eccezione.
	
	\item{\method{+ addContactInGroup(contact: Object(contact), group: Object(contact)): boolean}}\\
	Metodo per l'aggiunta di un contatto in un gruppo ben definito, riceve come parametro d'ingresso il contatto e il gruppo.
Il metodo richiama da prima \method{contactExistInGroup}, successivamente il flusso principale richiama la \inglese{servlet} \texttt{AddressBookDoInsertInGroupServlet} passandogli i parametri d'ingresso, quindi se la \inglese{servlet} notifica l'effettiva esistenza di un gruppo con il nome ricevuto, essa si occupa di aggiungervi il contatto e il flusso principale ritorna al metodo \method{setup()} per aggiornare la rubrica in locale, ritorna infine un \inglese{feedback} booleano a \textit{true} per dare notifica se l'operazione ha avuto buon fine, altrimenti solleva un'eccezione.
	
	\item{\method{+ deleteContactFromGroup(contact: Object(contact), group: Object(group)): boolean}}\\
	Metodo per l'eliminazione di un contatto presente in un gruppo, riceve come parametro d'ingresso il gruppo su cui effettuare la ricerca del contatto da eliminare e il contatto. 
Il metodo richiama da prima \method{contactExistInGroup}, poi Il flusso principale richiama la \inglese{servlet} \texttt{AddressBookDoRemoveFromGroupServlet} passandogli i parametri d'ingresso, la \inglese{servlet} notifica l'effettiva esistenza di un gruppo con il nome ricevuto che contiene il contatto candidato all'eliminazione, essa si occupa dell'eliminazione e il flusso principale ritorna al metodo \method{setup()} per aggiornare la rubrica in locale, ritorna infine un \inglese{feedback} booleano per dare notifica sull'esito dell'operazione: \textit{true}\textit{true} se è stato creato il nuovo gruppo, altrimenti solleva un'eccezione.
	
	\item{\method{- contactExistInGroup(contact: Object(contact), group: Object(group)): boolean}}\\
Metodo che verifica l'esistenza di un contato in un determinato gruppo, restituisce un \inglese{feedback} booleano per confermare o meno l'esistenza dell'utente ricercato nel gruppo designato.

\item{\method{+ applyFilterByString(param: String): array()}}\\
Metodo che data una stringa come parametro d'ingresso filtra la lista di contatti presente in locale. Non richiede interazione con il server, ovvero non vengono richiamate \inglese{servlet}.

In un primo step il programmatore dovrà creare un \inglese{array} atto a contenere l'elenco dei contatti da restituire, quindi il metodo procede con la definizione di un espressione regolare basata sul parametro ricevuto. Per tale operazione la sintassi obbligatoria è:\\

\verb|var pattern = new RegExp(param)|\\

Quindi viene avviato un ciclo \textit{for} che analizza ogni contatto presente in \memberdata{contacts}. Successivamente si effettua un filtraggio sui campi \textit{name, surname, email} del contatto attualmente in esame. Per fare tali operazioni si dovrà usare:\\

\verb|pattern.test(\*elemento su cui ricercare param*\)|\\

Se tale filtraggio da esito positivo, il contatto viene aggiunto all'\textit{array} dei contatti da restituire. Il metodo termina quando, dopo aver eseguito il ciclo \textit{for}, il metodo ritorna l'\textit{array} sopracitato.

\item{\method{+ applyFilterByGroup(idGroup: int): array()}}\\
Metodo che dato un intero come parametro d'ingresso filtra la lista di contatti presente in locale, restituendo un elenco dei contatti appartenenti al gruppo avente id uguale a idGroup, inoltre non deve interagire con il \inglese{server}, ovvero non vengono richiamate \inglese{servlet}. 

Come per \method{applyFilterByString()} anche in questo caso il metodo deve creare un \textit{array} da popolare con i contatti appartenente al gruppo idGroup. Quindi viene eseguito un ciclo \inglese{for} che ispezione tutti i contatti presenti in \memberdata{contacts} e per ogni ciclo il metodo controlla i gruppi in cui compare il contatto i-esimo (tale azione si attua per mezzo di un secondo ciclo innestato), se si verifica una corrispondenza tra l'id del gruppo in esame con l'identificativo idGroup, si aggiunge il contatto all'\textit{array} creato in precedenza. Il metodo termina restituendo l'\textit{array} sopra-citato.

\item{\method{+ showFilter(filteredContacts: Array()): void}}\\
Metodo usato per l'eliminare il contenuto preesistente nella rubrica e visualizza il nuovo contenuto filtrato. Il metodo riceve come parametro un \textit{array} che rappresenta la lista filtrata dei contatti da visualizzare. Per eseguire tali operazioni viene avviato un ciclo \inglese{for} per ogni contatto presente in \memberdata{contacts}. All'interno di tale ciclo deve essere eseguita la seguente istruzione:\\

\verb|addListItem(/*nome della variabile contenente il nodo dom AddressBookList*/,|\\
\verb|filtredContacts[contact])|\\

\item{\method{+ contactAlredyPresent(contact: Object(contact)): boolean}}\\
Metodo che dato un contatto verifica se appartiene alla rubrica. Per compiere tale azione viene da prima riscaricata la rubrica, mediante una chiamata a metodo \method{getAddressBookContacts()}. Quindi si usa un ciclo \inglese{for} per scorrere l'intera lista precedentemente ottenuta e ad ogni ciclo si controlla se il contatto i-esimo ha id uguale a quello del contatto ricevuto come parametro d'ingresso (\texttt{contact.id}). Se tale confronto ha esito positivo, il metodo termina ritornando \inglese{true}, altrimenti il metodo termina al termine del ciclo \inglese{for} ritornando \inglese{false}.

\item{\method{+ blockUser(contact: Object(contact)): boolean}}\\
Metodo utilizzato per bloccare (all'interno della rubrica utente) il contatto contact passato come parametro d'ingresso. Il primo step consiste nell'eseguire il controllo sul parametro ricevuto. Per farlo si osserva il valore ritornato da una chiamata a metodo \method{constactAlradyPresent()} passando come parametro d'ingresso contact. Se il valore di ritorno è \inglese{false} il metodo termina con l'eccezione ``Contatto non presente nella rubrica''. Altrimenti il controllo passa allo step successivo, in cui si verifica se l'attributo blocked di contact è impostato a \inglese{true}. Nel caso il metodo termina con l'eccezione ``Contatto già bloccato''. Se il flusso principale ha superato i test precedenti, allora viene creata un istanza di request \texttt{XMLHttpRequest()} usata per contattare la \inglese{servlet} avente nome:\\

\verb|urlServlet + operations[8] + suffix|\\

Per tale richiesta deve essere associata, all'evento \texttt{onreadystatechange}, una funzione che salva in un array il contenuto di una chiamata:\\

\verb|JSON.parse(request.responseText)|\\

Il flusso principale prosegue aprendo la request cosi definita (uso del metodo \texttt{open("POST", urlServlet + operations[8] + suffix, "true")}) e inviando il contatto mediante una chiamata \texttt{request.send("contactId=" + contact.id)}.

Il metodo termina eseguendo un controllo. Se il risultato result della chiamata a \inglese{servlet} ha valore \inglese{true} allora il metodo termina aggiornando la rubrica mediante una chiamata \method{setup()} e ritornando \inglese{true}. In caso contrario il metodo termina con un eccezione.

\item{\method{+ unlockUser(contact: Object(contact)): boolean}}\\
Metodo usato per sbloccare un utente all'interno di una rubrica. Il metodo ha un funzionamento analogo al precedente, nello specifico vengono attuati gli stessi controlli prima tramite \method{contactAlradyPresent()} e poi tramite la verifica del contenuto del di contact.block. Successivamente si crea una XMLHttpRequest che si appoggia alla \inglese{servlet} definita dalla stringa:\\

\verb|urlServlet + operations[9] + suffix|\\

Tale server ha il compito di sbloccare il contatto \textit{contact}. Quindi il metodo termina eseguendo un controllo sul tipo di risultato. Se l'operazione precedente è avvenuta con successo allora si procede aggiornando la rubrica (chiamata a metodo \method{setup()}), altrimenti il metodo termina con un eccezione.

\item{\method{+ getGroupsWhereContactsIs(contact: Object(contact)): Object(group)}}\\
Tale metodo ha il compito, dato un utente, di ritornare una lista di gruppi in cui esso è presente. Riceve come unico parametro il contatto da ricercare e successivamente vengono ciclati tutti i gruppi e per ogni gruppo si controlla se è inserito il contatto, in caso di match viene aggiunto il gruppo alla lista di gruppi da ritornare.

\end{description}

% TODO IBI ERIT FLETUS ET STRIDOR DENTIUM
\classsection{CommunicationPanelPresenter}

\subsubsection*{Funzione}
Questo presenter ha il compito di gestire tutte le comunicazioni che possono avvenire tra persone, quindi sia di natura testuale che di tipo audio - audio/video.

\subsubsection*{Relazioni d'uso}
Nessuna relazione d'uso evidenziata.

\subsubsection*{Classi estese ed interfacce implementate}

Nessuna relazione evidenziata.

\subsubsection*{Attributi}
\begin{description}
\item{\memberdata{-- chatElements: HTMLDivElements[]}}\\
  array associativo contenente tutte le chat aperte in un dato momento.
\end{description}

\subsubsection*{Metodi}
\begin{description}
\item{\method{- createLabel(user: Object(contact)): String}}\\
	Metodo che dato un utente (parametro del metodo), restituisce una stringa identificativa per il contatto in questione. Più precisamente, se il contatto contiene sia il nome che il cognome, la stringa restituita sarà composta da questi due elementi, altrimenti in base ai casi solo il nome o il cognome. Nell'ulteriore caso non siano presenti ne il nome ne il cognome, restituisce l'email, campo dati obbligatorio in fase di registrazione e quindi sempre presente.

\item{\method{- createChatItem(user: Object(contact)): DOM\_node}}\\
	Tale metodo restituisce un DOM\_element di tipo <li> contenente l'identificativo dell'utente con cui si sta comunicando (parametro del metodo). All'oggetto restituito viene assegnato un ``id'' corrispondente all'``id'' dell'utente con cui si vuole dialogare. All'evento \textit{onClick} viene collegato l'invocazione del metodo \method{displayChat(user)} che fa visualizzare la chat, viene inoltre creato il bottone per chiudere la chat che con l'evento ``onclick'' richiama il metodo \method{removeChat(user)}.

\item{\method{- createChatElement(user: Object(contact)): DOM\_node}}\\
	Tale metodo restituisce un DOM\_element contenente un \textit{div} che rappresenta l'area di chat. A tale elemento viene assegnato un ``id'' impostato a ``divContainerChat''. Deve essere creato anche una \textit{form} assegnando l'``id'' dell'utente come valore all'attributo ``id'' a tale \textit{form} va aggiunto una \textit{textarea} con \textit{id} impostato a \textit{chatText} dove andrà visualizzata la cronologia della chat; una casella di input con \textit{id} uguale \textit{input}; un bottone per inviare il messaggio assegnandogli all'evento \textit{onClick} l'invocazione del metodo \method{send(user, text)} del \textit{CommunicationPanelPresenter}. Infine viene aggiunta la \textit{form} all'elemento \textit{div} e ritornato.
	
\item{\method{+ displayChat(user: Object(contact)): void}}\\
Provoca la visualizzazione della chat con l'utente (ricevuto da parametro) nel \textit{div} contenitore pertinente.

\item{\method{+ addChat(user: Object(contact)): void}}\\
Provoca l'aggiunta di una chat, con l'utente ricevuto come parametro, andando ad aggiornare il campo privato \textit{chatElements} e visualizzandola nell'elemento ulOpenChat.
	
\item{\method{+ removeChat(user: Object(contact)): void}}\\
Comportamento simile al metodo scritto precedentemente. Viene chiusa una chat aperta con l'utente ricevuto da parametro andando ad eliminare l'elemento corrispondente da \textit{chatElements} e rimuovendo anche il DOM\_element dall'ulOpenChat.

\item{\method{+ appendToChat(user: Object(contact), text: String): void}}\\
Il metodo ha il compito di aggiunge una stringa \textit{text} all'interno dell'area di testo che è associata alla chat con l'utente \textit{user} (parametro del metodo).

\item{\method{+ createPanel(): DOM\_node}}\\
Con l'invocazione di tale metodo, viene inizializzato il pannello costruendo i \inglese{widget} grafici interni e lo restituisce in modo che possa essere inserito all'interno del pannello principale. Più dettagliatamente, deve costruire:
\begin{itemize}
\item div relativo alla chiamata assegnandogli come ``id'' CommunicationPanel;
\item div relativo alla gestione della chiamta assegnandogli come ``id'' divCall;
\item div relativo alla chat testuale assegnandogli come ``id'' divChat;
\item elemento video per la visualizzazione del proprio stream ricavato dalla videocamera con ``id'' myVideo;
\item elemento video per la visualizzazione dello \textit{stream} dell'altro utente con \textit{id} otherVideo;
\item div dove inserire la visualizzazione delle statistiche della chiamata tra cui:
\begin{itemize}
\item elemento ``span'' per byte Ricevuti con \textit{id}= statReceived;
\item elemento ``span'' per byte Inviati con \textit{id}= statSend;
\item elemento ``span'' per il tempo con \textit{id}= timerSpan;
\end{itemize}
\item bottone per la terminazione della chiamata che all'evento \textit{onClick} richiama il metodo \method{endCall()} di communicationcenter;
\item elemento div contenente le chat aperte.
\end{itemize}
Tutti gli elementi sopra descritti andranno poi aggiunti al parametro di ritorno \textit{element} tramite la funzione \method{appendChild()} fornita dai DOM\_element.

\item{\method{+ updateTimer(text: String): void}}\\
Il metodo aggiorna il tempo di comunicazione durante una chiamata. Imposta il campo \textit{value} dello \inglese{span} relativo alla visualizzazione del tempo chiamata con il valore del parametro ricevuto in input (text).

\item{\method{+ updateStats(text: String, isReceivedData: boolean): void}}\\
Il metodo controlla tramite il parametro booleano ricevuto in input (isReceivedData) se bisogna aggiornare lo span relativo ai dati inviati oppure lo span relativo ai dati ricevuti e lo imposta con il valore \textit{text} ricevuto per parametro.

\item{\method{+ getMyVideo(): DOM\_element}}\\
Il metodo ritorna l'elemento DOM con \textit{id} uguale ad ``myVideo''.

\item{\method{+ getOtherVideo(): DOM\_element}}\\
Il metodo ritorna l'elemento DOM con \textit{id} uguale ad ``otherVideo''.

\end{description}

% TODO IBI ERIT FLETUS ET STRIDOR DENTIUM
\classsection{ContactPanelPresenter}

\subsubsection*{Funzione}
Presenter incaricato di gestire il pannello che visualizza le informazioni di un singolo contatto.

\subsubsection*{Relazioni d'uso}
Nessuna relazione d'uso evidenziata.

\subsubsection*{Classi estese ed interfacce implementate}

Nessuna relazione evidenziata.

\subsubsection*{Attributi}
\begin{description}
\item{\memberdata{-- element: DOM\_element}}\\
  elemento controllato dal presenter in questione.
\end{description}

\subsubsection*{Metodi}
\begin{description}
\item{\method{- adjustBlockButtonDisplay(contact: Object(contact)): void}}\\
	Il metodo riceve come parametro un contatto (\textit{contact}), successivamente viene controllato se tale utente è bloccato o meno e imposta correttamente la proprietà \textit{style} agli elementi della vista. più dettagliatamente:\\
se il contatto è bloccato:
\begin{itemize}
\item nascondo il bottone per bloccare un utente;
\item mostro il bottone per sbloccare un utente
\end{itemize}
altrimenti\\
\begin{itemize}
\item mostro il bottone per bloccare un utente;
\item nascondo il bottone per sbloccare un utente
\end{itemize}


\item{\method{- buildGroupsDiv(contact Object(contact)): void}}\\
Il metodo, che riceve come parametro un contatto (\textit{contact}, ha il compito di recuperare tutti i gruppi 
 a cui il contatto appartiene e visualizzare nel \textit{div} opportuno una \textit{label} per ogni gruppo trovato indicandone il nome. 

\item{\method{+ createPanel():DOM\_element}}\\
Metodo richiamato quando viene selezionato un contatto dalla rubrica. Inizializza il pannello che mostra le informazioni del contatto selezionato. Più precisamente, verranno creati tanti listItem (<li>) quanti i seguenti elementi:
\begin{itemize}
\item nome;
\item cognome;
\item email;
\item avatar raffigurante l'immagine del contatto;
\item div contenente i gruppi a cui il contatto appartiene;
\item pulsante per la chiamata audio;
\item pulsante per la chiamata video;
\item pulsante per avviare una chat;
\item pulsante per bloccare il contatto;
\item pulsante per sbloccare il contatto.
\end{itemize}
Infine ritorno l'elemento creato.
	
\item{\method{+ display(contact: Object(contact)): void}}\\
Metodo che popola nel corretto modo tutti gli elementi necessari per la visualizzazione di un utente. Prima di tutto quindi si recupereranno tutti gli elementi necessari, ossia:
\begin{itemize}
\item elemento per nome;
\item elemento per cognome;
\item elemento per email;
\item elemento per avatar;
\item pulsante per aggiungere un contatto alla rubrica;
\item pulsante per la chiamata audio;
\item pulsante per la video-chiamata;
\item pulsante per avviare una chat;
\item pulsante per bloccare il contatto;
\item pulsante per sbloccare il contatto.
\end{itemize}
Successivamente si andranno a popolare tali elementi utilizzando le informazioni del contatto passato come parametro (\textit{contact}), infine si imposterà correttamente il comportamento dei bottoni associando le giuste chiamate a metodi del \textit{mediator}, ossia:
\begin{itemize}
\item aggiungi contatto --> onContactAdded(contact.id);
\item blocca contatto --> onBlockContact();
\item sblocca contatto --> onUnlockContact();
\item chiamata audio --> onCall(contact, false);
\item chiamata video --> onCall(contact, true);
\item avvia chat --> onChatStarted(contact);
\end{itemize}

\end{description}

\classsection{PresenterMediator}

\subsubsection*{Funzione}
Presenter incaricato di gestire la collaborazione tra i vari sotto-presenter.

\subsubsection*{Relazioni d'uso}
Nessuna relazione d'uso evidenziata.

\subsubsection*{Classi estese ed interfacce implementate}
Nessuna relazione evidenziata.

\subsubsection*{Attributi}
\begin{description}
\item{\memberdata{-- presenters: Array associativo}}\\
  \textit{array} contenente tutti i sotto-presenter.\\
  Presenter di primo livello:
  \begin{itemize}
  \item login --> LoginPanelPresenter;
  \item register --> RegisterPanelPresenter;
  \item addressbook --> AddressBookPanelPresenter;
  \item tools --> ToolsPanelPresenter;
  \item main --> MainPanelPresenter.
  \end{itemize}
  Presenter di secondo livello:
  \begin{itemize}
  \item accountsettingspp --> AccountSettingsPanelPresenter;
  \item communicationpp --> CommunicationPanelPresenter;
  \item contactpp --> ContactPanelPresenter;
  \item callhistorypp --> CallHistoryPanelPresenter;
  \item messagepp --> MessagePanelPresenter;;
  \item searchresultpp --> SearchResultPanelPresenter;
  \item grouppp --> GroupPanelPresenter.
  \end{itemize}
\end{description}

\subsubsection*{Metodi}
\begin{description}
\item{\method{- getCommunicationPP(): CommunicationPanelPresenter}}
Il metodo restituisce l'istanza dell'attributo di tipo  CommunicationPanelPresenter.

\item{\method{+ buildUI(): void}}\\
Tale metodo ha il compito di inizializzare l'interfaccia grafica delegando ai presenter il compito di disegnare gli elementi principali dell'interfaccia, incaricando i presenter di primo livello di creare e popolare i rispettivi pannelli. In particolare, si nasconderanno i seguenti pannelli:
\begin{itemize}
\item register;
\item login;
\item addressbook;
\item main;
\item tools.
\end{itemize}


\item{\method{+ buildLoginUI(): void}}\\
Visualizza l'interfaccia di autenticazione al sistema, che comprende il \textit{form} di \textit{login}. Viene chiamato quindi il metodo \method{initialize()} del LoginPanelPresenter per la sua costruzione.
	
\item{\method{+ buildRegistrationUI(): void}}\\
Visualizza il form di registrazione al sistema, utilizzato dagli utenti che vogliono creare un nuovo account, viene chiamato quindi il metodo \method{initialize()} del RegisterPanelPresenter per la sua costruzione.

\item{\method{+ onContactSelected(contact: Object(contact)): void}}\\
Metodo invocato nel momento in cui viene selezionato un contatto, tale evento invoca il metodo \method{display} del ContactPanelPresenter prendendo come parametro l'oggetto \textit{contact} ricevuto inizialmente. 


\item{\method{+ onContactAdded(contact: Object(contact)): void}}\\
Funzione di \textit{callback} richiamata dai pulsanti di SearchResultPanel che comunica all'AddressBookPanelPresenter di aggiungere un contatto. Riceve come parametro l'utente che si vuole aggiungere (\textit{contact}), questo parametro viene a suo volta passato alla funzione \texttt{addContact()} di addressBookPanelPresenter. In caso di errore, viene segnalato con un \textit{alert}.

\item{\method{+ onContactRemoved(userID: String): void}}\\
Funzione di \textit{callback} richiamata dai pulsanti di SearchResultPanel che comunica all'AddressBookPanelPresenter di rimuovere un contatto. Il metodo riceve come parametro l'utente che si vuole rimuovere (contact). Questo parametro viene a suo volta passato alla funzione \texttt{removeContact()} di addressBookPanelPresenter. In caso di errore, viene segnalato con un \textit{alert}.

\item{\method{+ onGroupAdded(name: String): void}}\\
Funzione di \textit{callback} richiamata dai pulsanti di SearchResultPanel che comunica all'AddressBookPanelPresenter di aggiungere un gruppo. Il metodo riceve come parametro il nome del nuovo gruppo che si vuole aggiungere (name). Questo parametro viene a suo volta passato alla funzione addGroup() di addressBookPanelPresenter.


\item{\method{+ onGroupRemoved(group: object(group)): void}}\\
Funzione di \textit{callback} richiamata dai pulsanti di SearchResultPanel che comunica all'AddressBookPanelPresenter di rimuovere un gruppo. Il metodo riceve come parametro il gruppo che si vuole rimuovere (\textit{group}). Questo parametro viene a suo volta passato alla funzione \texttt{removeGroup()} di addressBookPanelPresenter.


\item{\method{+ onContactAddeddInGroup(contact: Object(contact), group: Object(group)): void}}\\
Funzione di \textit{callback} richiamata dai pulsanti di SearchResultPanel che comunica all'AddressBookPanelPresenter di aggiungere un contatto ad un gruppo. Il metodo riceve come parametro il contatto che si vuole aggiungere (\textit{contact}) e il gruppo a cui aggiungere il contatto (\textit{group}). Questi parametri vengono a suo volta passati alla funzione \texttt{addContactInGroup()} di addressBookPanelPresenter.


\item{\method{+ onContactRemovedInGroup(contact: Object(contact), group: Object(group)): void}}\\
Funzione di \textit{callback} richiamata dai pulsanti di SearchResultPanel che comunica all'AddressBookPanelPresenter di rimuovere un contatto da un gruppo. Il metodo riceve come parametro il contatto che si vuole rimuovere (\textit{contact}) e il gruppo a cui rimuovere il contatto (\textit{group}). Questi parametri vengono a suo volta passati alla funzione \texttt{removeContactFromGroup()} di addressBookPanelPresenter.


\item{\method{+ onBlockedContact(contact: Object(contact)): void}}\\
Funzione di \textit{callback} che comunica all'AddressBookPanelPresenter di bloccare un contatto. Il metodo riceve come parametro il contatto che si vuole bloccare (\textit{contact}), questo parametro viene a suo volta passato alla funzione \texttt{blockUser()} di addressBookPanelPresenter. In caso di errore, viene notificato con un \textit{alert} riportando l'errore relativo.


\item{\method{+ onUnlockContact(contact: Object(contact)): void}}\\
Funzione di \textit{callback} che comunica all'AddressBookPanelPresenter di sbloccare un contatto. Il metodo riceve come parametro il contatto che si vuole sbloccare (\textit{contact}), questo parametro viene a suo volta passato alla funzione \texttt{unlockUser()} di addressBookPanelPresenter. In caso di errore, viene notificato con un \textit{alert} riportando l'errore relativo.


\item{\method{+ displayMessagePanel(): void}}\\
Provoca la creazione del pannello della segreteria e la sua visualizzazione all'interno del MainPanel come elemento figlio. La costruzione del pannello è affidata al metodo \textit{createPanel} che viene reso disponibile da tutti i presenter di secondo livello, viene quindi creato dal metodo \method{createPanel()} del presenter messagePanelPresenter un elemento di tipo DOM\_node che viene poi passato al presenter \textit{main} come parametro del metodo \method{displayChildPanel}.


\item{\method{+ displayAccountSettingsPanel(): void}}\\
Provoca la creazione del pannello delle impostazioni dell'utente e la sua visualizzazione all'interno del MainPanel. Come il metodo sopra citato, viene creato dal metodo \method{createPanel()} del presenter CallHistoryPanelPresenter un elemento di tipo DOM\_node che viene poi passato al presenter \textit{main} come parametro del metodo \method{displayChildPanel}.


\item{\method{+ displayCallHistoryPanel(): void}}\\
Provoca la creazione del pannello dello storico delle chiamate e la sua visualizzazione all'interno del MainPanel. Come il metodo sopra citato, viene creato dal metodo \method{createPanel()} del presenter CallHistoryPanelPresenter un elemento di tipo DOM\_node che viene poi passato al presenter \textit{main} come parametro del metodo \method{displayChildPanel}.


\item{\method{+ onFiltredApplyedByParam(param: String): void}}\\
Funzione di \textit{callback} richiamata dai pulsanti di SearchResultPanel che comunica all'AddressBookPanelPresenter di filtrare la lista dei contatti secondo il parametro ricevuto (param). Viene scatenata l'invocazione del metodo \method{applyFilterByString} del presenter AddressBookPresenter passandogli come parametro \textit{param}.

\item{\method{+ getGroupsWhereContactsIs(contact: Object(contact)): void}}\\
Funzione di \textit{callback} che comunica all'AddressBookPanelPresenter di cercare i gruppi a cui appartiene un utente nella propria rubrica ricevuto come parametro (contact). Viene quindi invocato il metodo \method{getGroupsWhereContactIs()} del presenter AddressBookPanelPresenter.

\item{\method{+ displaySearchResultPanel(): void}}\\
Provoca la creazione del pannello delle impostazioni del proprio account e la sua visualizzazione all'interno del MainPanel. Viene creato dal metodo \method{createPanel()} del presenter AccountSettingsPanelPresenter un elemento di tipo DOM\_node che viene poi passato al presenter \textit{main} come parametro del metodo \method{displayChildPanel}.


\item{\method{+ displayCommunicationPanel(): void}}\\
Provoca la creazione del pannello delle comunicazioni e la sua visualizzazione all'interno del MainPanel. Viene creato dal metodo \method{createPanel()} del presenter CommunicationPanelPanelPresenter un elemento di tipo DOM\_node che viene poi passato al presenter \textit{main} come parametro del metodo \method{displayChildPanel}.

\item{\method{+ displayContact(contact: Object(contact)): void}}\\
L'invocazione di tale metodo provoca la visualizzazione della scheda di un contatto (contact) nel CommunicationpanelPresenter. Viene richiamato il metodo \method{displayCommunicationPanel()} e successivamente il metodo \method{display} passando come contatto il parametro ricevuto precedentemente.

\item{\method{+ contactAlreadyPresent(contact: Object(contact)): boolean}}\\
Metodo che controlla se l'utente ricevuto come parametro del metodo (contact) è gia presente nella rubrica. Viene invocato il metodo \method{contactAlreadyPresent()} di AddressBookPanelPresenter passando come parametro il contatto in questione.

\item{\method{+ onChatStarted(user: Object(contact)): void}}\\
Coordina i presenter nel momento in cui ha inizio una nuova comunicazione testuale e incapsula la collaborazione fra ContactPanelPresenter e CommunicationPanelPresenter. Viene quindi
\begin{itemize}
\item invocato il metodo \method{createPanel()} di CommunicationPanelPresenter;
\item invocato il metodo \method{displayChildPanel()} di MainPanelPresenter passandogli come parametro l'oggetto creato precedentemente;
\item invocato il metodo \method{addChat()} di CommunicationPanelPresenter passando come parametro l'utente interessato;
\item invocato il metodo \method{displayChat()} di CommunicationPanelPresenter passando come parametro l'utente interessato.
\end{itemize} 

\item{\method{+ onCall(contact: Object(contact), onlyAudio: boolean): boolean}}\\
Metodo per gestire la chiamata, riceve come parametri l'utente che si vuole chiamare (\textit{contact}) e un \textit{flag} per segnalare se la chiamata è solo audio (\textit{onlyAudio}).

\end{description}


\subsection{Package org.softwaresynthesis.mytalk.clientpresenter.kernel}\label{sec:kernel}
\classsection{CommunicationCenter}

\subsubsection*{Funzione}
Classe logica che gestisce tutta la parte della comunicazione lato \textit{client}.

\subsubsection*{Relazioni d'uso}
Nessuna relazione d'uso evidenziata.

\subsubsection*{Classi estese ed interfacce implementate}
Nessuna relazione evidenziata.

\subsubsection*{Attributi}
\begin{description}
\item{\memberdata{+ videoComunication: String[]}}\\
Array che contiene i dati della video-chiamata
\item{\memberdata{+ openChat: DOM\_element[]}}\\
Array che contiene i dati della video-chiamata
\item{\memberdata{- urlServlet: String}}\\
Contiene l'indirizzo della \inglese{servlet} che gestisce la comunicazione
\item{\memberdata{- my: String[]}}\\
Array contenente i dati personali dell'utente
\item{\memberdata{- websocket: Object()}}\\
Oggetto che rappresenta la \textit{websocket} da utilizzare per comunicare con il server. Il costruttore di tale oggetto richiede l'URL della \inglese{servlet}.

\end{description}

\subsubsection*{Metodi}
\begin{description}
\item{\method{- formatBytes(bytes: int): String}}\\
Metodo che formatta i byte ricevuti ed inviati al fine di fornire una visualizzazione sensata delle statistiche.\\
Il metodo riceve un numero rappresentante i \textit{bytes} (\memberdata{bytes} e li converte in KB/s o MB/s a seconda della grandezza.

\item{\method{- formatTime(tempo: int): String}}\\
Metodo che ritorna nel formato ``hh:mm:ss'' il tempo della comunicazione.\\
Il metodo riceve un numero rappresentante il tempo espresso in secondi (\memberdata{tempo}, lo converte nel formalismo descritto prima e lo restituisce.

\item{\method{- stopTimer(): void}}\\
Metodo che ferma l'aggiornamento del timer, viene richiamato quando la chiamata termina.

\item{\method{- stopStats(): void}}\\
Metodo che ferma l'aggiornamento delle statistiche relative alla chiamata in corso, viene richiamato quando la chiamata termina.

\item{\method{- dumpStats(obj: Object()): void}}\\
 Metodo che permette l'estrazione dei dati rappresentanti le statistiche della chiamata, esso riceve come parametro un oggetto (\memberdata{obj}), lo elabora, ed estrae solo i byte ricevuti e inviati. Ricavati questi dati, li visualizza nello \textit{span} corrispondente.
 
\item{\method{- gotDescription(): void}}\\
Metodo utilizzato da WebRTC che imposta la propria descrizione e la invia al \textit{client} chiamato in modo tale da poter instaurare la comunicazione. Verrà quindi prima di tutto utilizzata la funzione proprietaria di WebRTC \method{setLocalDescription} per impostare la propria descrizione e successivamente, mediante la \textit{websocket}, inviata al client \underline{peer} chiamato.
 
\item{\method{+ connect(): void}}\\
Metodo utilizzato per la creazione della connessione con il \textit{server}. Viene inizializzata la variabile \memberdata{websocket} richiamando il costruttore con un parametro (URL della servlet). Successivamente viene specificato il comportamento della \textit{websocket} per determinati eventi quali:
\begin{itemize}
\item \textbf{onopen}: istruzioni da eseguire quando la \textit{websocket} viene aperta. In questo caso, deve essere inviato alla \textit{websocket} il proprio \textit{id} utente in modo tale da identificare univocamente il canale aperto con il server.
\item \textbf{onclose}: istruzioni da eseguire quando la \textit{websocket} viene chiusa per qualsiasi motivo, notificando il server di tale avvenimento.
\item \textbf{onerror}: istruzioni da eseguire quando avviene un errore con la \textit{websocket}. Tale evento verrà notificato all'utente con un messaggio d'errore che descrive l'errore e l'atteggiamento da seguire.
\item \textbf{onmessage}: metodo più corposo in quanto gestisce l'unico evento richiamato mentre avviene una comunicazione dal server verso il \textit{client}. Per diversificare i messaggi ricevuti, si è deciso di rappresentare il messaggio come un \textit{array} nel quale il primo elemento identifica il tipo di richiesta. In questo modo, avremo tre tipi:
\begin{itemize}
\item type 2: riceve e gestisce una richiesta di chiamata. Imposta quindi con il metodo nativo di WebRTC \method{setRemoteDescription} la descrizione del \underline{peer} chiamante (secondo elemento dell'\textit{array} ricevuto) e la comunicazione può iniziare.
\item type 3: riceve una richiesta con l'\textit{id} del chiamante come dato. Viene memorizzato nel \textit{client} in modo tale da conoscere il canale che il chiamante ha aperto con il server per future comunicazioni.
\item type 5: notifica il cambiamento di stato degli amici. I dati ricevuti contengono l'\textit{id} dell'utente e relativo nuovo stato, si procede quindi a modificare lo stato dell'utente nella rubrica.
\end{itemize}
\end{itemize}
 
\item{\method{+ disconnect(): void}}\\
Metodo utilizzato per disconnettersi dal sistema, tale evento deve essere notificato al \textit{server} tramite la \textit{websocket} inviando un messaggio di tipo 4 contenente il proprio \textit{id}.

\item{\method{+ call(isCaller: boolean, contact: Object(contact)): void}}\\
Metodo che gestisce la vera e propria chiamata tramite WebRTC.\\
Inizialmente, si inizializza la variabile \memberdata{pc} che è di tipo RTCPeerConnection, oggetto reso disponibile dalle librerie di WebRTC, successivamente si definiscono i comportamenti associati agli eventi di RTCPeerConnection, ossia:
\begin{itemize}
\item onicecandidate: evento scatenato quando un nuovo \underline{peer} si ``candida'' per poter chiamare, viene quindi inviata la propria ``descrizione'' all'altro \underline{peer}.
\item onaddstream: evento scatenato quando viene aggiunto uno \textit{stream} nell'oggetto RTCPeerConnection. In questa situazione, il chiamato riceve lo \textit{stream} del chiamante e lo visualizza nell'apposito tag video. In contemporanea, vengono richiamati i metodi che fanno partire la visualizzazione delle statistiche di chiamata.
\item onremovestream: evento scatenato quando viene rimosso uno \textit{stream} nell'oggetto RTCPeerConnection. In questa situazione, colui che effettua questa operazione vuole terminare la chiamata e deve quindi re-inviare la propria descrizione (attraverso il metodo \method{gotDescription()} in modo tale che l'altro \underline{peer} possa aggiornare la descrizione remota. Vengono richiamati alla fine i metodi che fermano la visualizzazione delle statistiche di chiamata e chiudono il canale di comunicazione tramite la chiamata del metodo \method{close()} dell'oggetto \memberdata{pc}.
\end{itemize}

\item{\method{+ webkitGetUserMedia(audio: boolean, video: boolean): void}}\\
Metodo richiamato nel momento in cui un utente vuole iniziare una chiamata. Il metodo, fornito dalle librerie native di \textit{Google Chrome}, cattura lo \textit{stream} della propria videocamera/microfono. Riceve due parametri che indicano se deve essere catturato audio e/o video. Procedendo, viene aggiunto lo \textit{stream} all'oggetto RTCPeerConnection tramite il metodo \method{addStream(localstream)} fornito dalle librerie di WebRTC.

\item{\method{+ endCall(): void}}\\
Metodo richiamato per terminare una chiamata, a cascata viene richiamato il metodo \method{removeStream(localstream)} e inviata la nuova descrizione al \underline{peer} remoto.
 
 \end{description}

\clearpage

\section{Specifica sotto-architettura \texttt{clientview}}\label{sec:clientviewarchitecture}
Il sistema dispone di dodici viste, ognuna associata ad uno specifico \inglese{presenter}. Tali viste vengono create e modificate da parte dei relativi \inglese{presenter} invocando i metodi che le viste rendono disponibili tramite la loro API pubblica.

La descrizione di ciascuna vista riportata nel presente documento comprende, oltre alla funzione, agli attributi e ai metodi che essa mette a disposizione del relativo \inglese{presenter}, anche la struttura degli elementi che la compongono.

\subsection{Package \texttt{org.softwaresynthesis.mytalk.clientview}}

\classsection{AccountSettingsView}
\subsubsection*{Funzione}
Vista associata al \inglese{presenter} \classname{AccountSettingsPresenter} che ha il compito di permettere la visualizzazione e la modifica dei dati dell'utente.

\subsubsection*{Attributi}
\begin{description}
  \item{\memberdata{-- thisPresenter: Object}}\\
  Riferimento al \inglese{presenter} associato alla vista al quale vengono inoltrati i comandi relativi alla logica applicativa tramite i gestori di eventi associati agli elementi grafici del pannello.
\end{description}

\subsubsection*{Metodi}
\begin{description}
\item{\method{+ display(): void}}\\
Tale metodo ha il compito di inizializzare la vista, ovvero l'insieme dei \inglese{widget} grafici che costituiscono il pannello per la visualizzazione delle impostazioni e dei dati dell'utente.

In particolare, la vista ha il compito di ottenere i dati relativi a nome, cognome e immagine del profilo dell'utente e popolare le apposite regioni del pannello con le informazioni corrispondenti, quindi configurare il comportamento del pulsante per la modifica dei dati alla cui pressione deve corrispondere un'invocazione del metodo \method{onChangeButtonPressed()}.

\item{\method{+ onChangeButtonPressed(): void}}\\
Gestisce la pressione del pulsante ``Modifica dati'' trasformando l'elemento \verb'<div>' che contiene l'immagine del profilo e la lista di elementi testuali in un \textit{form} da compilare per apportare i cambiamenti i propri dati di registrazione. È importante che quest'ultimo abbia per l'attributo \verb'enctype' il valore \verb'multipart/form-data'.

Il \inglese{form} è costituito da una lista i cui tre elementi \verb'<li>' contengono rispettivamente:
\begin{itemize}
  \item[--] una \verb'<label>' e il relativo campo \verb'<input>' per il nome, che inizialmente visualizza il nome precedentemente impostato nel client;
  \item[--] una \verb'<label>' e il relativo campo \verb'<input>' per il cognome dell'utente, che inizialmente visualizza il valore precedentemente impostato nel client;
  \item[--] una \verb'<label>' e il relativo campo \verb'<input>' per l'immagine del profilo (quest'ultimo dovrà avere l'attributo \verb'type' uguale a \verb'file').
\end{itemize}

Inoltre, il \inglese{form} presenta un pulsante per l'invio dei dati inseriti al server alla cui pressione deve essere attivato un gestore di evento che raccoglie i dati inseriti e li trasmette invocando il metodo \method{sendUserData} del \inglese{presenter} e quindi solleva un evento con identificativo \verb'showAccountSettingPanel' per ricaricare il pannello.

\end{description}

\subsubsection*{Elementi}
Il pannello è costituito da un elemento \verb'<div>' con identificativo uguale a \verb'AccountSettingsPanel' che contiene al proprio interno:
\begin{itemize}
  \item[--] un \verb'<div>' con attributo \verb'class' uguale a \verb'panelHeader' con all'interno un elemento \verb'<h1>' avente come figlio il solo nodo testo \textsc{<<dati personali>>};
  \item[--] un elemento \verb'<img>' con identificativo uguale alla stringa \verb'picture';
  \item[--] un elemento \verb'<ul>' contenente:
  \begin{itemize}
    \item[-] un \verb'<li>' con identificativo uguale a \verb'name';
    \item[-] un \verb'<li>' con identificativo uguale a \verb'surname';
  \end{itemize}
  \item[--] un elemento \verb'<button>' contenente il testo ``Modifica dati''.
\end{itemize}


% Qui ci andrà ADDRESSBOOKVIEW


\classsection{CallHistoryView}
\subsubsection*{Funzione}
Vista associata al \inglese{presenter} \classname{CallHistoryPresenter} utilizzata per la visualizzazione dello storico chiamate di un utente del sistema.

\subsubsection*{Attributi}
\begin{description}
  \item{\memberdata{-- thisPresenter: Object}}\\
  Riferimento al \inglese{presenter} associato alla vista al quale vengono inoltrati i comandi relativi alla logica applicativa tramite i gestori di eventi associati agli elementi grafici del pannello.
\end{description}

\subsubsection*{Metodi}
\begin{description}

  \item{\method{+ addListItem(call: Object): void}}\\
  Metodo che aggiunge alla lista delle chiamate visualizzata nel pannello una nuova voce che corrisponde alla chiamata passata come parametro.
  
  Ogni voce della lista contiene i dati relativi alla chiamata vale a dire l'altro utente coinvolto e la data della chiamata. Poiché a partire da questo pannello non è possibile effettuare alcuna azione, non è presente alcuna funzione di \inglese{callback} associata alle voci della lista.
  
  \item{\method{+ display(): void}}\\
  Tale metodo ha il compito di inizializzare il pannello della vista con i dati corretti popolando la lista delle chiamate in base ai dati ricevuti dal \inglese{presenter}.
  
\end{description}

\subsubsection*{Elementi}
Il pannello è costituito da un elemento \verb'<div>' avente come identificativo il valore \verb'CallHistoryPanel' il quale, a sua volta, contiene:
\begin{itemize}
  \item[--] un elemento \verb'<div>' avente attributo \verb'class' uguale al valore \verb'panelHeader' e con all'interno un elemento \verb'<h1>' contenente il titolo \textsc{<<storico chiamate>>};
  \item[--] un elemento \verb'<ul>' avente come identificativo la stringa \verb'ulHistory'.
\end{itemize}


% Qui ci andrà COMMUNICATIONVIEW


% Qui ci andrà CONTACTVIEW


\classsection{GroupView}
\subsubsection*{Funzione}
Vista associata al \classname{GroupPresenter} utilizzata per la visualizzazione dei gruppi della rubrica.

\subsubsection*{Attributi}
\begin{description}
  \item{\memberdata{-- thisPresenter: Object}}\\
  Riferimento al \inglese{presenter} associato alla vista al quale vengono inoltrati i comandi relativi alla logica applicativa tramite i gestori di eventi associati agli elementi grafici del pannello.
\end{description}

\subsubsection*{Metodi}
\begin{description}
  
  \item{\method{+ addListItem(group: Object): void}}\\
  Il metodo ha lo scopo di aggiungere una nuova voce alla lista con identificativo \verb'groupList' presente nel pannello. Ogni nuova \verb'<li>' dovrà essere a sua volta composta da:
  \begin{itemize}
    \item[--] un pulsante sotto forma di immagine per mostrare il contenuto del gruppo con l'attributo \texttt{class} uguale a \texttt{showGroupContacts} e percorso \texttt{img/expandGroupImg.png};
    \item[--] un elemento \verb'<span>' con un nodo di testo  che rappresenta il nome gruppo stesso;
    \item[--] un pulsante sotto forma di immagine con attributo \texttt{class} uguale a \texttt{deleteGroup} e percorso \texttt{img/deleteGroupImg.png} per la rimozione del gruppo;
    \item[--] un pulsante sotto forma di immagine con attributo \texttt{class} uguale a \texttt{addContactInGroup} e percorso \texttt{img/addToGroupImg.png} per l'aggiunta di un nuovo contatto all'interno del gruppo stesso;
    \item[--] un elemento \verb'<ul>' che rappresenta la lista (in forma estesa oppure compatta) dei contatti della rubrica che appartengono al gruppo in oggetto.
  \end{itemize}
  
  Il metodo dovrà inoltre associare ai tre pulsanti le funzioni di \inglese{callback} per la gestione dell'evento \verb'click', in particolare:
  \begin{itemize}
    \item[--] alla pressione del pulsante per invertire la modalità di visualizzazione del gruppo, dovrà essere cambiato l'attributo \verb'class' della lista dei contatti del gruppo da \verb'collapsedList' a \verb'uncollapsedList' (o viceversa) e, nel caso in cui il pulsante sia utilizzato per visualizzare un gruppo in forma compatta, dovrà essere rimosso il \inglese{form} mediante il quale sarà possibile aggiungere nuovi contatti al gruppo;
    \item[--] alla pressione del pulsante per aggiungere un nuovo contatto al gruppo dovrà essere impostata la visualizzazione compatta dei membri del gruppo, quindi dovrà essere creato il \inglese{form} per l'aggiunta di un nuovo contatto al gruppo utilizzando il metodo privato \method{createAddToGroupForm} della vista stessa;
    \item[--] alla pressione del pulsante per l'eliminazione di un gruppo, dovrà essere chiesta conferma all'utente con un'istruzione del tipo
    \begin{verbatim}
confirm("Sei sicuro di voler eliminare il gruppo" + group.name + "?");
    \end{verbatim}
    e, in caso di risposta affermativa, dovrà essere sollevato un evento con nome identificativo uguale a \verb'groupRemoved' cui risponderà il \inglese{presenter} incaricato di inviare al server la richiesta corrispondente.
  \end{itemize}
  
  \item{\method{+ display(): void}}\\
  Il metodo ha il compito di inizializzare la vista, ovvero l'insieme dei \inglese{widget} grafici che costituiscono il pannello per la gestione dei gruppi della rubrica.
  
  Dal momento che per popolare il pannello con i dati necessari è necessario prima scaricarli dal server, la vista invoca il metodo \method{displayGroupList()} del relativo \inglese{presenter} per portare a termine l'operazione.
  
  \item{\method{-- createAddToGroupForm(group: Object): HTMLElement}}\\
  Un'invocazione di tale metodo dovrà fare in modo che sia creato il \inglese{form} per l'aggiunta di un nuovo contatto al gruppo passato come parametro in ingresso. La vista dovrà innanzitutto ottenere un insieme dei contatti per cui è possibile l'aggiunta nel suddetto gruppo, mediante il metodo \method{selectCandidates} fornito dal \inglese{presenter} associato.
  
  In secondo luogo, dovrà essere creato un elemento \verb'<fieldset>' contenente una serie di \verb'<input>' (di tipo \verb'checkbox') e \verb'<label>' per ognuno dei contatti candidati all'aggiunta.
  
  Infine, dovranno essere aggiunti al \verb'<form>' due pulsanti, uno per determinare l'aggiunta di tutti i contatti selezionati all'interno del gruppo visualizzato e uno per terminare l'operazione di inserimento nascondendo il relativo \inglese{form}. Il metodo dovrà inoltre terminare restituendo al chiamante un riferimento all'elemento HTML \verb'<form>' appena creato.
  
  \item{\method{-- createContactList(group: Object): HTMLElement}}\\
  Questo metodo ha il compito di creare la lista dei contatti che appartengono al gruppo passato come parametro in ingresso, che dovrà avere di default il valore dell'attributo \verb'class' impostato a \verb'collapsedList'.
  
  Dovrà inoltre essere creato un nuovo elemento \verb'<ul>' all'interno del quale aggiungere un nodo \verb'<li>' per ogni contatto del gruppo. Le voci di questa lista saranno formate da un nodo di testo contenente il nome del contatto, più un pulsante sotto forma di immagine avente la classe \verb'deleteContact' e percorso uguale a \verb'img/deleteContactImg.png'.
  
  La lista così costituita dovrà, infine, essere restituita al chiamante del metodo.

\end{description}

\subsubsection*{Elementi}
Il pannello è rappresentato da un elemento \verb'<div>' con identificativo uguale a \verb'GroupPanel' ed è costituito dai seguenti elementi:
\begin{itemize}
  \item[--] un \verb'<div>' con attributo \verb'class' uguale a \verb'panelHeader' con all'interno un elemento \verb'<h1>' contenente il solo nodo di testo \textsc{gestione gruppi};
  \item[--] una \verb'<ul>' con identificativo uguale a \verb'groupList'.
\end{itemize}


\classsection{LoginView}
\subsubsection*{Funzione}
Vista associata al \inglese{presenter} \classname{LoginPresenter} utilizzata per effettuare il login nel sistema.

\subsubsection*{Attributi}
\begin{description}
  \item{\memberdata{-- thisPresenter: Object}}\\
  Riferimento al \inglese{presenter} associato alla vista al quale vengono inoltrati i comandi relativi alla logica applicativa tramite i gestori di eventi associati agli elementi grafici del pannello.
\end{description}

\subsubsection*{Metodi}
\begin{description}

  \item{\method{+ correctAnswer(): void}}\\
	Metodo usato per modificare il nodo DOM salvato nella variabile \memberdata{thisPanel} per comunicare all'utente che la risposta da lui inserita per il recupero password è corretta. 
	
Il metodo deve rimuovere il figlio con identificativo uguale a \texttt{passwordretrieval} dalla vista, quindi aggiunge un nuovo elemento \verb'<p>' contenente il messaggio:
\begin{verbatim}
	Recupero password avvenuto correttamente.
	Ti è stata inviata un'email contenente i dati richiesti.
\end{verbatim}
il messaggio citato è temporaneo e dovrà essere rimosso allo scadere di un \inglese{timeout} della durata di due secondi utilizzando la funzione \verb'window.setTimeout(funtion() {...}, 2000)'.

Il corpo della funzione anonima passata a \verb'setTimeout' avrà il compito di rimuovere dall'elemento \verb'thisPanel' il paragrafo creato in precedenza.

  \item{\method{+ destroy(): void}}\\
  Tale metodo ha il compito di eliminare il pannello dall'interfaccia grafica visualizzata. A tal fine, deve essere rimosso il nodo che corrisponde al pannello nell'albero DOM della UI\@.
  
  \item{\method{+ errorLogin(): void}}\\
  Il metodo ha lo scopo di segnalare visivamente all'utente il fatto che si è verificato un errore nell'elaborazione delle credenziali di autenticazione inserite nel \inglese{form}. Tale effetto dovrà essere ottenuto aggiungendo al valore dell'attributo \verb+class+ il valore \verb+error+.

	\item{\method{+ incorrectAnswer(): void}}\\
	Metodo che in caso di inserimento della risposta non corretta alla domanda segreta visualizza un messaggio di avvertimento all'utente per due secondi quindi lascia il controllo al \inglese{form} di inserimento della risposta alla domanda segreta.
	
	Il metodo deve cercare un elemento \verb'<p>' per la visualizzazione del messaggio, contenente il seguente testo:
\begin{verbatim}
	Dati non corretti. Inserire nuovamente la risposta.
\end{verbatim}
	
	Tale messaggio dovrà essere visualizzato con una procedura analoga a quella usata in \method{correctAnswer()}.

	\item{\method{+ initialize(view: HTMLElement): void}}\\
	Il metodo ha il compito di inizializzare il pannello di login dopo che è stato caricato il relativo \inglese{template} che crea di tutti i \inglese{widget} grafici contenuti al suo interno.
	
	Oltre ad aggiungere all'interfaccia grafica il sotto-albero DOM che riceve come parametro in ingresso, il metodo dovrà configurare il pannello assegnando ai pulsanti le funzioni di \inglese{callback} corrispondenti.
	
	In particolare, alla pressione del pulsante di login deve essere eseguita una funzione che recupera il nome utente e la password inseriti dall'utente (con i metodi \method{getUsername()} e \method{getPassword()} rispettivamente), quindi invocare il metodo \method{login} del \inglese{presenter} per la valutazione dei dati.
	
	Alla pressione del pulsante di registrazione deve invece essere sollevato un evento che provoca la visualizzazione del pannello di registrazione, vale a dire un evento con stringa identificativa uguale a \verb'showRegistrationPanel'.
	
	Infine, alla pressione del pulsante per il recupero della password deve essere costruito il \inglese{form} corrispondente, mediante una chiamata al metodo \method{buildRetrievePasswordForm()}, aggiungendo l'elemento restituito all'interfaccia grafica finale.
	
	\item{\method{-- buildRetrievePasswordForm(): HTMLElement}}\\
	Metodo usato per costruire il \inglese{form} per il recupero della password, che ha il compito di costruire un elemento \verb'fieldset', che contiene al suo interno:
	\begin{itemize}
	  \item[--] un \verb'<label>' contenente la domanda segreta;
	  \item[--] un elemento \verb'<input>' per inserire il testo della risposta;
	  \item[--] un \verb'<input>' di tipo \verb'submit' per l'invio dei dati.
	\end{itemize}
Alla pressione di quest'ultimo, la vista dovrà invocare il metodo \method{hasAnsweredCorrectly} del \inglese{presenter} associato e, in base al risultato positivo o negativo ottenuto, visualizzare il messaggio corrispondente tramite i metodi \verb'correctAnswer()' e \verb'incorrectAnswer()' rispettivamente.

Infine, il metodo restituisce il nodo DOM radice della nuova porzione di interfaccia grafica creata, vale a dire il \verb'<fieldset>'.

\end{description}

\subsubsection*{Elementi}
Il pannello contiene un elemento \verb'<div>' con identificativo uguale a \verb'LoginPanel' il quale a sua volta ha come figlio un \verb'<fieldSet>' con identificativo uguale a \verb'loginForm' che contiene:
  \begin{itemize}
    \item[--] un elemento \verb'<div>' per il logo;
    \item[--] un elemento \verb'<ul>' contenente i seguenti elementi:
    \begin{itemize}
      \item[-] un elemento \verb'<li>' per il nome utente, con all'interno un elemento \verb'<label>' e il relativo campo di inserimento \verb'<input>';
      \item[-] un elemento \verb'<li>' per la password, con all'interno un elemento \verb'<label>' e il relativo campo di inserimento \verb'<input>';
      \item[-] \verb'<li>' per i pulsanti contenente al suo interno:
      \begin{itemize}
        \item[$\cdot$] un elemento \verb'<button>' per effettuare il login;
        \item[$\cdot$] un elemento \verb'<button>' per effettuare la registrazione al sistema;
        \item[$\cdot$] un elemento \verb'<button>' per recuperare i dati dimenticati.
      \end{itemize}
    \end{itemize}
  \end{itemize}
  

\classsection{MainView}
\subsubsection*{Funzione}
Vista associata al \classname{MainPresenter} che contiene il pannello principale ed è responsabile del suo aggiornamento e della sua visualizzazione.

\subsubsection*{Attributi}
\begin{description}
  \item{\memberdata{-- thisPresenter: Object}}\\
  Riferimento al \inglese{presenter} associato alla vista al quale vengono inoltrati i comandi relativi alla logica applicativa tramite i gestori di eventi associati agli elementi grafici del pannello.
\end{description}

\subsubsection*{Metodi}
\begin{description}
  \item{\method{+ destroy(): void}}\\
  Metodo che ha il compito di rimuovere il pannello dall'interfaccia grafica visualizzata rimuovendo il nodo corrispondente al pannello principale dalla sua posizione all'interno dell'albero DOM\@.
  
  \item{\method{+ displayChildPanel(node: HTMLElement): void}}\\
  Provoca la visualizzazione di un pannello secondario all'interno della \classname{MainView} aggiungendo la sua radice, vale a dire il nodo DOM ricevuto come parametro in ingresso, come elemento figlio del \verb'<div>' che rappresenta il pannello principale.

  \item{\method{+ initialize(view: HTMLElement): void}}\\
  Tale metodo ha il compito di visualizzare del pannello principale, che viene passato tramite il parametro in ingresso \verb'view', aggiungendolo nella posizione corretta dell'albero DOM dell'interfaccia grafica in modo che venga ad occupare la posizione centrale.
  
\end{description}

\subsubsection*{Elementi}
Il pannello è costituito da un solo elemento \verb'<div>' avente come attributo identificativo la stringa \verb'MainPanel' all'interno del quale sarà visualizzata, di volta in volta, una delle sette viste di secondo livello, vale a dire:
\begin{itemize}
  \item[--] \classname{AccountSettingsView};
  \item[--] \classname{CallHistoryView};
  \item[--] \classname{CommunicationView};
  \item[--] \classname{ContactView};
  \item[--] \classname{GroupView};
  \item[--] \classname{MessageView};
  \item[--] \classname{SearchResultView}.
\end{itemize}


\classsection{MessageView}
\subsubsection*{Funzione}
Vista associata al \inglese{presenter} \classname{MessagePresenter} utilizzata per visualizzare i messaggi presenti nella segreteria telefonica dell'utente.

\subsubsection*{Attributi}
\begin{description}
  \item{\memberdata{-- thisPresenter: Object}}\\
  Riferimento al \inglese{presenter} associato alla vista al quale vengono inoltrati i comandi relativi alla logica applicativa tramite i gestori di eventi associati agli elementi grafici del pannello.
\end{description}

\subsubsection*{Metodi}
\begin{description}

\item{\method{+ addListItem(message: Object): void}}\\
  Tale metodo riceve in input un oggetto che corrisponde a un messaggio della segreteria, caratterizzato dalle proprietà \verb'sender', \verb'date', \verb'status', \verb'video' e \verb'src'.
  
  Sulla base delle proprietà del messaggio, il metodo ha il compito di creare un nuovo \verb'<li>' e aggiungere ad esso:
  \begin{itemize}
    \item[--] un elemento \verb'<img>' che corrisponde allo stato letto/non letto del messaggio impostato tramite una chiamata al metodo privato \method{getStatusSrc} della vista;
    \item[--] un nodo di testo con il nome dell'utente che ha inviato il messaggio;
    \item[--] un nodo di testo contenente la data del messaggio;
    \item[--] un elemento \verb'<img>' per l'eliminazione del messaggio, che deve avere attributo \verb'src'
  \end{itemize}
  
  Prima di aggiungere il nuovo \inglese{list item} in coda alla lista dei messaggi, il metodo dovrà configurare il comportamento delle parti dell'interfaccia. In particolare, un click sulla voce della lista marchi il messaggio come letto tramite il metodo \method{setStatus} del \inglese{presenter} e inizializzi l'elemento \verb'<video>' per la visualizzazione dello stesso.
  
  Un click sull'immagine per l'eliminazione del messaggio deve invece chiamare l'apposito metodo \method{deleteMessage} messo a disposizione dal \inglese{presenter} e, infine, un click sull'immagine per cambiare lo stato deve invertire lo stato del messaggio tramite il metodo \method{setStatus} del \inglese{presenter} e cambiare l'immagine nel pannello che rappresenta lo stato letto/non letto del messaggio.

  \item{\method{+ display(): void}}\\
  Il metodo ha il compito di inizializzare la vista, ovvero l'insieme dei \inglese{widget} grafici che costituiscono il pannello per la visualizzazione della segreteria telefonica dell'utente.
  
  Il metodo ha quindi il compito di popolare la lista dei messaggi della segreteria con i dati corretti che sono stati scaricati dal server.
  
  \item{\method{-- getStatusSrc(message: Object): String}}\\
  Il metodo ha il compito di restituire la stringa corretta che rappresenta il percorso relativo dell'immagine associata allo stato del messaggio. Se il messaggio è nuovo (non letto) il risultato dovrà essere \verb'img/unreadmessage.png', in caso contrario  \verb'img/readmessage.png'.
  
\end{description}

\subsubsection*{Elementi}
Il pannello è costituito da un elemento \verb'<div>' con identificativo uguale a \verb'MessagePanel' che contiene al proprio interno:
\begin{itemize}
  \item[--] un \verb'<div>' con attributo \verb'class' uguale a \verb'panelHeader' con al proprio interno un \verb'<h1>' con il solo nodo di testo \textsc{<<segreteria>>};
  \item[--] un elemento \verb'<video>' con identificativo uguale a \verb'messageVideo' e attributo \verb'controls' uguale a \verb'controls', contenente un elemento \verb'<source>';
  \item[--] un \verb'<div>' con identificativo uguale a \verb'divMessage' contenente un elemento \verb'<ul>' con identificativo uguale a \verb'messageList'.
\end{itemize}


\classsection{RegisterView}
\subsubsection*{Funzione}
Vista associata al \inglese{presenter} \classname{RegisterPresenter} utilizzata per effettuare la registrazione.

\subsubsection*{Attributi}
\begin{description}
  \item{\memberdata{-- thisPresenter: Object}}\\
  Riferimento al \inglese{presenter} associato alla vista al quale vengono inoltrati i comandi relativi alla logica applicativa tramite i gestori di eventi associati agli elementi grafici del pannello.
\end{description}

\subsubsection*{Metodi}
\begin{description}
  \item{\method{+ destroy(): void}}\\
  Tale metodo ha il compito di eliminare il pannello dall’interfaccia grafica visualizzata. A tal fine, deve essere rimosso il nodo che corrisponde al pannello nell’albero DOM della UI\@.
  
  \item{\method{+ initialize(view: HTMLElement): void}}\\
  	Il metodo ha il compito di inizializzare il pannello di registrazione al sistema una volta caricato il relativo \inglese{template} che crea di tutti i \inglese{widget} grafici contenuti al suo interno.
	
	Oltre ad aggiungere all'interfaccia grafica il sotto-albero DOM che riceve come parametro in ingresso, il metodo dovrà configurare il pannello assegnando ai pulsanti le funzioni di \inglese{callback} corrispondenti.
	
	In particolare, alla pressione del pulsante ``Indietro'' deve essere sollevato un evento avente come nome identificativo la stringa \verb'showLoginPanel'.
	
  Alla pressione del pulsante per la registrazione, invece, devono essere recuperati i valori di tutti i campi dati del form mediante i metodi \method{getUsername()}, \method{getPassword()}, \method{getQuestion()}, \method{getAnswer()}, \method{getName()}, \method{getSurname()}, \method{getPicturePath()} e, in un secondo momento, deve essere richiamato il metodo \method{register} del \inglese{presenter}.

\end{description}

\subsubsection*{Elementi}
Il pannello è costituito da n elemento di tipo \verb'<div>' con identificativo uguale a \verb'LoginPanel' che contiene un elemento \verb'<form>' con al proprio interno:
\begin{itemize}
  \item[--] un elemento \verb'<ul>' il quale contiene:
  \begin{itemize}
    \item[-] un \verb'<li>' per il nome utente, costituito a sua volta da una \verb'<label>' e dal relativo campo \verb'<intput>' avente attributo \verb'type' uguale a \verb'email';
    \item[-] un \verb'<li>' per la password, costituito a sua volta da una \verb'<label>' e dal relativo campo \verb'<intput>' avente attributo \verb'type' uguale a \verb'password';
    \item[-] un \verb'<li>' per la domanda segreta, costituito a sua volta da una \verb'<label>' e dal relativo campo \verb'<intput>';
    \item[-] un \verb'<li>' per la risposta, costituito a sua volta da una \verb'<label>' e dal relativo campo \verb'<intput>';
    \item[-] un \verb'<li>' per il nome, costituito a sua volta da una \verb'<label>' e dal relativo campo \verb'<intput>';
    \item[-] un \verb'<li>' per il cognome, costituito a sua volta da una \verb'<label>' e dal relativo campo \verb'<intput>';
    \item[-] un \verb'<li>' per l'immagine del profilo, costituito a sua volta da una \verb'<label>' e dal relativo campo \verb'<intput>' avente attibuto \verb'type' uguale a \verb'file';
    \item[-] un \verb'<li>' che contiene un avviso relativo all'obbligatorietà dei campi dati nome utente, password, domanda e risposta;
  \end{itemize}
  \item[--] un \verb'<div>' per i pulsanti che ha al proprio interno:
  \begin{itemize}
    \item[-] un elemento \verb'<button>' per tornare al pannello di login;
    \item[-] un elemento \verb'<button>' per portare a termine la registrazione al sistema.
  \end{itemize}
\end{itemize}


\classsection{SearchResultView}
\subsubsection*{Funzione}
Vista associata al \inglese{presenter} \classname{SearchResultPresenter} utilizzata per visualizzare i risultati di una ricerca fra gli utenti registrati al sistema.

\subsubsection*{Attributi}
\begin{description}
  \item{\memberdata{-- thisPresenter: Object}}\\
  Riferimento al \inglese{presenter} associato alla vista al quale vengono inoltrati i comandi relativi alla logica applicativa tramite i gestori di eventi associati agli elementi grafici del pannello.
\end{description}

\subsubsection*{Metodi}
\begin{description}
  
  \item{\method{+ addListItem(contact: Object): void}}\\
  Provoca la visualizzazione nel pannello di una nuova voce della lista contenente i dati del contatto passato come parametro. L'elemento \verb'<li>' che dovrà essere aggiunto in coda alla lista deve avere al proprio interno un nodo di testo corrispondente al nome del contatto, un elemento \verb'<img>' corrispondente allo stato del contatto e un elemento \verb'<img>' corrispondente all'immagine dello stesso.
  
  Inoltre, all'evento \verb'click' sollevato sul \verb'<li>' dovrà essere associato un gestore che a sua volta avrà il compito di procurarsi un riferimento a un evento con identificativo uguale a \verb'showContactPanel', assegnare alla proprietà \verb'contact' di quest'ultimo il contatto passato come parametro del metodo e sollevare l'evento corrispondente.
  
  \item{\method{+ display(): void}}\\
  Il metodo configura il comportamento della vista per la visualizzazione dei risultati di una ricerca. In particolare, deve creare un gestore per l'evento identificato dalla stringa \verb'click' associato al pulsante di ricerca che avvii una ricerca in base al contenuto inserito nell'\verb'<input>' del pannello e visualizzi i risultati mediante il metodo \method{displayResult}.

  \item{\method{-- getImageSrc(contact: Object): String}}\\
  Il metodo ha il compito di fornire l'indirizzo sul server dell'immagine corrispondente a uno stato utente. In base al valore della proprietà \verb'state' del contatto passato come parametro, il metodo dovrà restituire una fra le stringhe:
  \begin{itemize}
    \item[-] \verb'img/stateavailable.png';
    \item[-] \verb'img/stateoccupied.png';
    \item[-] \verb'img/stateoffline.png'.
  \end{itemize}

\end{description}

\subsubsection*{Elementi}
Il pannello è costituito da un elemento \verb'<div>' avente come identificativo \verb'SearchResultPanel' il quale contiene:
\begin{itemize}
  \item[--] un elemento \verb'<div>' avente l'attributo \verb'class' uguale a \verb'panelHeader', che contiene a sua volta al proprio interno un elemento \verb'<h1>' con il titolo \textsc{<<ricerca utenti>>};
  \item[--] un elemento \verb'<input>' con identificativo uguale a \verb'searchInputField';
  \item[--] un elemento \verb'<button>' avente l'attributo \verb'class' uguale a \verb'searchInputButton' e identificativo \verb'searchInputButton';
  \item[--] un elemento \verb'<ul>' per presentare i risultati della ricerca, avente come identificativo il valore \verb'userlList'.
\end{itemize}


\classsection{ToolsView}
\subsubsection*{Funzione}
Vista associata al \inglese{presenter} \classname{ToolsPresenter} utilizzata per visualizzare gli strumenti e le funzionalità rese disponibili all'utente.

\subsubsection*{Attributi}
\begin{description}
  \item{\memberdata{-- thisPresenter: Object}}\\
  Riferimento al \inglese{presenter} associato alla vista al quale vengono inoltrati i comandi relativi alla logica applicativa tramite i gestori di eventi associati agli elementi grafici del pannello.
\end{description}

\subsubsection*{Metodi}
\begin{description}

  \item{\method{+ addCommunicationFunction(): void}}\\
  Questo metodo dovrà provocare l'aggiunta di un nuovo elemento alla lista degli strumenti per permettere di ritornare al pannello delle comunicazioni.
  
  Dovrà pertanto creare un nuovo elemento \verb'<li>', assegnare alla sua proprietà identificativa il valore \verb'CallFunction', configurarne il comportamento alla pressione del mouse facendo in modo che sia sollevato un evento con nome \verb'showCommunicationPanel' quindi procurarsi un riferimento al nodo DOM che corrisponde alla lista delle funzionalità e aggiungere in coda il nuovo \inglese{list item} creato con il metodo \verb'appendChild'.
  
  \item{\method{+ destroy(): void}}\\
  Tale metodo ha il compito di provocare la rimozione del pannello degli strumenti governato da questa vista, rimuovendo in maniera definitiva dall'albero DOM il nodo che corrisponde alla radice del pannello.
  
  \item{\method{+ initialize(view: HTMLElement): void}}\\
  Il metodo ha la funzione di aggiungere nel punto corretto dell'albero DOM che rappresenta l'interfaccia utente il nodo passato come parametro, vale a dire il pannello degli strumenti, quindi di configurare il comportamento quest'ultimo assegnando a tutti i \inglese{list item} delle funzionalità i gestori di eventi corrispondenti.
  
  In particolare:
  \begin{itemize}
    \item[--] alla pressione della voce con identificativo \verb'liAnswering' deve essere sollevato un evento con nome \verb'showMessagePanel';
    \item[--] alla pressione della voce con identificativo \verb'liSetting' deve essere sollevato un evento con nome \verb'showAccountSettingPanel';
    \item[--] alla pressione della voce con identificativo \verb'liCallList' deve essere sollevato un evento con nome \verb'showCallHistoryPanel';
    \item[--] alla pressione della voce con identificativo \verb'liGroup' deve essere sollevato un evento con nome \verb'showGroupPanel';
    \item[--] alla pressione della voce con identificativo \verb'liSearch' deve essere sollevato un evento con nome \verb'showSearchResultPanel';
    \item[--] alla pressione della voce con identificativo \verb'liLogout' deve essere visualizzata una finestra di dialogo che richiede conferma all'utente sulla reale intenzione di uscire e, in caso di risposta affermativa, solleva sia l'evento \verb'logout' che \verb'showLoginPanel'.
  \end{itemize}
  
  Infine, prima di ritornare deve essere inizializzato l'elemento \verb'<select>' per la gestione dello stato utente, mediante l'opportuno metodo privato \method{initializeSelectState()}.
  
  \item{\method{+ removeCommunicationFunction(): void}}\\
  Il metodo ha il compito di rimuovere il pulsante per ritornare al pannello della comunicazione dalle opzioni disponibili nel momento in cui non sussiste più alcuna comunicazione (audio, audio/video o testuale) attiva per il client
  
   Il metodo dovrà pertanto ottenere un riferimento al nodo DOM che rappresenta la lista delle funzioni, uno a quello del \inglese{list item} corretto e quindi utilizzare il metodo \verb'removeChild' per operare la rimozione.
  
  \item{\method{+ updateStateValue(): HTMLElement}}\\
  Questo metodo ha il compito di ottenere un riferimento all'elemento \verb'<select>' utilizzato per il cambio di stato, estrarre il valore associato all'opzione selezionata mediante un'istruzione del tipo
  \begin{verbatim}
    var state = select.options[selectedIndex].value;
  \end{verbatim}
  quindi sollevare un opportuno evento avente come identificativo la stringa \verb'changeMyState' e come proprietà \verb'state' il valore dello stato recuperato in precedenza.
  
  Il metodo, infine, dovrà restituire al chiamante un riferimento all'elemento \verb'<select>' che è stato manipolato.
  
  \item{\method{-- initializeSelectState(): HTMLElement}}\\
  Il metodo ha il compito di inizializzare l'elemento \verb'<select>' per il cambio di stato presente sul pannello degli strumenti aggiungendo ad esso, se non sono già presenti, le opzioni che corrispondono allo stato ``disponibile'' e ``occupato'' utilizzando una sintassi del tipo:
  \begin{verbatim}
    var availableOption = new Option("Disponibile", "available");
    selectState.add(availableOption, null);
  \end{verbatim}
  
  È importante, inoltre, associare all'evento \verb'change' associato all'elemento \verb'<select>' una funzione anonima che provoca l'invocazione del metodo \method{updateStateValue()} della vista stessa.
  
  Infine, l'elemento \verb'<select>' che è stato manipolato da questo metodo dovrà essere restituito al chiamante.

\end{description}

\subsubsection*{Elementi}
Il pannello è costituito da un elemento \verb'<div>' avente come identificativo la stringa \verb'ToolsPanel', il quale contiene al proprio interno:
\begin{itemize}
  \item[--] un elemento \verb'<div>' avente attributo \verb'class' uguale a \verb'panelHeader';
  \item[--] un \verb'<div>' che contiene un elemento \verb'<select>' con identificativo uguale a \verb'selectState';
  \item[--] un terzo \verb'<div>' con all'interno:
  \begin{itemize}
    \item[-] un elemento \verb'<ul>' con identificativo uguale a \verb'ToolsList' per la lista degli strumenti, contenente:
      \begin{itemize}
        \item[$\cdot$] un \verb'<li>' con identificativo \verb'liAnswering' contenente il nodo di testo <<Segreteria>>;
        \item[$\cdot$] un \verb'<li>' con identificativo \verb'liSetting' contenente il nodo di testo <<Impostazioni>>;
        \item[$\cdot$] un \verb'<li>' con identificativo \verb'liCallList' contenente il nodo di testo  <<Lista chiamate>>;
        \item[$\cdot$] un \verb'<li>' con identificativo \verb'liGroup' contenente il nodo di testo <<Gruppi>>;
        \item[$\cdot$] un \verb'<li>' con identificativo \verb'liSearch' contenente il nodo di testo <<Ricerca>>;
      \end{itemize}
    \item[-] un elemento \verb'<ul>' con all'interno un solo \verb'<li>' avente come identificativo \verb'liLogout' e contenente come unico figlio il nodo di testo \verb'Logout'.
  \end{itemize}
\end{itemize}


%***********************************************************************
%*                QUESTO È IL CONFINE FRA IL BENE E IL MALE            *
%***********************************************************************
% TODO IBI ERIT FLETUS ET STRIDOR DENTIUM
\subsection{AddressBookView}
\subsubsection*{Funzione}
Vista associata al presenter \classname{AddressBookPresenter}. Utilizzata per la visualizzazione della rubrica personale.

\subsubsection*{Elementi}
\begin{itemize}
\item \textit{div} contenente un elemento di tipo \textit{text} e un \textit{button} per una ricerca veloce nella propria rubrica;
\item \textit{div} contenente un \textit{ul} dove verranno visualizzati tutti i contatti della rubrica;
\item \textit{div} contenente una \textit{select} utilizzata per ordinare la rubrica in base ai parametri dati;
\item \textit{div} contenente una \textit{select} dove sono elencati i gruppi della propria rubrica.
\end{itemize}

% TODO IBI ERIT FLETUS ET STRIDOR DENTIUM
\subsection{CommunicationView}
\subsubsection*{Funzione}
Vista associata al presenter \classname{CommunicationPresenter}. Utilizzata inizialmente durante una chiamata.

\subsubsection*{Elementi}
\begin{itemize}
\item \textit{div} contenente un elemento di tipo \textit{text} e un \textit{button} per una ricerca veloce nella propria rubrica;
\item \textit{div} contenente due elementi \textit{video} rispettivamente utilizzati per visualizzare lo stream della propria fotocamera e dell'altro client;
\item \textit{div} contenente 3 \textit{span} utilizzati per visualizzare le statistiche, ossia i dati inviati, i dati ricevuto e il tempo di comunicazione;
\item \textit{button} per terminare la chiamata.
\end{itemize}

% TODO IBI ERIT FLETUS ET STRIDOR DENTIUM
\subsection{ContactView}
\subsubsection*{Funzione}
Vista associata al presenter \classname{ContactPresenter}. Utilizzate per visualizzare il profilo di un utente.

\subsubsection*{Elementi}
\begin{itemize}
\item \textit{div} contenente un elemento di tipo \textit{ul} con 4 \textit{li}:
\begin{itemize}
\item nome;
\item cognome;
\item email;
\item gruppi a cui appartiene se il contatto è presente nella mia rubrica.
\end{itemize}
\item \textit{button} per avviare una chiamata audio;
\item \textit{button} per avviare una video chiamata;
\item \textit{button} per avviare una chat;
\item \textit{button} per lasciare un messaggio in segreteria;
\item \textit{button} per aggiungere il contatto nella propria rubrica se non presente;
\item \textit{button} per rimuovere il contatto dalla propria rubrica se presente;
\item \textit{button} per bloccare il contatto se presente nella propria rubrica e non bloccato;
\item \textit{button} per sbloccare il contatto se presente nella propria rubrica e bloccato.
\end{itemize}

\clearpage

\section{Tracciamenti}

Al fine di rendere evidente l'associazione di necessità e sufficienza che intercorre tra i requisiti (emersi in attività di analisi) e le classi predisposte per soddisfarli, riportiamo di seguito due tabulati di tracciamento. Il primo, requisiti-classi, mostra il soddisfacimento di ogni requisito identificato. Il secondo, classi-requisiti, evidenzia l'utilità delle classi e dimostra che ogni una di esse è stata creata con uno scopo ben preciso.

\subsection{Requisiti - classi}

\subsection{Classi - requisiti}

\end{document}
